\documentclass{article}
\PassOptionsToPackage{numbers, compress}{natbib}
% if you need to pass options to natbib, use, e.g.:
%     \PassOptionsToPackage{numbers, compress}{natbib}
% before loading neurips_2026

% The authors should use one of these tracks.
% Before accepting by the NeurIPS conference, select one of the options below.
% 0. "default" for submission
% \usepackage{neurips_2026}
% the "default" option is equal to the "main" option, which is used for the Main Track with double-blind reviewing.
% 1. "main" option is used for the Main Track
%  \usepackage[main]{neurips_2026}
% 2. "position" option is used for the Position Paper Track
%  \usepackage[position]{neurips_2026}
% 3. "eandd" option is used for the Evaluations & Datasets Track
 % \usepackage[eandd]{neurips_2026}
 % if you want to opt-in for a de-anomymized submission (not recommended, but OK):
 %\usepackage[eandd, nonanonymous]{neurips_2026}
% 4. "creativeai" option is used for the Creative AI Track
%  \usepackage[creativeai]{neurips_2026}
% 5. "sglblindworkshop" option is used for the Workshop with single-blind reviewing
 % \usepackage[sglblindworkshop]{neurips_2026}
% 6. "dblblindworkshop" option is used for the Workshop with double-blind reviewing
%  \usepackage[dblblindworkshop]{neurips_2026}

% After being accepted, the authors should add "final" behind the track to compile a camera-ready version.
% 1. Main Track
 % \usepackage[main, final]{neurips_2026}
% 2. Position Paper Track
%  \usepackage[position, final]{neurips_2026}
% 3. Evaluations & Datasets Track
 % \usepackage[eandd, final]{neurips_2026}
% 4. Creative AI Track
%  \usepackage[creativeai, final]{neurips_2026}
% 5. Workshop with single-blind reviewing
%  \usepackage[sglblindworkshop, final]{neurips_2026}
% 6. Workshop with double-blind reviewing
%  \usepackage[dblblindworkshop, final]{neurips_2026}
% Note. For the workshop paper template, both \title{} and \workshoptitle{} are required, with the former indicating the paper title shown in the title and the latter indicating the workshop title displayed in the footnote.
% For workshops (5., 6.), the authors should add the name of the workshop, "\workshoptitle" command is used to set the workshop title.
% \workshoptitle{WORKSHOP TITLE}

% "preprint" option is used for arXiv or other preprint submissions
 \usepackage[preprint]{neurips_2026}

% to avoid loading the natbib package, add option nonatbib:
%    \usepackage[nonatbib]{neurips_2026}

\usepackage[utf8]{inputenc} % allow utf-8 input
\usepackage[T1]{fontenc}    % use 8-bit T1 fonts
\usepackage{hyperref}       % hyperlinks
\usepackage{url}            % simple URL typesetting
\usepackage{booktabs}       % professional-quality tables
\usepackage{array}          % m{} column type for vertical centering
\usepackage{amsmath}        % \eqref, align, cases, etc.
\usepackage{amsfonts}       % blackboard math symbols
\usepackage{amsthm}         % theorem/proof environments
\usepackage{nicefrac}       % compact symbols for 1/2, etc.
\usepackage{microtype}      % microtypography
\usepackage[table]{xcolor}  % colors (table option for \cellcolor)
\usepackage{graphicx}       % image inclusion
\usepackage{wrapfig}        % text-wrapping figures/tables

\usepackage{tabularx}
\usepackage{array}
\usepackage{caption}
\usepackage{pifont}
\usepackage{makecell}
\usepackage{multirow}
\usepackage{placeins}        % \FloatBarrier to pin floats within sections
\usepackage{float}           % [H] placement specifier

\newtheorem{lemma}{Lemma}
\newtheorem{proposition}{Proposition}

\graphicspath{{images/}{figures/}}    % default image paths

% Ranking colors for quantitative tables (shared with RALU-style tables)
\definecolor{tabfirst}{rgb}{1, 0.7, 0.7}
\definecolor{tabsecond}{rgb}{1, 0.85, 0.7}
\definecolor{tabthird}{rgb}{1, 1, 0.7}

\newtheorem{theorem}{Theorem}

\definecolor{mygray}{gray}{0.95}

% Note. For the workshop paper template, both \title{} and \workshoptitle{} are required, with the former indicating the paper title shown in the title and the latter indicating the workshop title displayed in the footnote. 
% \title{Wavelet Diffusion: Efficient Image and Video Generation using Multi-scale Spectral Representation}
\title{Spectral Progressive Diffusion \\ for Efficient Image and Video Generation}


% The \author macro works with any number of authors. There are two commands
% used to separate the names and addresses of multiple authors: \And and \AND.
%
% Using \And between authors leaves it to LaTeX to determine where to break the
% lines. Using \AND forces a line break at that point. So, if LaTeX puts 3 of 4
% authors names on the first line, and the last on the second line, try using
% \AND instead of \And before the third author name.


\begin{document}
\author{%
  Howard Xiao \quad Brian Chao \quad Lior Yariv \quad Gordon Wetzstein\vspace{1ex}\\
  Stanford University \\
  \url{https://howardxiao.ca/speed/}
}

\maketitle

\makeatletter
\let\arxiv@savedaddcontentsline\addcontentsline
\def\arxiv@tocfile{toc}
\renewcommand{\addcontentsline}[3]{%
  \def\arxiv@targetfile{#1}%
  \ifx\arxiv@targetfile\arxiv@tocfile
  \else
    \arxiv@savedaddcontentsline{#1}{#2}{#3}%
  \fi
}
\makeatother





\definecolor{mg}{rgb}{0.639,0.984,0.722}
\definecolor{my}{rgb}{0.996,0.875,0.643}
\definecolor{mr}{rgb}{0.941,0.561,0.620}
\newcommand{\ccg}{\cellcolor{mg}}  % cell color green
\newcommand{\ccy}{\cellcolor{my}}  % cell color yellow
\newcommand{\ccr}{\cellcolor{mr}}  % cell color red

\newcommand{\cmark}{\ding{51}}%
\newcommand{\xmark}{\ding{55}}%
\newcommand{\bluecheck}{{\color{blue}\checkmark}}
\newcommand{\blackcheck}{{\color{black}\cmark}}
\newcommand{\redx}{{\Large\color{\red}\xmark} }

% table font size
\newcommand{\headfontsize}{\scriptsize}
\newcommand{\tablerefsize}{\tiny}
\newcommand{\vertice}{\mathbf{v}}
\newcommand{\normal}{\mathbf{n}}



% cell colors + check / triangle / x mark
\newcommand*\colourcheck[1]{%
  \expandafter\newcommand\csname #1check\endcsname{\textcolor{#1}{\ding{51}}}%
}
\colourcheck{red}
\colourcheck{green}

\newcommand*\colourxx[1]{%
  \expandafter\newcommand\csname #1xx\endcsname{\textcolor{#1}{\ding{55}}}%
}
\colourxx{red}
\colourxx{green}

\newcommand*\colourtriangle[1]{%
  \expandafter\newcommand\csname #1triangle\endcsname{\textcolor{#1}{$\varDelta$}}%
}
\colourtriangle{yellow}


%%%%%%% Color cells, black font
\newcommand{\greenvcell}{\ccg \cmark}
\newcommand{\redxcell}{\ccr \xmark}
\newcommand{\yellowtrcell}{\ccy $\varDelta$} %\triangle
\newcommand{\greenxcell}{\ccg \xmark}
\newcommand{\redvcell}{\ccr \cmark}

%%%%%%% No background, color font
% \newcommand{\greenvcell}{\greencheck}
% \newcommand{\redxcell}{\redxx}
% \newcommand{\yellowtrcell}{\yellowtriangle} %\triangle
% \newcommand{\redvcell}{\redcheck}
% \newcommand{\greenxcell}{\greenxx}

\newcommand{\hlineB}[1]{\noalign{\hrule height #1\arrayrulewidth}}

\newcommand{\bc}[1]{\textcolor{blue}{bc: #1}}
\newcommand{\hx}[1]{\textcolor{purple}{hx: #1}}
\newcommand{\ly}[1]{\textcolor{teal}{ly: #1}}

\newcommand{\methodname}{\textbf{Spectral Progressive Diffusion}}
\begin{abstract}
Diffusion models have been shown to implicitly generate visual content autoregressively in the frequency domain, where low-frequency components are generated earlier in the denoising process while high-frequency details emerge only in later timesteps. 
This structure offers a natural opportunity for efficient generation, as high-resolution computation on noise-dominated frequencies is largely redundant.
We propose \emph{Spectral Progressive Diffusion}, a general framework that progressively grows resolution along the denoising trajectory of pretrained diffusion models. 
To this end, we develop a spectral noise expansion mechanism and derive an optimal resolution schedule from the model's power spectrum. Our framework supports training-free acceleration and a novel fine-tuning recipe that further improves efficiency and quality.
We demonstrate significant speedups on state-of-the-art pretrained image and video generation models while preserving visual quality.

% In this work, we propose \methodname, a token-efficient visual generation method based on progressive frequency processing. We exploit the spectral autoregressive structure of diffusion models to integrate frequency-domain transitions \textit{within a single denoising trajectory}, eliminating computation on noise-dominated high-frequency components. We derive optimal resolution transition times from the resolution schedule with minimal, \textit{single-hyper-parameter} tuning for our training-free setting, and introduce an optional fine-tuning procedure to further improve quality, achieving significant speedups when applied to state-of-the-art pretrained image and video generation models.

\end{abstract}

\section{Introduction}
\label{sec:intro}


Computational demands for visual generative models are increasing rapidly as image resolutions and video sequence lengths continue to grow. This trend reveals a fundamental scaling crisis: while scaling model and data sizes consistently yields better generation quality, the underlying cost of self-attention in Diffusion Transformers (DiTs)~\citep{dit} scales quadratically with the number of generated tokens. This creates a growing conflict between the desire for high-fidelity, long-duration content and the compute cost of existing architectures. A shift toward token-efficient representations offers a flexible and broadly compatible path for further scaling image and video generation.

Rissanen et al.~\citep{rissanen2023generative} and Dieleman~\citep{dieleman2024spectral} observed that diffusion models implicitly learn to generate visual content autoregressively in the frequency domain, where low-frequency components are generated earlier in the denoising process while high-frequency details emerge only in later timesteps.  While this observation has motivated a variety of frequency-domain designs for diffusion models \cite{wavediff, spectral-reg, fourierspace, frecas}, they either do not support existing pretrained models or provide limited efficiency gains. A natural way to exploit spectral autoregression is through progressive resolution growth, since the frequency content representable by a signal is intrinsically tied to its spatial resolution. However, prior progressive resolution approaches \cite{pyrflow, scalespace, ralu, lssgen} either require significant modifications to the model architecture or rely heavily on heuristics for when and how to upsample, limiting their compatibility with state-of-the-art pretrained models \cite{zimage, flux-2, wan, cogvideox}.

In this paper, we propose \emph{Spectral Progressive Diffusion}: a general framework for progressive-resolution generation, guided by the spectral structure of the denoising process. Motivated by the spectral autoregression property of diffusion models \cite{dieleman2024spectral}, we expand the resolution only at the timesteps where high-frequency content begins to emerge from noise (see Fig.~\ref{fig:method}). This keeps early denoising steps in a reduced token space and avoids redundant computation on noise-dominated frequencies. 
To achieve this, we introduce a spectral noise expansion mechanism that uses a spectral transformation to inject high-frequency noise at the correct level while preserving the partially-denoised low-frequency content. We further derive optimal resolution transition times directly from the model's power spectrum, determined by a single error-tolerance hyperparameter. 

Our framework applies directly to state-of-the-art pretrained diffusion and flow-matching models without architectural modifications. We demonstrate the flexibility of our framework across three visual generation modalities: latent-space image generation~\cite{black2024flux, zimage}, pixel-space image generation~\cite{pixelgen}, and latent-space video generation~\cite{wan}. In a training-free setting, our method delivers immediate speedups on existing pretrained models. We additionally introduce a fine-tuning recipe that further improves efficiency and quality, a direction previously unexplored for progressive-resolution generation. Extensive experiments show speedups of up to $7\times$ on image generation and $2.5\times$ on video generation, outperforming prior spatial acceleration methods in both runtime and visual fidelity.

To summarize: \vspace{-3pt}
\begin{itemize}
\item We propose a progressive-resolution generation framework guided by the spectral autoregression of diffusion models, that applies directly to pretrained models and improves efficiency.
\item We introduce a principled spectral noise expansion mechanism and derive optimal resolution transition times from the model's power spectrum.
\item We demonstrate our framework on both image and video generation models, supporting both training-free inference and fine-tuning while maintaining high generation quality.
\end{itemize}



% The resulting frequency-growing process significantly reduces the effective number of tokens processed within the denoising trajectory.


% In a related line of work, seminal discussions by \citet{dieleman2024spectral} and \citet{fourierspace} motivated substantial subsequent research on spectral analysis of diffusion models. Most notably, diffusion models exhibit an implicit autoregressive structure in the frequency domain, recovering low-frequency components early in denoising and high-frequency details only at later timesteps. This perspective aligns with cascaded diffusion approaches, where progressively increasing resolution corresponds to generating increasingly higher-frequency content. However, prior methods do not leverage this spectral autoregressive structure to accelerate generation \textit{within a single denoising trajectory}, leading to suboptimal efficiency improvements as they require low-resolution images to be fully denoised first. 

%%%%%%%%%%%%% old begin %%%%%%%%%%%%%%%%%
% Computational demands for visual generative models are increasing rapidly as image resolutions and video sequence lengths continue to grow. This challenge is exacerbated in modern generative model architectures based on Diffusion Transformers (DiTs) \cite{dit}, where the cost of attention scales quadratically with the number of generated tokens, leading to significant memory and latency constraints. Prior progressive or cascaded diffusion approaches \cite{cdm, relaydiff, pyrflow, matryoshka, pyrdiff, cogview3, lssgen, scalespace} improve efficiency by conditioning high-resolution synthesis on low-resolution images. However, they require significant modifications to the model architecture or the diffusion process, rendering them incompatible with existing, powerful pretrained models \cite{zimage, flux-2, wan2025, cogvideox}.

% In a related line of work, seminal discussions by \citet{dieleman2024spectral} and \citet{fourierspace} motivated substantial subsequent research on spectral analysis of diffusion models. Most notably, diffusion models exhibit an implicit autoregressive structure in the frequency domain, recovering low-frequency components early in denoising and high-frequency details only at later timesteps. This perspective aligns with cascaded diffusion approaches, where progressively increasing resolution corresponds to generating increasingly higher-frequency content. However, prior methods do not leverage this spectral autoregressive structure to accelerate generation \textit{within a single denoising trajectory}, leading to suboptimal efficiency improvements as they require low-resolution images to be fully denoised first. 

%Furthermore, there is no approach that exploits frequency-ordered generation to improve efficiency while remaining compatible with existing pretrained diffusion models.

% this can just be explained later to save space. We start from a detailed analysis of the power spectral density of noisy pixels and latents along the diffusion denoising trajectory. Our analysis reveals a stark disparity in signal-to-noise ratio (SNR) evolution: low-frequency structural components are robust to high noise, whereas high-frequency details remain completely swamped by noise during the early denoising steps.

%%%%%%%%%%%%% old end %%%%%%%%%%%%%%%%%
\iffalse
In this paper, we propose \emph{Spectral Progressive Diffusion}: a token-efficient generation framework driven by progressive frequency processing. Motivated by the spectral autoregression structure of diffusion models \cite{fourierspace, dieleman2024spectral}, we directly integrate spectral transformations into a single denoising trajectory for progressive resolution growing within a single diffusion pass. By isolating the low-frequency components with spectral transformations, we restrict early denoising steps to a heavily compressed token space, bypassing redundant computation on noise-dominated high-frequency tokens. We introduce a principled spectral noise expansion mechanism and derive optimal resolution transition times to progressively introduce high-frequency noise without corrupting the partially-denoised low-frequency content. This frequency-growing process significantly reduces the effective number of tokens processed within a single denoising trajectory. %We also design an optional fine-tuning procedure to further enhance image quality. Extensive experiments on pretrained image and video generation models demonstrate that our method achieves an X times speedup compared to na\"ive high-resolution generation, and outperforms existing acceleration methods in both runtime and visual fidelity. 

Our approach neither changes diffusion trajectories nor requires a specialized model architecture; therefore, it is backward-compatible with state-of-the-art pretrained diffusion and flow-matching models in a plug \& play manner. We demonstrate this capability by fine-tuning several pretrained latent and pixel-space image and video generation models into more efficient spectrally progressive variants of themselves. A comparable fine-tuning strategy has, to our knowledge, not been demonstrated to date. 

Extensive experiments on pretrained image and video generation models demonstrate that our method achieves an X times speedup compared to na\"ive high-resolution generation, and outperforms existing acceleration methods in both runtime and visual fidelity. 

%Furthermore, we demonstrate novel applications uniquely enabled by our spectral-transform-based diffusion framework, including X.


%motivate the need for efficient generation, for diffusion and flow models  

% motivate spectral autoregression (blog post etc.); this is based on the insight that the frequency distribution of natural images follow a power law; convince people here that some of the latent representations follow the same power power law

% outline desirable properties of any specific method that might implement the idea of ``efficient multi-scale generation''; review existing methods, such as scale-space diffusion, casceaded pixel-level diffusion, a possible naive Fourier representation, and others based on these criteria and convince people that these do not match your criteria. possible includes: (1) pretrained models and fine-tuning; (2) diffusion vs. flow; (3) ``distillability'' into few-step; (4) ... other arguments

% now describe idea of wavelet diffusion and define what this means and how it's implemented; why is this the right approach, i.e., why do we need wavelets, and others not? make it super intuitive! be able to convince a skeptical person here

% discuss the experiments we show and what the insights are we validate; also discuss some practical benefits we show in terms of speedup factors we show

To summarize,
\begin{itemize}
    \item We propose a token-efficient, spectral-transformation-based progressive generation framework, compatible with training-free inference and LoRA fine-tuning across pretrained image and video model backbones.
    \item We derive optimal resolution transition times from the resolution schedule and error thresholds, enabling training-free speedups with minimal hyperparameter tuning of just one parameter.
    \item
    We demonstrate that our framework significantly accelerates pretrained diffusion models with negligible quality degradation, achieving up to \hx{Y$\times$} and \hx{Z$\times$} speedups on image and video generation, respectively.
\end{itemize}
\fi

% !TEX root = ../main.tex
\begin{figure}[t!]
  \centering
  \includegraphics[width=\textwidth]{figures/pipeline.001.jpeg}\vspace{-2pt}
  \caption{\textbf{Spectral Progressive Diffusion.} We progressively grow the resolution along the denoising trajectory using an \textit{optimal resolution schedule} derived from the spectral power of pretrained models \textbf{(left)}. At each scheduled transition, our \textit{spectral noise expansion} mechanism \textbf{(right)} injects high-frequency noise at the correct level while preserving the partially-denoised low-frequency content.}
  \label{fig:method}
\end{figure}


\section{Related Work}

\label{sec:related}

%\hx{I suggest that we change 2.3 to something like "Other Efficient Generation Methods". We should talk about feature caching (TeaCache, ToCa, etc..), Token merging/pruning/Sparse Attention, and Distillation. I included the table to distinguish us from others in 2.1. But this may not be needed in the paper, for our benefits for now.}

%\subsection{Progressive and multi-scale generation}
%\label{subsec:related_prog_gen}
%\hx{Maybe we can consider saying the customized upsampler modules break i.i.d Gaussian distribution the model was pretrained with.}

\vspace{-4pt}
\paragraph{Spectral-domain designs in diffusion models.} The seminal discussion by \citet{dieleman2024spectral} on the spectral autoregression property of diffusion has inspired many spectral-domain designs in diffusion models. Most of these models \cite{wavediff, wavescore, dimsum, wfvae, wavedm, du2025flowk, dctdiff} require training from scratch, making them incompatible with existing pretrained image and video generation models \cite{flux-2,zimage,wan,cogvideox}. Other methods use frequency-aware losses~\cite{lwd, spectral-reg}, autoencoder regularization~\cite{diffusability}, or noise schedules~\cite{beta-sampling, spectralschedule, spectrallyguided, fourierspace} to improve the generation quality of existing models but do not improve efficiency. Recent training-free methods use spectral-domain designs for efficient high-resolution generation \cite{frecas, diffusehigh, hiwave}. However, as they require fully denoised lower-resolution images, their overall efficiency gains are limited. Our method leverages similar insights of spectral autoregression but progressively grows resolution within a single denoising trajectory, achieving greater efficiency gains while maintaining compatibility with pretrained models. 

% support fine-tuning existing models by reweighting the loss \cite{lwd, spectral-reg} to refine high-frequency details, but do not improve token efficiency. Spectral-analysis-based noise schedules \cite{beta-sampling, spectralschedule, fourierspace} improve image quality on small, low-resolution datasets (e.g. CIFAR-10, ImageNet $64\times64$), but they also don't improve token efficiency and have not been validated on modern pretrained models. 


% HiWave \cite{hiwave} and DiffuseHigh \cite{diffusehigh} enhance high-frequency details after upsampling a generated low-resolution image, while FreCaS \cite{frecas} uses a multi-stage cascade with repeated VAE decode–encode and renoising steps. 

% Spatial upsampling approaches reduce the effective number of processed tokens throughout the generation process by denoising low-resolution tokens and upsampling them \cite{bottleneck, ralu, demofusion, fresco, lssgen, freescale}. 

\vspace{-3pt}
\paragraph{Progressive resolution for efficient generation.} 
% Progressively changing resolution during the diffusion process has been proposed as a method to reduce the effective number of processed tokens and improve the efficiency of image and video generation. 
% Prior works have explored progressive resolution as a way to process fewer tokens for more efficient image and video generation. 
Progressive resolution approaches have been proposed to improve generation efficiency by reducing token counts early in the denoising process.
However, most existing methods~\cite{cdm, jing2022subspace, edifyimage, relaydiff, pyrflow, matryoshka, pyrdiff, cogview3, scalespace} require training from scratch due to model-specific architectures, limiting their compatibility with pretrained models. Current training-free~\cite{bottleneck, ralu, demofusion, fresco, freescale} and fine-tuning progressive resolution approaches~\cite{selfcascade, lssgen} either require significant heuristics-based hyperparameter tuning or specialized modules that limit their effectiveness and their compatibility with adaptation methods such as LoRA~\cite{hu2021lora}. Progressive resolution generation has also been explored in visual autoregressive models \cite{var, infinity, flowar, huang2025spectralar}, where next-token prediction is reformulated as next-scale prediction. Our method uniquely leverages the spectral autoregression property of diffusion models to design a principled progressive resolution growing mechanism for efficient image and video generation using minimal hyperparameters. Furthermore, our method is fully compatible with current pretrained models for both training-free inference and fine-tuning. 

% Other approaches \cite{selfcascade, lssgen} suggest fine-tuning approach but rely on specialized upsampler modules tightly coupled to specific models and architectures, making them less flexible and incompatible with adaptation methods such as LoRA~\cite{hu2022lora}. Training-free progressive resolution approaches \cite{bottleneck, ralu, demofusion, fresco, freescale} reduce the effective number of processed tokens throughout the denoising process, but require significant heuristics-based hyperparameter tuning. 


% Scale Space Diffusion \cite{scalespace} is most closely related to our approach, but requires a new UNet-based architecture trained from scratch, making it incompatible with modern DiT backbones.

% While the aforementioned works all realize progressive generation through cascaded, multi-stage conditional diffusion processes, where the diffusion process in each stage is conditioned on the lower-resolution output from the previous stage, they do not implement progressive generation \textit{within} a single diffusion process. This multi-stage framework greatly complicates the overall model design and results in sub-optimal speedups for high-resolution generation. The recently proposed Scale Space Diffusion \cite{scalespace} is the most similar in concept to ours where progressive resolution growing is directly baked in a single diffusion process, yet they also require a completely new UNet-based architecture design for pretraining, rendering it incompatible with modern DiT architectures. Unlike these prior approaches, our method is completely agnostic of the model architecture and compatible with existing, powerful pretrained models.

%\bc{add more differences in spectral motivation here}

%\subsection{Frequency-domain designs in visual generation}
%\bc{need to better motivate token efficiency here. token efficiency is not directly connected to freq.domain analysis here}


% motivate with diffusion is spectral autoregression blogposts, fourier space paper

% scheduling designs, but no speedups
% what papers actually leverage different velocity for low and high res 

%- diffusion is spectral AR
%- latent wavelet diffusion, diffuse high, hiwave, freeU, dimsum, frecas, %multiscale wave, sagd, spectral schedule, wavediff, wdno, wfvae
%- focus on speedup design
%- other spectral AR works, wala, 

%\subsection{Other methods for efficient generation}

\paragraph{Other methods for efficient generation.} Aside from progressive resolution pipelines, a wide variety of acceleration methods have been proposed for efficient generation. Model distillation approaches reduce the number of denoising steps by training a student model to approximate the outputs of a pretrained teacher \cite{dmd, dmd2, causvid, lcm, pcm, prog-dist}. Feature caching approaches reuse features across denoising timesteps to reduce redundant computation \cite{teacache, toca, freqca, ma2023deepcache}. Token merging techniques identify redundant tokens and merge similar tokens based on predefined heuristics \cite{lu2025toma, haurum2024agglomerative, bolya2023token, kim2024tokenfusion, wu2025importance, lee2025local, fang2025attend}. Sparse attention mechanisms identify sparse patterns in the attention maps and use block-wise calculations to improve efficiency \cite{li2025radial, xia2025trainingfree, zhan2025bidirectional, katharopoulos202linear, wang2020linformer, choromanski2021performer, beltagy2020longformer}. In this paper, we present an orthogonal and complementary progressive resolution generation approach that accelerates both image and video generation through spectral-domain transformations.



%To work with a vast 
%- spatial upsampling: bottleneck sampling, ralu, demofusion, fresco, lssgen, "a fourier space perspective of diffusion models", freescale
%- feature caching; teacache, taylorseer, toca, freqca, t-stitch 
%- others: token merging, sparse attention, feature caching, token and layer routing, etc

\iffalse
others: cascaded diffusion needs to go somewhere


\subsection{Scale space generation}
- overview a bunch of methods
- then say explicitly what the goal of our method is (might be different from some of the prior works' goals) and that we're using wavelets

\subsection{Wavelets and their applications in Generative Modelling}
- now talk about all the other ways wavelets have been used in generative models, including GANs
- our goals might be different from these other methods

Prior work has injected wavelet structure into generative pipelines at three distinct points, none of which coincides with ours. \emph{Pixel-space wavelet diffusion}~\citet{phung2023wavediff, kiruluta2025hybridwf} runs the denoiser on DWT coefficients of RGB images for a $\sim$$2$--$3\times$ speedup, but cannot leverage pretrained latent backbones and is restricted to small image benchmarks. \emph{Wavelet-structured VAEs}~\citet{li2025wfvae} push the transform inside the autoencoder for cheaper video encoding, leaving the downstream diffusion model unaware of the frequency structure --- an approach that is orthogonal and complementary to ours. \emph{Wavelet-guided training of latent diffusion}~\citet{sigillo2025lwd}, the closest prior work, computes a wavelet saliency map of latents to re-weight the flow-matching loss for ultra-high-resolution images, but the denoiser, the latent representation, and the sampling procedure are unchanged; the wavelet appears nowhere at inference and yields no compute savings. A parallel autoregressive line~\citet{huang2025spectralar, mattar2024wavelets} corroborates the coarse-to-fine view of generation~\citet{dieleman2024spectral, rissanen2023generative} but in image-only AR settings disjoint from latent video diffusion. In contrast, we decompose video \emph{latents} into wavelet sub-bands as the actual diffusion target, derive a per-sub-band flow trajectory in which each band activates only when its SNR crosses a threshold, and translate this into a $5.3\times$ attention speedup at sampling time --- all realized through LoRA fine-tuning on a frozen video DiT.
\fi

% \newcolumntype{Y}{>{\centering\arraybackslash}X}
% \setlength{\tabcolsep}{4pt}
% \setlength{\textfloatsep}{10pt}
% \renewcommand{\cellalign}{tc}
% \renewcommand{\theadalign}{cc}
% \renewcommand{\theadfont}{\bfseries\footnotesize}
% \begin{table}[!t]
%     \centering
%     \captionsetup{aboveskip=0pt, belowskip=3pt}
%     {\footnotesize
%     \caption{\textbf{Comparison with progressive-resolution acceleration methods for native-resolution generation.} Our framework is the only one that combines spectral-domain transitions with LoRA fine-tuning on arbitrary pretrained backbones, with a spectrum-derived resolution schedule, for both image and video.}
%     \label{tab:method_comparison}
%     \renewcommand{\arraystretch}{1.4}
%     \begin{tabularx}{\linewidth}{>{\centering\arraybackslash}m{3.0cm}|YYYYYY}
%         \thead{Method} & \thead{Spectral\\transition} & \thead{Training-\\free} & \thead{LoRA\\fine-tuning} & \thead{Resolution\\schedule} & \thead{Image} & \thead{Video} \\
%         \hline
%         Jenga \cite{jenga}                    & \redxcell   & \greenvcell & \redxcell   & \redxcell   & \redxcell   & \greenvcell \\
%         \hline
%         RALU \cite{ralu}                      & \redxcell   & \greenvcell & \redxcell   & \redxcell   & \greenvcell & \redxcell \\
%         \hline
%         Bottleneck Sampling \cite{bottleneck} & \redxcell   & \greenvcell & \redxcell   & \redxcell   & \greenvcell & \greenvcell \\
%         \hline
%         Pyramid Flow \cite{pyrflow}           & \redxcell   & \redxcell   & \redxcell   & \redxcell   & \redxcell   & \greenvcell \\
%         \hline
%         Matryoshka \cite{matryoshka}          & \redxcell   & \redxcell   & \redxcell   & \redxcell   & \greenvcell & \redxcell \\
%         \hline
%         LSSGen \cite{lssgen}                  & \redxcell   & \redxcell   & \redxcell   & \redxcell   & \greenvcell & \redxcell \\
%         \hlineB{3}
%         \textbf{Ours}                         & \greenvcell & \greenvcell & \greenvcell & \greenvcell & \greenvcell & \greenvcell \\
%         \hline
%     \end{tabularx}
%     }
% \end{table}

% Previous algorithm-class overview table (kept for reference):
% \newcolumntype{Y}{>{\centering\arraybackslash}X}
% \setlength{\tabcolsep}{4pt}
% \setlength{\textfloatsep}{10pt}
% \begin{table}[t!]
%     \centering
%     \captionsetup{aboveskip=0pt, belowskip=3pt}
%     {\footnotesize
%     \caption{\textbf{Overview of efficient visual generation with spatial upsampling.} \bc{i need to think of this table more}}
%     \label{tab:method_comparison_old}
%     \begin{tabularx}{\linewidth}{m{4cm}|YYYY}
%         \textbf{Algorithm class} & \textbf{Post-trainable / distill-able} & \textbf{Training-free Possible} & \textbf{Efficient Training} & \textbf{Efficient Inference} \\
%         \hline
%         Architectural changes & \redxcell & \redxcell & \yellowtrcell & \yellowtrcell \\
%         \hline
%         Training-free acceleration & \redxcell & \greenvcell & N/A & \greenvcell \\
%         \hline
%         %Inference acceleration (temporal) & \redxcell & \greenvcell & N/A & \greenvcell \\ \hline
%         %Adaptive token \& layer routing & \greenvcell & \redxcell & \yellowtrcell & \greenvcell \\ \hline
%         Frequency domain modelling & \yellowtrcell & \redxcell & \yellowtrcell & \greenvcell \\
%         \hlineB{3}
%         Ours & \greenvcell & \greenvcell & \greenvcell & \greenvcell \\
%         \hline
%     \end{tabularx}
%     }
% \end{table}

% !TEX root = ../main.tex
%\section{Method}
%\label{sec:method}


% % ----- Brian's standing comments on this section (collected at section top) -----
% \bc{I suggest that we simply have two sections: Preliminaries and <Method Name>. At the very beginning of the methods section, we give a high-level overview and summarize what we will derive and explain in each of the sections.}

% \bc{As shown in Fig.~X, higher frequencies of latents emerge naturally throughout the denoising process. To fully leverage this spectral autoregression property of diffusion models for massive computational savings, we propose \methodname, a novel dito seeffusion pipeline in which progressive resolution growth occurs within a single denoising trajectory via spectral transformations and noise padding. During denoising, partially denoised low-resolution latents are transformed to the spectral domain at closed-form transition times, augmented with high-frequency noise, and mapped back to the spatial domain via an inverse transform. This progressively increases resolution while avoiding computation on noise-dominated high-frequency tokens.}

% \bc{(on the power-spectrum paragraph, now in Preliminaries.) I would actually move this section to preliminaries. At the very beginning of Sec.~\ref{sec:method:noise-injection}, we can then clearly motivate our method with this subsection. See my rewrite of Sec.~\ref{sec:method:noise-injection}. The ``noise-floor'' plot where specific frequencies emerge should also happen before this section.}

% \bc{(on the training-free subsection.) I changed injection to expansion because this signals resolution growth.}
% % --------------------------------------------------------------------------------

\section{Preliminaries}
\label{sec:method:prelim}

\subsection{Diffusion Models and Flow Matching}
\label{sec:method:prelim:basics}

Diffusion models~\citep{ho2020denoising,song2021score} define a generative process that gradually transforms samples from a simple prior (e.g. Gaussian) to data samples via a learned reverse-time process. Flow Matching \cite{liu2022rectified,lipman2023flow} reformulates the diffusion process as a continuous-time optimal transport problem and learns a velocity field that deterministically transforms noisy samples to data samples along straight paths.

Specifically, given a clean data point $\mathbf{x}_0$ sampled from a data distribution $p_{\text{data}}$, and a noise sample $\mathbf{x}_1 = \boldsymbol{\epsilon} \sim \mathcal{N}(0, I)$, the noise-to-data path is a straight line defined by
\begin{equation}
  \mathbf{x}_t \;=\; (1-t)\, \mathbf{x}_0 \,+\, t\,\boldsymbol{\epsilon},
  \qquad t \in [0, 1].
  \label{eq:flow-forward}
\end{equation}
A neural network
$\mathbf{v}_\theta(\mathbf{x}_t, t)$ is trained to predict the target velocity $\dot{\mathbf{x}}_t = \boldsymbol{\epsilon} - \mathbf{x}_0$ with the following objective:
\begin{equation}
  \mathcal{L}_\text{flow}(\theta) \;=\; \mathbb{E}_{t, \mathbf{x}_0, \boldsymbol{\epsilon}}\!\left[\,\|\mathbf{v}_\theta(\mathbf{x}_t, t) - (\boldsymbol{\epsilon} - \mathbf{x}_0)\|^2\,\right].
  \label{eq:flow-loss}
\end{equation}

A network with sufficient capacity that is trained on enough data converges to the \textit{Bayes-optimal velocity predictor} $\mathbf{v}_\theta(\mathbf{x}_t, t) = \mathbf{v}^*(\mathbf{x}_t, t)$~\citep{liu2022rectified, lipman2023flow}.

\iffalse

Under the flow-matching objective in Eq.~\eqref{eq:flow-loss}, the Bayes-optimal velocity predictor is the conditional expectation
\begin{equation}
  \mathbf{v}^*(\mathbf{x}_t, t) \;=\; \mathbb{E}\!\left[\,\boldsymbol{\epsilon} - \mathbf{x}_0 \,\big|\, \mathbf{x}_t\,\right],
  \label{eq:optimal-v}
\end{equation}
which is the unique minimizer of Eq.~\eqref{eq:flow-loss}~\citep{liu2022rectified,lipman2023flow}. A network with sufficient capacity trained on enough data converges to $\mathbf{v}^*$ in expectation.
\fi


During inference, data samples are generated through sampling from the prior $\mathbf{x}_1 \sim \mathcal{N}(0, I)$ and solving the probability-flow Ordinary Differential Equation (ODE) $\dot{\mathbf{x}}_t = \mathbf{v}_\theta(\mathbf{x}_t, t)$. Due to its faster training and stable convergence, recent image and video generative models \cite{flux-2,zimage,sd3,wan} have increasingly adopted the flow matching paradigm.


Most state-of-the-art diffusion and flow matching models adopt the Diffusion Transformer (DiT) architecture \citep{dit}, which processes tokens through a sequence of self-attention and Multilayer Perceptron (MLP)-based operations. This motivates our approach to reduce the number of tokens processed along the denoising trajectory using the spectral autoregression property of diffusion models \citep{dieleman2024spectral}.

%costing one forward pass of $v_\theta$ per step. 

%Our framework applies to both diffusion \eqref{eq:diff-forward} and flow-matching \eqref{eq:flow-forward} backbones.


% Modern image and video diffusion models often operate in a compressed latent space~\citep{rombach2022ldm} rather than on raw pixels.

\iffalse
Diffusion models~\citep{ho2020denoising,song2021score} define a forward process that progressively corrupts a clean sample $x_0 \sim p_\text{data}$ by adding i.i.d.\ Gaussian noise,
\begin{equation}
  x_t \;=\; \alpha_t\, x_0 \,+\, \sigma_t\,\epsilon,
  \qquad t \in [0, 1], \qquad \epsilon \sim \mathcal{N}(0, I),
  \label{eq:diff-forward}
\end{equation}
with schedule $(\alpha_t, \sigma_t)$ moving from $(1, 0)$ at $t = 0$ (clean data) to $(0, 1)$ at $t = 1$ (pure noise). A network $\epsilon_\theta(x_t, t)$ is trained to predict the noise by minimizing
\begin{equation}
  \mathcal{L}_\text{diff}(\theta) \;=\; \mathbb{E}_{t, x_0, \epsilon}\!\left[\,\|\epsilon_\theta(x_t, t) - \epsilon\|^2\,\right].
  \label{eq:diff-loss}
\end{equation}
Generation starts from a random sample $x_1 = \epsilon \sim \mathcal{N}(0, I)$ and numerically integrates the corresponding reverse-time process (the probability-flow ODE) backward from $t = 1$ to $t = 0$, with the score $\nabla_{x_t}\!\log p_t(x_t) = -\epsilon_\theta(x_t, t)/\sigma_t$ obtained via Tweedie's identity~\citep{efron2011tweedie}. 
\fi

\iffalse
\paragraph{Diffusion Transformers.}
\bc{can probably remove this section, since we will explain this in intro already}
Modern image and video generation models are predominantly built on Diffusion Transformer (DiT) backbones~\citep{dit}, where image and video pixels or latents are encoded into a sequence of $N$ tokens, then processed by a stack of self-attention and feed-forward MLP blocks. In particular, the computational cost of self-attention is quadratic in $N$, making efficient, high-resolution generation especially challenging. This motivates our approach: by restricting early denoising steps where high-frequency components remain noise-dominated to a reduced spatial resolution, we avoid unnecessary computation on uninformative high-frequency tokens, thereby significantly improving generation efficiency.
\fi

% Let $\mathbf{x}_0 \in \mathbb{R}^{C \times T \times H \times W}$ denote a clean data sample drawn from $p_\text{data}$ ($C$ channels and $T$ frames, where $T = 1$ for images and $T > 1$ for videos).


% !TEX root = ../main.tex
\begin{figure}[t]
  \centering
  \includegraphics[width=\linewidth]{figures/power_spectrum.pdf}
  \caption{\textbf{Diffusion process in the spectral domain.} Latent power spectra in both image and video models decay rapidly with frequency (Fig.~\textbf{(a)}), consistent with natural images. Diffusion exhibits a frequency-domain autoregressive structure (Fig.~\textbf{(b)}) due to the aforementioned property: low frequencies emerge early in the denoising process, while high frequencies remain noise-dominated.}
  \label{fig:wavelet-spectral}
\end{figure}


\subsection{Spectral Autoregression in Diffusion Models}
\label{sec:method:prelim:spectral}

Natural images exhibit a characteristic power-law decay \citep{ruderman1997origins}, which is the foundation of the spectral autoregression property of diffusion models \citep{dieleman2024spectral}. Specifically, let $\Phi = \{\phi_\omega\}_{\omega \in \Omega}$ be a spectral basis indexed by frequency $\omega \in \Omega$ (e.g., Fourier~\citep{fourier1822theorie}, discrete cosine~\citep{ahmed1974dct}, or Haar wavelet~\citep{haar1909theorie}), and denote the associated spectral transformation by $T_\Phi$. We define the per-frequency signal power, or \emph{power spectrum}, of a data distribution $p_{\text{data}}$ as
\begin{equation}
  P_\omega \;:=\; \mathbb{E}_{\mathbf{x}_0}\!\left[\,\bigl|x_0^{(\omega)}\bigr|^2\,\right],
  \quad \text{where  } x_0^{(\omega)} := \langle \mathbf{x}_0, \phi_\omega \rangle \text{ under $T_\Phi$}.
  \label{eq:Pb-def}
\end{equation}

When $p_\text{data}$ represents natural images, $P_\omega$ is known to exhibit a power-law decay, that is
\begin{equation}
  P_\omega \;\propto\; |\omega|^{-\beta},
  \label{eq:power-law}
\end{equation}
with the exponent term typically in the range $\beta \in [2, 3]$~\citep{dieleman2024spectral,fourierspace}. We validated similar trends in both image and video latent representations of modern latent-space diffusion models (see Fig. \ref{fig:wavelet-spectral}).

% Although modern diffusion models operate in VAE latent space rather than pixel space \citep{rombach2022ldm}, their frequency statistics remain largely preserved. We compute radially averaged power spectra of latent representations of image and video generation models and observe a consistent power-law decay (Fig. \ref{fig:wavelet-spectral}). 

%We estimate exponents of $\beta \approx 1.92$ for FLUX.1-dev image latents \citep{black2024flux} and $\beta \approx 2.42$ for WAN 2.1 video latents \citep{wan} (Fig.~\ref{fig:wavelet-spectral}).

The power-law decay of $P_\omega$ in Eq.~\eqref{eq:power-law} implies that high frequencies carry substantially weaker signal than low frequencies. Since the Gaussian noise $\boldsymbol{\epsilon}$ added during the forward process has a flat, i.e. frequency-independent, power spectrum, higher frequencies are dominated by noise early in the denoising trajectory. In this work, we leverage this spectral autoregression property to progressively grow resolution within a single denoising trajectory via spectral domain transformations for efficient image and video generation.

% As observed by \citet{dieleman2024spectral} and \citet{fourierspace}, diffusion models exhibit the spectral autoregression property, recovering low-frequency signals early in the denoising process and generating higher-frequency details only in the later denoising steps.


\section{Spectral Progressive Diffusion}
In this section, we present a spectral-transformation-based progressive generation framework that enables arbitrary resolution transitions within a single denoising trajectory for pretrained image and video generation models (Sec.~\ref{sec:method:noise-injection}). We then derive an optimal resolution transition schedule based on power spectrum analysis and resolution-dependent frequency limits (Sec.~\ref{sec:method:schedule:signal}). Finally, we introduce a principled fine-tuning strategy to further improve both generation quality and efficiency (Sec.~\ref{sec:method:lora}).

\label{sec:method:schedule}

% ----- Commented-out intro (previous 3.2 opening paragraph), kept for reference: -----
% This section develops the theoretical foundation of our resolution schedule. We project the latent onto an orthonormal basis of the frequency domain, decomposing flow matching into a collection of independent per-frequency subproblems, and show that the expected velocity prediction error on each frequency is determined entirely by the signal's power spectrum. This yields a principled per-frequency activation criterion based on when the velocity is reliably predictable from the noisy input. In this section we include our derivations for the flow-matching parameterization of Eq.~\eqref{eq:flow-forward}; the analogous derivation for diffusion models is provided in the supplementary material.
% ------------------------------------------------------------------------------------


\subsection{Training-Free Inference with Spectral Noise Expansion}
\label{sec:method:noise-injection}

% In practice, Eq.~\eqref{eq:stage-transition} is discretized into $S$ resolution stages $(0 < s_1 < s_2 < \cdots < s_S = 1)$, with stage $i$ held for $t \in (t_{s_i}^{\text{transition}},\, t_{s_{i-1}}^{\text{transition}}]$ and matched to the denoising timesteps of the pretrained models.

Motivated by the spectral autoregression property of diffusion (Sec.~\ref{sec:method:prelim:spectral}), we progressively increase image resolution by injecting higher-frequency components along the denoising trajectory using a spectral transformation. Our framework applies natively to existing pretrained image and video generation models, without additional custom modules or architectural modifications.

Given an orthonormal spectral basis $\Phi$, let $T_\Phi$ and $T_\Phi^{-1}$ denote the forward and inverse spectral transformations. We define $S$ progressive resolution scales $s_{1:S}$, where $0 < s_1 < s_2 < \cdots < s_S = 1$, and their corresponding resolution transition times $t_{1:S-1}$, with $1 > t_1 > \cdots > t_{S-1} > 0$, that are matched to the noise schedules of pretrained models. We denote the final full-resolution state as $\mathbf{x}_0^{s_S} \in \mathbb{R}^{C \times T \times H \times W}$, and assume stage $i$ runs at spatial resolution $(s_iH, s_iW)$ for $t \in (t_i,\, t_{i-1}]$. 
% In practice, we select $s_i$ such that resolutions $(s_i H, s_i W)$ align with the multi-resolution training distribution of pretrained generative models. For example, we set $s_i$ to be powers of $1/2$ for FLUX.1-dev, as it is primarily trained on image latents with power-of-two resolutions.

At each transition time $t_i$, we expand the current low-resolution state $\mathbf{x}^{s_i}_{t_i} \in \mathbb{R}^{C \times T \times (s_iH) \times (s_iW)}$ to the next resolution state $\mathbf{x}^{s_{i+1}}_{t_i} \in \mathbb{R}^{C \times T \times (s_{i+1}H) \times (s_{i+1} W)}$ using \emph{spectral noise expansion} and timestep alignment (see Fig.~\ref{fig:method}).

\paragraph{Spectral noise expansion.}
Given a low-resolution $\mathbf{x}^{s_i}_{t_i}$ at transition time $t_i$, we expand its resolution to the next scale $s_{i + 1}$ through the following steps:
\begin{enumerate}
  \item[(i)] Compute the spectrum $\boldsymbol{\xi}^{s_i}_{t_i} = T_\Phi(\mathbf{x}^{s_i}_{t_i})$ supported on the frequency set $\Omega_{s_i}$ representable at scale $s_i$.
  \item[(ii)] Embed $\boldsymbol{\xi}^{s_i}_{t_i}$ in the lower-frequency part of the spectrum for $\Omega_{s_{i + 1}}$ and fill the slots $\Omega_{s_{i+1}} \setminus \Omega_{s_i}$ with $t_i \epsilon$, $\epsilon \sim \mathcal{N}(0, 1)$ to get $\boldsymbol{\xi}^{s_{i + 1}}_{t_i}$. 
  \item[(iii)] Convert the spectrum back to spatial domain $\mathbf{x}^{s_{i+1}}_{t_i}$ at resolution $(s_{i+1}H, s_{i+1}W)$ via $\mathbf{x}^{s_{i+1}}_{t_i} = T_\Phi^{-1}(\boldsymbol{\xi}^{s_{i + 1}}_{t_i})$.
\end{enumerate}
Steps~(i)--(iii) expand the resolution scale from $s_i$ to $s_{i + 1}$ for $\mathbf{x}_{t_i}^{s_{i + 1}}$ by injecting high frequencies at the correct noise level while preserving the partially-denoised low-frequency content of $\mathbf{x}^{s_i}_{t_i}$.

\paragraph{Timestep alignment.}
As $\Phi$ is orthonormal, the high-frequency noise padding in $\Omega_{s_{i + 1}} \setminus \Omega_{s_i}$ raises the overall noise level of the enlarged $\mathbf{x}_{t_i}^{{s_{i+1}}}$, which makes it no longer correspond to the original timestep $t_i$. We therefore rescale the inverse-transformed output and re-index the timestep by:
\begin{equation}
  \tilde{\mathbf{x}}^{s_{i+1}}_{\tilde{t}_i} := \frac{s_{i+1}/s_i}{1 + ((s_{i+1}/s_i) - 1)t_i} \cdot \mathbf{x}^{s_{i+1}}_{t_i},
  % \tilde{\mathbf{x}}^{s_{i+1}}_{\tilde{t}_i} = \frac{s_{i+1}/s_i}{1 + ((s_{i+1}/s_i) - 1)t_i} \cdot \mathbf{x}^{s_{i+1}}_{t_i},
  \label{eq:scale-align}
\end{equation}
where
\begin{equation}
  \tilde{t}_i := \frac{(s_{i+1}/s_i)\,t_i}{1 + ((s_{i+1}/s_i) - 1)t_i}.
  \label{eq:time-align}
\end{equation}
Eq.~\ref{eq:scale-align} and Eq.~\ref{eq:time-align} give us the aligned state $\tilde{\mathbf{x}}^{s_{i+1}}_{\tilde{t}_i}$ at the aligned timestep $\tilde{t}_i$. The integration of the probability-flow ODE then resumes at resolution $(s_{i+1}H, s_{i+1}W)$ until the next transition time $t_{i+1}$. 
The full derivation of Eqs.~\eqref{eq:scale-align}--\eqref{eq:time-align} is provided in Appendix~\ref{app:time-align}. In practice, $S$ and $s_{1:S}$ are chosen to align with the multi-resolution training distribution of pretrained generative models. By default, we choose $T_\Phi$ as the Discrete Cosine Transform (DCT)~\citep{ahmed1974dct} with discrete cosine basis $\Phi$. Alternative spectral transforms (such as Discrete Wavelet Transform (DWT) and Fourier Transform (FFT)) are compared and evaluated in Sec.~\ref{sec:exp:ablations}.


\renewcommand{\arraystretch}{1.0}
\setlength{\aboverulesep}{0.2ex}
\begin{table*}[!t]
  \centering
  \scriptsize
  \begin{tabular*}{\textwidth}{@{\extracolsep{\fill}} l ccccccc @{}}
    \toprule
    \multirow{2}{*}{Method} & \multirow{2}{*}{\makecell{Speedup (s) $\uparrow$}} & \multirow{2}{*}{\makecell{TFLOPs $\downarrow$}} & Overall & \multicolumn{2}{c}{Image quality} & \multicolumn{2}{c}{Text alignment} \\
    \cmidrule(lr){4-4} \cmidrule(lr){5-6} \cmidrule(lr){7-8}
     & & & ImageReward $\uparrow$ & CLIP-IQA $\uparrow$ & NIQE $\downarrow$ & T2I-Comp. $\uparrow$ & GenEval $\uparrow$ \\
    \midrule
    FLUX (50 steps)             & 1.00$\times$              & 2991.01              & \textbf{1.095}    & 0.707             & 6.75             & \textbf{0.634}    & \textbf{0.698} \\
    RALU~\citep{ralu}     & \underline{1.58}$\times$              & \textbf{1749.94}  & 1.028             & \underline{0.712} & \textbf{6.07}    & 0.613                 & 0.648     \\
    \textbf{Ours ($S=2$)} & \textbf{1.66}$\times$  & \underline{1755.22}              & \underline{1.049} & \textbf{0.719}    & \underline{6.43}             & \underline{0.617} & \underline{0.654}     \\
    % \textbf{Ours ($S=3$)} & \textbf{1.73}$\times$     & \textbf{1.79}$\times$     & 1.015             & 0.711             & \underline{6.33} & 0.593             & 0.640     \\
    \midrule[0.3pt]
    FLUX (10 steps) & 4.84$\times$ & 610.02 & 0.981 & 0.679 & 6.93 & 0.618 & 0.647 \\
    Bottleneck~\citep{bottleneck} & 4.67$\times$              & 571.23              & 0.889          & 0.661             & 9.16          & 0.620             & \textbf{0.687}    \\
    RALU~\citep{ralu}             & 4.98$\times$              & 540.47              & 1.022          & \underline{0.700} & \textbf{6.43} & \textbf{0.626}    & 0.652             \\
    \textbf{Ours ($S=2$)}         & \underline{5.77}$\times$  & \underline{500.34}  & \textbf{1.059} & 0.696             & 6.69          & \underline{0.624} & \underline{0.655} \\
    \textbf{Ours ($S=3$)}         & \textbf{6.09}$\times$     & \textbf{469.15}     & \underline{1.042} & \textbf{0.701} & \underline{6.53} & 0.623          & 0.637             \\
    \midrule[0.3pt]
    FLUX (7 steps) & 6.62$\times$ & 431.45 & 0.920 & 0.660 & 8.25 & 0.594 & 0.583 \\
    Bottleneck~\citep{bottleneck} & 6.64$\times$              & 431.52              & 0.792          & 0.631             & 8.71          & 0.605             & \underline{0.672} \\
    RALU~\citep{ralu}             & 6.69$\times$              & \underline{426.01}  & 0.999          & 0.681             & 6.87          & \textbf{0.633}    & \textbf{0.682}    \\
    \textbf{Ours ($S=2$)}         & \underline{6.78}$\times$  & 427.03              & \textbf{1.039} & \underline{0.689} & \underline{6.78} & 0.620          & 0.667             \\
    \textbf{Ours ($S=3$)}         & \textbf{7.09}$\times$     & \textbf{406.24}     & \underline{1.015} & \textbf{0.694} & \textbf{5.99} & \underline{0.627} & 0.637             \\
    \bottomrule
  \end{tabular*}
  \vspace{-2pt} 
  \caption{\textbf{Training-free quantitative comparisons on FLUX.1-dev~\citep{black2024flux}.} Image resolution is $1024^2$ and baseline rows are copied from Table 1 in RALU~\citep{ralu}. We group results based on $1\times, 5\times, 7\times$ speedup settings where progressive resolution approaches are applied in combination with reducing denoising steps. Our method improves the speed--quality tradeoff across all speedup settings compared to baseline approaches. Further evaluation details and comparisons are shown in Appendix Sec.~\ref{subsec:image_latent}.}
  % replicating RALU's~\citep{ralu} Table~1 protocol; we keep only spatial-acceleration baselines.
  \label{tab:flux-main}
\end{table*}



\subsection{Optimal Resolution Schedule}

\label{sec:method:schedule:signal}

Section~\ref{sec:method:noise-injection} described our progressive resolution transition mechanism during the denoising process given a predetermined resolution transition schedule $t_{1:S-1}$. In practice, while increasing the number of resolution scales $S$ improves token efficiency, it introduces additional transition times as hyperparameters to tune. As a result, existing progressive resolution approaches typically limit $S = 2$ or leave $t_{1:S-1}$ as a brittle, model-specific design choice~\citep{bottleneck, ralu}. 
In contrast, our method supports an arbitrary number of resolutions $S$ and arbitrary resolution scales $s_{1:S}$. In the following section, we derive a $\delta$-\textit{optimal resolution schedule} for pretrained flow matching models that determines the optimal transition times $t^*_{1:S-1}$ from a single error threshold parameter $\delta$ independent of $S$.

%We include similar derivations for standard diffusion models in the supplemental materials.

\paragraph{Spectral-domain flow matching.}
The spatial-domain flow matching forward process defined in Eq.~\eqref{eq:flow-forward} can be written in the spectral domain under transformation $T_\Phi$ with orthonormal basis $\Phi = \{\phi_\omega\}_{\omega \in \Omega}$:
\begin{equation}
  x_t^{(\omega)} \;=\; (1-t)\,x_0^{(\omega)} \,+\, t\,\epsilon^{(\omega)},
  \qquad \omega \in \Omega,
  \label{eq:forward-mode}
\end{equation}
where $\epsilon^{(\omega)} := \langle \boldsymbol{\epsilon}, \phi_\omega \rangle$. By orthonormality of $\Phi$, $\epsilon^{(\omega)} \sim \mathcal{N}(0, 1)$ is i.i.d.\ Gaussian noise for each $\omega \in \Omega$. We denote the per-frequency component of the Bayes-optimal velocity predictor described in Sec. \ref{sec:method:prelim:basics} by $v^{*(\omega)}(\mathbf{x}_t, t)$. We then determine when a spectral component $\omega$ is noise-dominated in Proposition 1 and link this to spatial resolution to derive the $\delta$-optimal resolution schedule in Proposition 2. 

% We can then decompose the flow matching training objective defined in Eq.~\eqref{eq:flow-loss} into the sum of individual objectives for each $\omega \in \Omega$ using Parseval's identity~\citep{mallat1989wavelet}:
% \begin{equation}
%   \mathcal{L}_\text{flow}(\theta) \;=\; \sum_{\omega \in \Omega}\, \mathbb{E}_{t,\mathbf{x}_0,\mathbf{\epsilon}}\!\left[\,\bigl|v_\theta^{(\omega)}(\mathbf{x}_t, t) - (\epsilon^{(\omega)} - x_0^{(\omega)})\bigr|^2\,\right].
% \end{equation}

\textbf{Proposition 1 (Per-frequency $\delta$-optimal activation time):} Under the simplified modelling assumption $x_0^{(\omega)} \sim \mathcal{N}(0, P_\omega)$ with $P_\omega > 0$, for all $t \ge t_\omega$ and $\delta\in(0, 1)$, we have:
\begin{equation}
  \mathbb{E}\!\left[\bigl|v^{*(\omega)}(\mathbf{x}_t, t) - \epsilon^{(\omega)}\bigr|^2\right] \;\le\; \delta,
  \label{eq:noise-dominated}
\end{equation}
where $t_\omega$ is the $\delta$-optimal activation time for $\omega$:
\begin{equation}
  t_\omega \;:=\; \frac{1}{1 + \sqrt{\dfrac{\delta}{P_\omega\,(1 + P_\omega - \delta)}}}.
  \label{eq:activation}
\end{equation}
Intuitively, for $t \geq t_\omega$, spectral component $\omega$ is noise-dominated as its optimal velocity predictions are approximately noise $\epsilon^{(\omega)}$. After $t_\omega$, the signal power overcomes the error tolerance $\delta$ and structural information is recovered. The proof of Proposition 1 is given in Appendix Sec.~\ref{app:proof-prop1}. We further include empirical evidence supporting Proposition 1 in Appendix Sec.~\ref{sec:motivation}.


\textbf{Proposition 2 (Per-resolution $\delta$-optimal transition time):} Under the setting of Proposition~1, and assuming $P_\omega$ is monotonically decreasing in $|\omega|$ (consistent with the power-law decay Eq.~\eqref{eq:power-law}), for any resolution scale $s_i, s_{i + 1}$, the optimal transition time from scale $s_i$ up to $s_{i+1}$ is
\begin{equation}
  t_i^* \;:=\; \min_{\omega \in \Omega_{s_i}}\, t_\omega \;=\; t_{\omega=s_i \cdot \omega_{\max}(H,W)},
  \label{eq:stage-transition}
\end{equation}
where $\Omega_{s_i}$ is the set of frequencies representable on the $(s_iH, s_iW)$ grid. The maximum representable frequency of the full-resolution spatial grid $\Omega_{s_S} = \omega_{\max}(H, W) = \min(H, W)/2$ is given by the Nyquist--Shannon sampling theorem~\citep{nyquist1928certain,shannon1949communication}. 

Proposition 2 connects frequency activation times with spatial resolution transitions through the Nyquist limit. It implies that high-resolution computation in early denoising steps is redundant as most representable frequencies $\omega$ are still noise-dominated, motivating our approach of Spectral Progressive Diffusion. We apply spectral noise expansion at $t_i^*$ to progressively increase resolution precisely when finer details emerge from noise, thereby maximizing token efficiency. The hyperparameter $\delta$ is independent of $S$, making our method particularly robust to tuning compared to prior methods. The proof of Proposition 2 is deferred to Appendix~\ref{app:proof-prop2} and we analyze the effect of the number of resolution stages $S$ and the tolerance $\delta$ in Sec.~\ref{sec:exp:ablations}.

% In practice, Eq.~\eqref{eq:stage-transition} is discretized into $S$ resolution stages $(0 < s_1 < s_2 < \cdots < s_S = 1)$, with stage $i$ held for $t \in (t_{s_i}^{\text{transition}},\, t_{s_{i-1}}^{\text{transition}}]$ and matched to the denoising timesteps of the pretrained models.

% \paragraph{Per-frequency SNR.}
% Under the basis $\Phi$, the rectified-flow forward process in Eq.~\eqref{eq:flow-forward} factors into independent per-frequency equations in which the noise component $\epsilon^{(\omega)} \sim \mathcal{N}(0, 1)$ is i.i.d.\ Gaussian (full derivation in Sec.~\ref{sec:method:schedule:signal}). Following~\citet{kingma2021vdm}, we define the \emph{per-frequency signal-to-noise ratio} as the ratio of the clean- and noise-component second moments in this process:
% \begin{equation}
%   \mathrm{SNR}_\omega(t)
%   \;:=\;
%   \frac{\mathbb{E}[\,|(1-t)\,x_0^{(\omega)}|^2\,]}{\mathbb{E}[\,|t\,\epsilon^{(\omega)}|^2\,]}
%   \;=\;
%   \underbrace{\frac{(1-t)^2}{t^2}}_{\mathrm{SNR}(t)}\;\cdot\;P_\omega,
%   \label{eq:snr}
% \end{equation}
% which factorises into the frequency-independent timestep-SNR term $\mathrm{SNR}(t) = (1-t)^2/t^2$ and the frequency-specific signal power $P_\omega$.

% \paragraph{Frequency limit from resolution.}
% By the Nyquist--Shannon sampling theorem~\citep{nyquist1928certain,shannon1949communication}, when generating on a spatial grid of size $H \times W$ the representable 2D frequencies form the rectangle $|\omega_x| \le W/2$, $|\omega_y| \le H/2$. For the radial analysis used to derive our resolution schedule, we define the maximum representable radial frequency as
% \begin{equation}
%   \omega_{\max}(H, W) \;=\; \frac{\min(H, W)}{2}.
%   \label{eq:nyquist-cap}
% \end{equation}
% A stage running at reduced resolution $(sH, sW) \le (H, W)$ for $s \leq 1$ can represent only the frequency subset
% \begin{equation}
%   \Omega_s \;:=\; \{\,\omega \in \Omega : |\omega| \le s \cdot \omega_{\max}(H, W)\,\} \;\subseteq\; \Omega.
%   \label{eq:nyquist-set}
% \end{equation}



% To characterize how much useful signal each frequency carries at time $t$, we measure how far $v^*$ deviates from the null prediction $\hat v = \epsilon$. Under the simplified modelling assumption that $x_0^{(\omega)} \sim \mathcal{N}(0, P_\omega)$, we get:
% \begin{equation}
%   \mathbb{E}\!\left[\bigl|v^{*(\omega)} - \epsilon^{(\omega)}\bigr|^2\right]
%   \;=\; \frac{\mathrm{SNR}_\omega(t)\,(1 + P_\omega)}{1 + \mathrm{SNR}_\omega(t)},
%   \label{eq:velocity-error-optimal}
% \end{equation}
% We defer the derivations of above to Appendix~\ref{app:bayes-velocity}. 

% \paragraph{Noise-dominated frequencies.}
% Given a tolerance $\delta > 0$, we say frequency $\omega$ is \emph{noise-dominated at time $t$} if the Bayes-optimal velocity deviates from the null prediction by at most $\delta$ in expectation:
% \begin{equation}
%   \mathbb{E}\!\left[\bigl|v^{*(\omega)} - \epsilon^{(\omega)}\bigr|^2\right] \;\le\; \delta.
%   \label{eq:noise-dominated}
% \end{equation}
% Substituting the Bayes-optimal error of Eq.~\eqref{eq:velocity-error-optimal} and the SNR factorisation from Eq.~\eqref{eq:snr} into this definition, and solving for $t$, turns the noise-dominated condition into an inequality on the timestep:
% \begin{equation}
%   t \;\ge\; \frac{1}{1 + \sqrt{\dfrac{\delta}{P_\omega\,(1 + P_\omega - \delta)}}}
%   \;=:\; t_\omega,
%   \label{eq:activation}
% \end{equation}
% defining the frequency-specific \emph{activation time} $t_\omega$ as the right-hand side. We assume $\delta < 1$, and the right hand side of Eq.~\eqref{eq:activation} is well defined with $0 < t_\omega < 1$. Frequency $\omega$ is thus noise-dominated for $t \ge t_\omega$ and transitions to signal-carrying for $t < t_\omega$ under tolerance $\delta$.

% \paragraph{$\delta$-optimum resolution schedule.}
% A \emph{resolution schedule} assigns each timestep $t$ a resolution scale $s_t \in (0, 1]$ such that the generator runs on the grid $(s_t H,\, s_t W)$ and represents the frequencies in $\Omega_{s_t}$ of Eq.~\eqref{eq:nyquist-set}. We call a schedule \emph{$\delta$-optimum} if at every $t$ it uses the smallest scale that still represents every frequency already signal-carrying under tolerance $\delta$. Formally,
% \begin{equation}
%   s^*_t \;=\; \min\,\bigl\{\,s \in (0, 1] : \{\omega \in \Omega : t_\omega > t\} \subseteq \Omega_s\,\bigr\}
%   \;=\; \frac{\omega^*(t)}{\omega_{\max}(H, W)},
%   \label{eq:delta-opt-schedule}
% \end{equation}
% where $\omega^*(t) := \max\{|\omega| : t_\omega > t\}$ is the highest frequency that has crossed its activation time $t_\omega$ (i.e., frequencies with $t < t_\omega$ are already signal-carrying at time $t$).

% Equivalently, the transition time from resolution scale $s$ up to a finer resolution scale $s' > s$ is exactly when every $\omega \in \Omega_s$ has become signal-carrying, i.e., when $t$ falls to the smallest activation time within $\Omega_s$:
% \begin{equation}
%   t_s^{\text{transition}} \;:=\; \min_{\omega \in \Omega_s}\, t_\omega \;=\; t_{\,s \cdot \omega_{\max}(H,W)},
%   \label{eq:stage-transition}
% \end{equation}
% where the second equality follows from the fact that $t_\omega$ is decreasing in $|\omega|$ as $P_\omega$ decays with $|\omega|$ (Eq.~\eqref{eq:power-law}), and Eq.~\eqref{eq:activation} is monotonically increasing in $P_\omega$.

\renewcommand{\arraystretch}{1.0}
\setlength{\doublerulesep}{0.5pt}
\setlength{\aboverulesep}{0.2ex}
\begin{table*}[!t]
  \centering
  \scriptsize
  \begin{tabular*}{\textwidth}{@{\extracolsep{\fill}} l ccccccc @{}}
    \toprule
    Method & \makecell{Speedup (s) $\uparrow$} & \makecell{TFLOPs $\downarrow$} & ImageReward $\uparrow$ & CLIP-IQA $\uparrow$ & NIQE $\downarrow$ & T2I-Comp. $\uparrow$ & GenEval $\uparrow$ \\
    \midrule
    Z-Image (50 steps)           & 1.00$\times$             & 4941.23             & 0.965    & 0.700    & 5.41    & 0.731    & 0.745 \\
    % Z-Image (32 steps)           & {1.56}$\times$             & {3166.62}             & 0.957    & 0.697    & {5.44}    & 0.686    & 0.725 \\
    \midrule
    \textbf{Ours} (TF, $S=2$)  & {1.65}$\times$ & {3132.03} & {0.904}             & {0.688}             & {5.87}             & {0.658} & {0.730} \\
    % \textbf{Ours ($S=3$)}  & \textbf{1.74}$\times$    & \textbf{1.72}$\times$    & 0.875             & 0.690             & {5.59} & 0.650             & 0.682 \\
    \textbf{Ours} (LoRA, $S = 2$)   & {1.65}$\times$             & {3132.03}             & \textbf{0.982} & \textbf{0.699} & \textbf{5.72}             & \textbf{0.725}                 & \textbf{0.731} \\
    \midrule
    \textbf{Ours} (TF, $S=3$)  & {1.74}$\times$ & {2871.09} & {0.875}             & {0.690}             & \textbf{5.59}             & {0.650} & {0.682} \\
    % \textbf{Ours ($S=3$)}  & \textbf{1.74}$\times$    & \textbf{1.72}$\times$    & 0.875             & 0.690             & {5.59} & 0.650             & 0.682 \\
    \textbf{Ours} (LoRA, $S = 3$)   & {1.74}$\times$             & {2871.09}             & \textbf{0.954} & \textbf{0.697} & {5.75}             & \textbf{0.717}                 & \textbf{0.728} \\
    % \textbf{Ours} (LoRA, $\delta=0.05$)   & \textbf{2.01}$\times$             & \textbf{2436.19}             & {0.919} & 0.684 & 6.12             & {0.661}                 & 0.718 \\
    % \midrule
    % \textbf{Ours} (TF, $S=3$)  & \textbf{1.74}$\times$ & \textbf{2871.09} & 0.875            & 0.690            & 5.59             & 0.650 & 0.682 \\
    % \textbf{Ours} (LoRA, $S=3$)  & \textbf{1.74}$\times$ & \textbf{2871.09} & \textbf{0.954}             &  \textbf{0.697}            & 5.75             & - & - \\

    % \textbf{Ours}
    \midrule
    \midrule
    % \textbf{Ours} (TF, $S=3$)  & \textbf{1.74}$\times$ & \textbf{2871.09} & 0.904             & 0.688             & 5.87             & {0.658} & {0.730} \\
    % % \textbf{Ours ($S=3$)}  & \textbf{1.74}$\times$    & \textbf{1.72}$\times$    & 0.875             & 0.690             & {5.59} & 0.650             & 0.682 \\
    % \textbf{Ours} (LoRA, $S=3$)   & \textbf{1.74}$\times$             & \textbf{2871.09}             & \textbf{0.982} & {0.699} & \textbf{5.72}             & \hx{TODO}                 & \hx{TODO} \\
    PixelGen (25 steps)         & 1.00$\times$             & 65.36             & 0.921    & 0.734    & 5.95    & 0.574    & 0.794 \\
    % PixelGen (13 steps)         & \textbf{1.92}$\times$             & {34.16}             & 0.886    & 0.726    & {5.89}    & 0.571    & 0.781 \\
    \midrule
    \textbf{Ours} (TF, $S=2$) & {1.60}$\times$             & {33.72}             & {0.799}    & {0.718}    & {6.10}    &  {0.568}  & \textbf{0.782} \\
    \textbf{Ours} (LoRA, $S=2$) & {1.55}$\times$             & {33.72}             & \textbf{0.913}    & \textbf{0.728}    & \textbf{5.87}    & \textbf{0.580}    & {0.776} \\
    \bottomrule
  \end{tabular*}
  % \vspace{3pt}
  \caption{\textbf{Fine-tuning quantitative comparisons on latent- (Z-Image~\citep{zimage}) and pixel-space image generation (PixelGen~\citep{pixelgen}).} Image resolution is $1024^2$ for Z-Image and $512^2$ for PixelGen. Across both latent- and pixel-space image generation, our fine-tuning method bridges the gap between model pretraining and progressive resolution inference, further improving generation quality while preserving efficiency gains. Further evaluation details and comparisons are shown in Appendix Sec.~\ref{subsec:image_pixel}.}
  \label{tab:zimage-main}
\end{table*}

\renewcommand{\arraystretch}{1.0}
\setlength{\aboverulesep}{0.2ex}
\begin{table*}[!t]
  \centering
  \scriptsize
  \begin{tabular*}{\textwidth}{@{\extracolsep{\fill}} l cccccccc @{}}
    \toprule
    \makecell{Method} &
    \makecell{Speedup (s) $\uparrow$} &
    \makecell{TFLOPs $\downarrow$} &
    \makecell{Subject\\Consistency $\uparrow$} &
    \makecell{Background\\Consistency $\uparrow$} &
    \makecell{Motion\\Smoothness $\uparrow$} &
    \makecell{Dynamic\\Degree $\uparrow$} &
    \makecell{Aesthetic\\Quality $\uparrow$} &
    \makecell{Image\\Quality $\uparrow$} \\
    \midrule
    WAN~2.1 (50)           & 1.00$\times$ & 119292         & 0.9492          & 0.9621          & 0.9874          & 0.4800          & 0.5993          & 0.6133          \\
    \midrule[0.3pt]
    WAN~2.1 (25)           & 2.00$\times$ & 59646          & 0.9434          & 0.9597          & \textbf{0.9879} & 0.4250          & 0.5893          & 0.5706          \\
    \textbf{Ours ($S=2$)} & \textbf{2.03}$\times$ & \textbf{57417} & \textbf{0.9462} & \textbf{0.9598} & 0.9859          & \textbf{0.4950} & \textbf{0.5975} & \textbf{0.6114} \\
    \midrule[0.3pt]
    \textbf{Ours ($S=3$)} & 2.54$\times$ & 45953 & 0.9299          & 0.9499          & 0.9815          & 0.5700 & 0.5696          & 0.5990 \\
    \bottomrule
  \end{tabular*}
  \caption{\textbf{Quantitative comparison on latent-space video generation on WAN~2.1~\citep{wan}.} Video resolution is maintained at 720P. Our training-free approach demonstrates more than $2\times$ speedup while maintaining high generation quality and outperforming the 25-step full-resolution baseline.}
  \label{tab:wan_vbench_720p}
\end{table*}




% \renewcommand{\arraystretch}{1.0}
% \setlength{\aboverulesep}{0.2ex}
% \begin{table*}[!t]
%   \centering
%   \scriptsize
%   \begin{tabular*}{\textwidth}{@{\extracolsep{\fill}} l ccccccc @{}}
%     \toprule
%     \makecell{Method} &
%     \makecell{Speedup $\uparrow$} &
%     \makecell{Subject\\Consistency $\uparrow$} &
%     \makecell{Background\\Consistency $\uparrow$} &
%     \makecell{Motion\\Smoothness $\uparrow$} &
%     \makecell{Dynamic\\Degree $\uparrow$} &
%     \makecell{Aesthetic\\Quality $\uparrow$} &
%     \makecell{Image\\Quality $\uparrow$} \\
%     \midrule
%     WAN2.1 (50 steps)           & 1.00$\times$ & 0.9492          & 0.9621          & 0.9874 & 0.4800          & 0.5993          & 0.6133          \\
%     \midrule[0.3pt]
%     WAN2.1 (25 steps)           & 2.00$\times$ & 0.9434          & 0.9597          & \textbf{0.9879} & 0.4250          & 0.5893          & 0.5706          \\
%     \textbf{Ours ($S=2$)} & 2.03$\times$ & \textbf{0.9462} & \textbf{0.9598} & 0.9859 & \textbf{0.4950} & \textbf{0.5975} & \textbf{0.6114} \\
%     \midrule[0.3pt]
%     \textbf{Ours ($S=3$)} & 2.54$\times$ & 0.9299          & 0.9499          & 0.9815 & 0.5700 & 0.5696          & 0.5990 \\
%     \bottomrule
%   \end{tabular*}
%   \vspace{3pt}
%   \caption{\textbf{Quantitative comparison on latent-space video generation on WAN2.1~\citep{wan}.} Video resolution is maintained at 720P. Our training-free approach demonstrates more than $2\times$ speedup while maintaining high generation quality and outperforming the 25-step full-resolution baseline.}
%   \label{tab:wan_vbench_720p}
% \end{table*}


\subsection{Spectral-transformation-based Fine-tuning}
\label{sec:method:lora}
While Spectral Progressive Diffusion is compatible with pretrained image and video generation models in a training-free manner (Sec.~\ref{sec:method:noise-injection}), it implicitly assumes multi-resolution generation capability of pretrained models across scales $s_{1:S}$. Moreover, our optimal resolution schedule (Sec.~\ref{sec:method:schedule:signal}) may not accurately reflect pretrained models' training dynamics, potentially introducing training-inference gaps. 
In this subsection, we outline a spectral-transformation-based fine-tuning approach that directly follows our Spectral Progressive Diffusion framework and our optimal resolution schedule. We include additional implementation details of our fine-tuning framework in Sec.~\ref{sec:exp:image} and Appendix Sec.~\ref{subsec:image_pixel}.

\paragraph{Resolution-specific velocity targets.} 
We assume $S$ progressive resolution scales $s_{1:S}$ and their corresponding $\delta$-optimal resolution transition times $t_{1:S-1}$ from Sec.~\ref{sec:method:schedule:signal} (denoted here as $t_i$ instead of $t^*_i$ for simplicity). For resolution stage $i$, the model operates at scale $s_i$ for $t \in (t_i, \tilde{t}_{i-1}]$, where $\tilde{t}_{i-1}$ is the aligned timestep obtained from Eq.~\eqref{eq:time-align} after expanding from $s_{i-1}$ to $s_i$ at transition time $t_{i-1}$. The state at the beginning of the stage $i$ is hence the enlarged and timestep-aligned $\tilde{\mathbf{x}}_{\tilde{t}_{i - 1}}^{s_{i}}$. The model is fine-tuned to follow a straight-line path from $\tilde{\mathbf{x}}_{\tilde{t}_{i - 1}}^{s_{i}}$ to the standard flow-matching state at the next transition time:
\begin{equation}
  \mathbf{x}_{t_i}^{s_i}
  =
  (1-t_i)\mathbf{x}_0^{s_i}
  +
  t_i\boldsymbol{\epsilon}^{s_i}.
  \label{eq:lora-stagewise}
\end{equation}
where $\mathbf{x}_0^{s_i}$ denotes the clean data sample at scale $s_i$.
The corresponding stage-specific velocity target is therefore:
\begin{equation}
  \mathbf{v}^{s_i}
  =
  \frac{
  \tilde{\mathbf{x}}_{\tilde{t}_{i - 1}}^{s_{i}}
  -
  \mathbf{x}_{t_i}^{s_i}
  }{
  \tilde{t}_{i - 1} - t_i
  }.
  \label{eq:lora-stagewise-velocity}
\end{equation}

\paragraph{Training and inference.}
At each training step we sample $\mathbf{x}_0 \sim p_\text{data}$ and $t \sim \mathcal{U}(0,1)$. We assign $t$ to the resolution stage $i$ satisfying $t \in (t_i,t_{i-1}]$, and then sample $\boldsymbol{\epsilon}^{s_i} \sim \mathcal{N}(0, I)$ at scale $s_i$. We construct the stage input $\tilde{\mathbf{x}}_{\tilde{t}_{i - 1}}^{s_{i}}$ by
applying the same spectral noise expansion and timestep alignment used at
inference, where newly introduced spectral coefficients are filled using
$t_{i-1} \cdot T_\Phi(\boldsymbol{\epsilon}^{s_i})$. We use the same
$\boldsymbol{\epsilon}^{s_i}$ to construct the endpoint
$\mathbf{x}^{s_i}_{t_i}$ from Eq.~\eqref{eq:lora-stagewise}, so that $\tilde{\mathbf{x}}_{\tilde{t}_{i - 1}}^{s_{i}}$ and $\mathbf{x}^{s_i}_{t_i}$ are correlated through a shared noise realization. The training sample is the point on the straight path between them:
\begin{equation}
  \mathbf{x}_t^{s_i}
  =
  \mathbf{x}_{t_i}^{s_i}
  +
  (t-t_i)\,\mathbf{v}^{s_i},
\end{equation}
with target velocity $\mathbf{v}^{s_i}$ from
Eq.~\eqref{eq:lora-stagewise-velocity}.

We fine-tune a pretrained model $\mathbf{v}_\theta$ by minimizing
\begin{equation}
  \mathcal{L}(\theta) \;=\; \mathbb{E}_{t,\, \mathbf{x}_0,\, \boldsymbol{\epsilon}}\!\left[\,\|\,\mathbf{v}_\theta(\mathbf{x}_t^{s_i}, t) - \mathbf{v}^{s_i}\,\|^2\,\right].
  \label{eq:lora-loss}
\end{equation}
At inference, we use the same procedure as in Sec.~\ref{sec:method:noise-injection} with the fine-tuned model, using the same resolution schedule $t_{1:S-1}$ and transform $T_\Phi$.

% !TEX root = ../main.tex
\section{Experiments}
\label{sec:exp}

% \newcommand{\placeholderfig}[1]{\parbox[c][#1][c]{0.95\linewidth}{\centering\includegraphics[width=0.95\linewidth,height=#1,keepaspectratio]{example-image-a}}}


\begin{figure}[!t]
  \centering
  \includegraphics[width=\linewidth]{figures/main_qual.pdf}
  \caption{\textbf{Visual Generation Qualitative Comparisons.} For the main comparison of latent-space image generation, our method outperforms the state-of-the-art spatial acceleration method RALU \cite{ralu} in both visual fidelity and inference speed. Across all evaluated modalities (latent/pixel-space image generation and latent-space video generation), we achieve substantial acceleration over standard high-resolution baselines while preserving generation quality.}
  \label{fig:qual}
\end{figure}


\subsection{Setup}
\label{sec:exp:setup}

We extensively evaluate our Spectral Progressive Diffusion framework on latent-space image generation with \textbf{FLUX.1-dev}~\citep{black2024flux} and \textbf{Z-Image}~\citep{zimage}, covering both training-free inference acceleration and fine-tuning with LoRA~\citep{hu2021lora}. We further demonstrate that our framework applies to pixel-space image generation with \textbf{PixelGen-XXL/16 T2I}~\citep{pixelgen} and latent-space video generation with \textbf{WAN~2.1-T2V-1.3B}~\citep{wan}. Each pretrained backbone uses its native resolution and default model-specific inference schedule.

For latent- and pixel-space image generation, we follow RALU's~\citep{ralu} evaluation protocol, reporting ImageReward~\citep{imagereward} for overall quality, CLIP-IQA~\citep{wang2023clipiqa} and NIQE~\citep{mittal2013niqe} for image quality together with T2I-CompBench~\citep{t2icompbench} and GenEval~\citep{geneval} for prompt alignment. For video generation, we report VBench~\citep{vbench} scores across its standard evaluation dimensions. We report total FLOPs integrated over the full denoising trajectory using the same convention as RALU and end-to-end normalized wall-clock speedup. We use $\delta=0.01$ and Discrete Cosine Transform (DCT) as our spectral transformation $T_\Phi$ across the experiments and include ablation studies on $\delta$, $S$, and $T_\Phi$ in Sec.~\ref{sec:exp:ablations}. We include additional implementation and experimental details in Appendix Secs.~\ref{sec:image_generation} and~\ref{sec:video_generation}.

% \hx{commented out per Howard -- configuration / activation-threshold table deferred}
% Table~\ref{tab:setup} summarises the per-backbone configuration together with the activation thresholds $t_{0}$ (Eq.~\eqref{eq:activation}) induced by the default tolerance $\delta$.
%
% \begin{table}[t]
%   \centering
%   \small
%   \caption{\textbf{Per-backbone configuration and activation thresholds} at the default tolerance $\delta$. \textit{Res.\ ratio} is the spatial downsampling between adjacent stages; \textit{Act.\ time $t_{0}$} and \textit{Act.\ step} follow from Eq.~\eqref{eq:activation}. Placeholder values (``-'') to be filled.}
%   \label{tab:setup}
%   \setlength{\tabcolsep}{5pt}
%   \begin{tabular}{l c c c c c c c c}
%     \toprule
%     Backbone & Resolution & Steps & Shift & $\delta$ & Basis & Res.\ ratio & $t_{0}$ & Act.\ step \\
%     \midrule
%     PixelGen         & - & - & - & - & - & - & - & - \\
%     SD3              & - & - & - & - & - & - & - & - \\
%     FLUX.1-dev       & - & - & - & - & - & - & - & - \\
%     WAN 2.1-T2V-1.3B & - & - & - & - & - & - & - & - \\
%     \bottomrule
%   \end{tabular}
% \end{table}

% \iffalse
% \subsection{Empirical Validation of Power Spectrum and Resolution Schedule}
% \label{sec:exp:validation}

% \bc{I think this whole section can be moved to the supplement, and move Fig.2 to preliminaries}

% We validate the two theoretical pillars of Section~\ref{sec:method:schedule} on our target backbones: (i) the power-law decay of the per-frequency signal power, and (ii) the prediction of Eq.~\eqref{eq:velocity-error-optimal} for when a pretrained velocity network actually starts to deviate from the input noise on each frequency. Figure~\ref{fig:wavelet-spectral} shows both.

% \paragraph{Power spectra of natural images and latents.}
% The left column of Fig.~\ref{fig:wavelet-spectral} overlays the radial power spectrum $P(|\omega|)$ measured on natural images, SD3 latents, FLUX.1-dev latents, and WAN~2.1 latents, with error bars across latent channels. All four curves admit a clean power-law fit $P(|\omega|) \propto |\omega|^{-\beta}$ with the exponent $\beta$ reported inline.

% \paragraph{Per-frequency signal content in pretrained denoisers.}
% The middle column of Fig.~\ref{fig:wavelet-spectral} measures, for a \emph{pretrained} velocity network $v_\theta$ without any modifications, the per-frequency squared deviation $\mathbb{E}\!\left[\bigl|v_\theta^{(\omega)}(x_t, t) - \epsilon^{(\omega)}\bigr|^2\right]$ as a function of timestep $t$, stratified by Haar sub-band on FLUX.1-dev and WAN~2.1. Comparing the model output against $\epsilon$ itself (rather than against the regression target $\epsilon - x_0$) directly probes whether the model is extracting any useful signal on that frequency: if the model simply outputs the input noise, the curve stays at zero. The measured curves are overlaid with the analytical prediction of Eq.~\eqref{eq:velocity-error-optimal} using the $\{P_\omega\}$ measured in the left column. Pretrained denoisers closely track the analytical curve: each sub-band stays at zero deviation while $\mathrm{SNR}_\omega(t)$ is small and rises sharply to $\sim 1 + P_\omega$ once $\mathrm{SNR}_\omega(t)$ crosses unity, confirming that noise-dominated frequencies are genuinely uninformative at their active times.

% % !TEX root = ../main.tex
\begin{figure}[t]
  \centering
  \includegraphics[width=\linewidth]{figures/power_spectrum.pdf}
  \caption{\textbf{Diffusion process in the spectral domain.} Latent power spectra in both image and video models decay rapidly with frequency (Fig.~\textbf{(a)}), consistent with natural images. Diffusion exhibits a frequency-domain autoregressive structure (Fig.~\textbf{(b)}) due to the aforementioned property: low frequencies emerge early in the denoising process, while high frequencies remain noise-dominated.}
  \label{fig:wavelet-spectral}
\end{figure}


% Further empirical validation comes from the spectral-passthrough experiment reported in Section~\ref{sec:exp:image} (Fig.~\ref{fig:passthrough} and the ``spectral passthrough'' rows of Tables~\ref{tab:flux-main} and~\ref{tab:wan-main}): replacing noise-dominated sub-bands with fresh Gaussian noise at every denoising step leaves generation quality essentially unchanged, confirming that at their active times these bands carry no information beyond noise.
% \fi

\vspace{-5pt}
\subsection{Latent- and Pixel-Space Image Generation}
\vspace{-5pt}
\label{sec:exp:image}

% 
\begin{table*}[t!]
    \centering
    \caption{\textbf{Extended quantitative comparisons on training-free latent-space image generation.} We evaluate FLUX.1-dev at $1024^2$ resolution and group methods into the same $1\times$, $5\times$, and $7\times$ wall-clock speedup tiers as Table~\ref{tab:flux-main}. Rows shared with Table~\ref{tab:flux-main} use the same values. Speedup is calculated relative to the 50-step FLUX.1-dev baseline wall-clock runtime on a single A100 GPU. \textcolor{blue}{\textbf{T}} and \textcolor{red}{\textbf{S}} denote temporal and spatial acceleration methods, respectively. Baseline metrics are copied from RALU~\citep{ralu} under the same evaluation protocol.}
    \label{tab:flux-extended}
    \resizebox{\textwidth}{!}{
    \begin{scriptsize}
    \begin{tabular}{cccccccccc}
         \toprule
        \multirow{2}{*}{Method} & \multirow{2}{*}{Accel.} & \multirow{2}{*}{Latency (s) $\downarrow$} & \multirow{2}{*}{Speedup (s) $\uparrow$} & \multirow{2}{*}{TFLOPs $\downarrow$} & Overall & \multicolumn{2}{c}{Image quality} & \multicolumn{2}{c}{Text alignment} \\
        \cmidrule(lr){6-10}
        & & & & & ImageReward $\uparrow$ & CLIP-IQA $\uparrow$ & NIQE $\downarrow$ & T2I-Comp. $\uparrow$ & GenEval $\uparrow$  \\
        \midrule
        FLUX (50 steps) & - & 25.1 & 1.00$\times$ & 2991.01 & \textbf{1.095} & 0.707 & 6.75 & \textbf{0.634} & \textbf{0.698}  \\
        RALU~\citep{ralu}   & \textcolor{red}{\textbf{S}}  & 15.85 & 1.58$\times$ & \underline{1749.94} & 1.028 & \underline{0.712} & \textbf{6.07} & 0.613 & 0.648 \\
        \textbf{Ours ($S=2$)}     & \textcolor{red}{\textbf{S}}       & \underline{15.16} & \underline{1.66}$\times$ & 1755.22 & \underline{1.049} & \textbf{0.719} & 6.43 & \underline{0.617} & \underline{0.654} \\
        \textbf{Ours ($S=3$)}     & \textcolor{red}{\textbf{S}}       & \textbf{14.54} & \textbf{1.73}$\times$ & \textbf{1672.04} & 1.015 & 0.711 & \underline{6.33} & 0.593 & 0.640 \\
        \cmidrule(lr){1-10}
        FLUX (10 steps) & \textcolor{blue}{\textbf{T}} & 5.18 & 4.84$\times$ & 610.02 & 0.981 & 0.679 & 6.93 & 0.618 & 0.647 \\
        $\Delta$-DiT~\citep{deltadit}      & \textcolor{blue}{\textbf{T}} & 7.42 & 3.38$\times$ & 772.10 & 0.102 & 0.487 & 9.60 & 0.306 & 0.397 \\
        ToCa~\citep{toca}                  & \textcolor{blue}{\textbf{T}} & 15.5 & 1.62$\times$ & 601.12 & -1.827 & 0.253 & 10.6 & 0.259 & 0.137 \\
        TeaCache~\citep{teacache}          & \textcolor{blue}{\textbf{T}} & 5.23 & 4.80$\times$ & 610.59 & 0.944 & 0.665 & 7.92 & 0.620 & 0.647 \\
        TaylorSeer~\citep{taylorseer}      & \textcolor{blue}{\textbf{T}} & 9.34 & 2.69$\times$ & 556.72 & 0.972 & 0.684 & 6.77 & 0.594 & 0.619 \\
        Bottleneck~\citep{bottleneck}      & \textcolor{red}{\textbf{S}}  & 5.37 & 4.67$\times$ & 571.23 & 0.889 & 0.661 & 9.16 & 0.620 & \textbf{0.687} \\
        RALU~\citep{ralu}                  & \textcolor{red}{\textbf{S}}  & 5.04 & 4.98$\times$ & 540.47 & 1.022 & \underline{0.700} & \textbf{6.43} & \textbf{0.626} & 0.652 \\
        \textbf{Ours ($S=2$)}     & \textcolor{red}{\textbf{S}}       & \underline{4.35} & \underline{5.77}$\times$ & \underline{500.34} & \textbf{1.059} & 0.696 & 6.69 & \underline{0.624} & \underline{0.655} \\
        \textbf{Ours ($S=3$)}     & \textcolor{red}{\textbf{S}}       & \textbf{4.12} & \textbf{6.09}$\times$ & \textbf{469.15} & \underline{1.042} & \textbf{0.701} & \underline{6.53} & 0.623 & 0.637 \\
        \cmidrule(lr){1-10}
        FLUX (7 steps) & \textcolor{blue}{\textbf{T}} & 3.79 & 6.62$\times$ & 431.45 & 0.920 & 0.660 & 8.25 & 0.594 & 0.583 \\
        TeaCache~\citep{teacache}          & \textcolor{blue}{\textbf{T}} & 4.21 & 5.96$\times$ & 431.83 & 0.733 & 0.623 & 13.7 & 0.599 & 0.594 \\
        TaylorSeer~\citep{taylorseer}      & \textcolor{blue}{\textbf{T}} & 7.00 & 3.59$\times$ & 431.74 & 0.660 & 0.646 & 9.43 & 0.514 & 0.446 \\
        Bottleneck~\citep{bottleneck}      & \textcolor{red}{\textbf{S}}  & 3.78 & 6.64$\times$ & 431.52 & 0.792 & 0.631 & 8.71 & 0.605 & \underline{0.672} \\
        RALU~\citep{ralu}                  & \textcolor{red}{\textbf{S}}  & 3.75 & 6.69$\times$ & \underline{426.01} & 0.999 & 0.681 & 6.87 & \textbf{0.633} & \textbf{0.682} \\
        \textbf{Ours ($S=2$)}     & \textcolor{red}{\textbf{S}}       & \underline{3.70} & \underline{6.78}$\times$ & 427.03 & \textbf{1.039} & \underline{0.689} & \underline{6.78} & 0.620 & 0.667 \\
        \textbf{Ours ($S=3$)}     & \textcolor{red}{\textbf{S}}       & \textbf{3.54} & \textbf{7.09}$\times$ & \textbf{406.24} & \underline{1.015} & \textbf{0.694} & \textbf{5.99} & \underline{0.627} & 0.637 \\
        \bottomrule
    \end{tabular}
    \end{scriptsize}}
\end{table*}

% \renewcommand{\arraystretch}{1.0}
\setlength{\aboverulesep}{0.2ex}
\begin{table*}[!t]
  \centering
  \scriptsize
  \begin{tabular*}{\textwidth}{@{\extracolsep{\fill}} l ccccccc @{}}
    \toprule
    \multirow{2}{*}{Method} & \multirow{2}{*}{\makecell{Speedup (s) $\uparrow$}} & \multirow{2}{*}{\makecell{TFLOPs $\downarrow$}} & Overall & \multicolumn{2}{c}{Image quality} & \multicolumn{2}{c}{Text alignment} \\
    \cmidrule(lr){4-4} \cmidrule(lr){5-6} \cmidrule(lr){7-8}
     & & & ImageReward $\uparrow$ & CLIP-IQA $\uparrow$ & NIQE $\downarrow$ & T2I-Comp. $\uparrow$ & GenEval $\uparrow$ \\
    \midrule
    FLUX (50 steps)             & 1.00$\times$              & 2991.01              & \textbf{1.095}    & 0.707             & 6.75             & \textbf{0.634}    & \textbf{0.698} \\
    RALU~\citep{ralu}     & \underline{1.58}$\times$              & \textbf{1749.94}  & 1.028             & \underline{0.712} & \textbf{6.07}    & 0.613                 & 0.648     \\
    \textbf{Ours ($S=2$)} & \textbf{1.66}$\times$  & \underline{1755.22}              & \underline{1.049} & \textbf{0.719}    & \underline{6.43}             & \underline{0.617} & \underline{0.654}     \\
    % \textbf{Ours ($S=3$)} & \textbf{1.73}$\times$     & \textbf{1.79}$\times$     & 1.015             & 0.711             & \underline{6.33} & 0.593             & 0.640     \\
    \midrule[0.3pt]
    FLUX (10 steps) & 4.84$\times$ & 610.02 & 0.981 & 0.679 & 6.93 & 0.618 & 0.647 \\
    Bottleneck~\citep{bottleneck} & 4.67$\times$              & 571.23              & 0.889          & 0.661             & 9.16          & 0.620             & \textbf{0.687}    \\
    RALU~\citep{ralu}             & 4.98$\times$              & 540.47              & 1.022          & \underline{0.700} & \textbf{6.43} & \textbf{0.626}    & 0.652             \\
    \textbf{Ours ($S=2$)}         & \underline{5.77}$\times$  & \underline{500.34}  & \textbf{1.059} & 0.696             & 6.69          & \underline{0.624} & \underline{0.655} \\
    \textbf{Ours ($S=3$)}         & \textbf{6.09}$\times$     & \textbf{469.15}     & \underline{1.042} & \textbf{0.701} & \underline{6.53} & 0.623          & 0.637             \\
    \midrule[0.3pt]
    FLUX (7 steps) & 6.62$\times$ & 431.45 & 0.920 & 0.660 & 8.25 & 0.594 & 0.583 \\
    Bottleneck~\citep{bottleneck} & 6.64$\times$              & 431.52              & 0.792          & 0.631             & 8.71          & 0.605             & \underline{0.672} \\
    RALU~\citep{ralu}             & 6.69$\times$              & \underline{426.01}  & 0.999          & 0.681             & 6.87          & \textbf{0.633}    & \textbf{0.682}    \\
    \textbf{Ours ($S=2$)}         & \underline{6.78}$\times$  & 427.03              & \textbf{1.039} & \underline{0.689} & \underline{6.78} & 0.620          & 0.667             \\
    \textbf{Ours ($S=3$)}         & \textbf{7.09}$\times$     & \textbf{406.24}     & \underline{1.015} & \textbf{0.694} & \textbf{5.99} & \underline{0.627} & 0.637             \\
    \bottomrule
  \end{tabular*}
  \vspace{-2pt} 
  \caption{\textbf{Training-free quantitative comparisons on FLUX.1-dev~\citep{black2024flux}.} Image resolution is $1024^2$ and baseline rows are copied from Table 1 in RALU~\citep{ralu}. We group results based on $1\times, 5\times, 7\times$ speedup settings where progressive resolution approaches are applied in combination with reducing denoising steps. Our method improves the speed--quality tradeoff across all speedup settings compared to baseline approaches. Further evaluation details and comparisons are shown in Appendix Sec.~\ref{subsec:image_latent}.}
  % replicating RALU's~\citep{ralu} Table~1 protocol; we keep only spatial-acceleration baselines.
  \label{tab:flux-main}
\end{table*}

% \renewcommand{\arraystretch}{1.0}
\setlength{\doublerulesep}{0.5pt}
\setlength{\aboverulesep}{0.2ex}
\begin{table*}[!t]
  \centering
  \scriptsize
  \begin{tabular*}{\textwidth}{@{\extracolsep{\fill}} l ccccccc @{}}
    \toprule
    Method & \makecell{Speedup (s) $\uparrow$} & \makecell{TFLOPs $\downarrow$} & ImageReward $\uparrow$ & CLIP-IQA $\uparrow$ & NIQE $\downarrow$ & T2I-Comp. $\uparrow$ & GenEval $\uparrow$ \\
    \midrule
    Z-Image (50 steps)           & 1.00$\times$             & 4941.23             & 0.965    & 0.700    & 5.41    & 0.731    & 0.745 \\
    % Z-Image (32 steps)           & {1.56}$\times$             & {3166.62}             & 0.957    & 0.697    & {5.44}    & 0.686    & 0.725 \\
    \midrule
    \textbf{Ours} (TF, $S=2$)  & {1.65}$\times$ & {3132.03} & {0.904}             & {0.688}             & {5.87}             & {0.658} & {0.730} \\
    % \textbf{Ours ($S=3$)}  & \textbf{1.74}$\times$    & \textbf{1.72}$\times$    & 0.875             & 0.690             & {5.59} & 0.650             & 0.682 \\
    \textbf{Ours} (LoRA, $S = 2$)   & {1.65}$\times$             & {3132.03}             & \textbf{0.982} & \textbf{0.699} & \textbf{5.72}             & \textbf{0.725}                 & \textbf{0.731} \\
    \midrule
    \textbf{Ours} (TF, $S=3$)  & {1.74}$\times$ & {2871.09} & {0.875}             & {0.690}             & \textbf{5.59}             & {0.650} & {0.682} \\
    % \textbf{Ours ($S=3$)}  & \textbf{1.74}$\times$    & \textbf{1.72}$\times$    & 0.875             & 0.690             & {5.59} & 0.650             & 0.682 \\
    \textbf{Ours} (LoRA, $S = 3$)   & {1.74}$\times$             & {2871.09}             & \textbf{0.954} & \textbf{0.697} & {5.75}             & \textbf{0.717}                 & \textbf{0.728} \\
    % \textbf{Ours} (LoRA, $\delta=0.05$)   & \textbf{2.01}$\times$             & \textbf{2436.19}             & {0.919} & 0.684 & 6.12             & {0.661}                 & 0.718 \\
    % \midrule
    % \textbf{Ours} (TF, $S=3$)  & \textbf{1.74}$\times$ & \textbf{2871.09} & 0.875            & 0.690            & 5.59             & 0.650 & 0.682 \\
    % \textbf{Ours} (LoRA, $S=3$)  & \textbf{1.74}$\times$ & \textbf{2871.09} & \textbf{0.954}             &  \textbf{0.697}            & 5.75             & - & - \\

    % \textbf{Ours}
    \midrule
    \midrule
    % \textbf{Ours} (TF, $S=3$)  & \textbf{1.74}$\times$ & \textbf{2871.09} & 0.904             & 0.688             & 5.87             & {0.658} & {0.730} \\
    % % \textbf{Ours ($S=3$)}  & \textbf{1.74}$\times$    & \textbf{1.72}$\times$    & 0.875             & 0.690             & {5.59} & 0.650             & 0.682 \\
    % \textbf{Ours} (LoRA, $S=3$)   & \textbf{1.74}$\times$             & \textbf{2871.09}             & \textbf{0.982} & {0.699} & \textbf{5.72}             & \hx{TODO}                 & \hx{TODO} \\
    PixelGen (25 steps)         & 1.00$\times$             & 65.36             & 0.921    & 0.734    & 5.95    & 0.574    & 0.794 \\
    % PixelGen (13 steps)         & \textbf{1.92}$\times$             & {34.16}             & 0.886    & 0.726    & {5.89}    & 0.571    & 0.781 \\
    \midrule
    \textbf{Ours} (TF, $S=2$) & {1.60}$\times$             & {33.72}             & {0.799}    & {0.718}    & {6.10}    &  {0.568}  & \textbf{0.782} \\
    \textbf{Ours} (LoRA, $S=2$) & {1.55}$\times$             & {33.72}             & \textbf{0.913}    & \textbf{0.728}    & \textbf{5.87}    & \textbf{0.580}    & {0.776} \\
    \bottomrule
  \end{tabular*}
  % \vspace{3pt}
  \caption{\textbf{Fine-tuning quantitative comparisons on latent- (Z-Image~\citep{zimage}) and pixel-space image generation (PixelGen~\citep{pixelgen}).} Image resolution is $1024^2$ for Z-Image and $512^2$ for PixelGen. Across both latent- and pixel-space image generation, our fine-tuning method bridges the gap between model pretraining and progressive resolution inference, further improving generation quality while preserving efficiency gains. Further evaluation details and comparisons are shown in Appendix Sec.~\ref{subsec:image_pixel}.}
  \label{tab:zimage-main}
\end{table*}

% \renewcommand{\arraystretch}{1.0}
\setlength{\aboverulesep}{0.2ex}
\begin{table*}[!t]
  \centering
  \scriptsize
  \begin{tabular*}{\textwidth}{@{\extracolsep{\fill}} l cccccccc @{}}
    \toprule
    \makecell{Method} &
    \makecell{Speedup (s) $\uparrow$} &
    \makecell{TFLOPs $\downarrow$} &
    \makecell{Subject\\Consistency $\uparrow$} &
    \makecell{Background\\Consistency $\uparrow$} &
    \makecell{Motion\\Smoothness $\uparrow$} &
    \makecell{Dynamic\\Degree $\uparrow$} &
    \makecell{Aesthetic\\Quality $\uparrow$} &
    \makecell{Image\\Quality $\uparrow$} \\
    \midrule
    WAN~2.1 (50)           & 1.00$\times$ & 119292         & 0.9492          & 0.9621          & 0.9874          & 0.4800          & 0.5993          & 0.6133          \\
    \midrule[0.3pt]
    WAN~2.1 (25)           & 2.00$\times$ & 59646          & 0.9434          & 0.9597          & \textbf{0.9879} & 0.4250          & 0.5893          & 0.5706          \\
    \textbf{Ours ($S=2$)} & \textbf{2.03}$\times$ & \textbf{57417} & \textbf{0.9462} & \textbf{0.9598} & 0.9859          & \textbf{0.4950} & \textbf{0.5975} & \textbf{0.6114} \\
    \midrule[0.3pt]
    \textbf{Ours ($S=3$)} & 2.54$\times$ & 45953 & 0.9299          & 0.9499          & 0.9815          & 0.5700 & 0.5696          & 0.5990 \\
    \bottomrule
  \end{tabular*}
  \caption{\textbf{Quantitative comparison on latent-space video generation on WAN~2.1~\citep{wan}.} Video resolution is maintained at 720P. Our training-free approach demonstrates more than $2\times$ speedup while maintaining high generation quality and outperforming the 25-step full-resolution baseline.}
  \label{tab:wan_vbench_720p}
\end{table*}




% \renewcommand{\arraystretch}{1.0}
% \setlength{\aboverulesep}{0.2ex}
% \begin{table*}[!t]
%   \centering
%   \scriptsize
%   \begin{tabular*}{\textwidth}{@{\extracolsep{\fill}} l ccccccc @{}}
%     \toprule
%     \makecell{Method} &
%     \makecell{Speedup $\uparrow$} &
%     \makecell{Subject\\Consistency $\uparrow$} &
%     \makecell{Background\\Consistency $\uparrow$} &
%     \makecell{Motion\\Smoothness $\uparrow$} &
%     \makecell{Dynamic\\Degree $\uparrow$} &
%     \makecell{Aesthetic\\Quality $\uparrow$} &
%     \makecell{Image\\Quality $\uparrow$} \\
%     \midrule
%     WAN2.1 (50 steps)           & 1.00$\times$ & 0.9492          & 0.9621          & 0.9874 & 0.4800          & 0.5993          & 0.6133          \\
%     \midrule[0.3pt]
%     WAN2.1 (25 steps)           & 2.00$\times$ & 0.9434          & 0.9597          & \textbf{0.9879} & 0.4250          & 0.5893          & 0.5706          \\
%     \textbf{Ours ($S=2$)} & 2.03$\times$ & \textbf{0.9462} & \textbf{0.9598} & 0.9859 & \textbf{0.4950} & \textbf{0.5975} & \textbf{0.6114} \\
%     \midrule[0.3pt]
%     \textbf{Ours ($S=3$)} & 2.54$\times$ & 0.9299          & 0.9499          & 0.9815 & 0.5700 & 0.5696          & 0.5990 \\
%     \bottomrule
%   \end{tabular*}
%   \vspace{3pt}
%   \caption{\textbf{Quantitative comparison on latent-space video generation on WAN2.1~\citep{wan}.} Video resolution is maintained at 720P. Our training-free approach demonstrates more than $2\times$ speedup while maintaining high generation quality and outperforming the 25-step full-resolution baseline.}
%   \label{tab:wan_vbench_720p}
% \end{table*}


% We evaluate both training-free and LoRA fine-tuned variants of our Spectral Progressive Diffusion framework and report quantitative results on latent space image generation in Tables ~\ref{tab:flux-main} and \ref{tab:zimage-main} and qualitative comparisons in Fig.~\ref{fig:qual}.

\paragraph{Training-free acceleration.} In Table~\ref{tab:flux-main} and Fig.~\ref{fig:qual}, we compare our training-free acceleration approach (Sec.~\ref{sec:method:noise-injection}) on FLUX.1-dev with training-free acceleration baselines that implement progressive resolution growing, including Bottleneck Sampling~\citep{bottleneck} and RALU~\citep{ralu}. Additional FLUX.1-dev comparisons are included in Appendix Sec.~\ref{subsec:image_latent}. Our approach consistently improves the speed-quality tradeoff over other methods across all three evaluated speedup regimes. Increasing $S=2$ to $S=3$ further improves efficiency, offering up to $7.36 \times$ FLOPs speedup and $7.09 \times$ wall clock speedup while maintaining high generation quality.

% Thanks to our Spectral Progressive Diffusion framework and derived $\delta$-optimal resolution schedule (Secs.~\ref{sec:method:noise-injection}, \ref{sec:method:schedule:signal}), 

% i removed: $\delta$ being the sole hyperparameter for our resolution schedule allows us to easily increase. i think claiming "sole hyperparam" should be toned down

% "as it shares the same VAE backbone and hence the same resolution schedule as FLUX.1.dev -- no need for this IMO". the only reason is that we dont need to remeasure power spectrum, but realisticalluy we can do it on any model so this is not a strong argument (rather a neutral one, so i removed to save space)

\vspace{-5pt}
\paragraph{Spectral-transformation-based fine-tuning.} In Table~\ref{tab:zimage-main}, we demonstrate that our spectral-transformation-based fine-tuning method (Sec.~\ref{sec:method:lora}) further improves image quality and efficiency compared to training-free acceleration. Since FLUX.1-dev is guidance-distilled and its base model is not open-sourced, we select Z-Image \citep{zimage} as the base model for our latent-space image fine-tuning experiments. All models are fine-tuned with LoRA~\citep{hu2021lora} using rank 32 for 2000 iterations in both latent-space and pixel-space image generation experiments. Across both latent- and pixel-space image generation, our fine-tuning framework narrows the gap between training and inference, further improving generation quality while maintaining efficiency gains. % In particular, for latent-space image generation, our fine-tuned model with $S = 3$ resolution stages improves both the efficiency and quality compared to training-free inference with $S = 2$ where more denoising steps are in higher resolutions.
We include additional implementation details and results in Appendix Sec.~\ref{subsec:image_pixel}.

% Both training-free variants match or outperform all training-free baselines at equivalent or superior speedups (Table~\ref{tab:flux-main}). Increasing resolution stages from $S=2$ to $S=3$ further boosts efficiency with negligible image quality loss. Furthermore, the fine-tuned variant enables operation at a higher error threshold ($\delta=0.05$ vs. the training-free $\delta=0.01$), yielding simultaneous improvements in acceleration and image quality (Table \ref{tab:zimage-main}).

% \paragraph{Training-free inference.}
% The training-free setting tests whether our spectrum-derived schedule can be used as a plug-and-play modification to a pretrained generator, without changing its weights or sampler. On FLUX.1-dev, our method consistently improves the speed--quality tradeoff over spatial acceleration baselines across all evaluated budget regimes. At moderate acceleration, Ours ($S=2$) reaches $1.66\times$ wall-clock speedup and $1.70\times$ FLOP speedup while improving ImageReward and CLIP-IQA over RALU; at stronger acceleration, Ours ($S=2$) achieves $5.77\times$ wall-clock speedup with higher ImageReward and comparable text alignment to RALU. Increasing to $S=3$ further reduces compute, reaching up to $7.09\times$ wall-clock and $7.36\times$ FLOP speedup, with only gradual degradation in alignment metrics. These results support the central premise of our method: early high-resolution tokens are largely wasted on noise-dominated frequencies, so removing them until their predicted activation times yields substantial acceleration without requiring backbone-specific retraining.

% \paragraph{LoRA fine-tuning.}
% We next evaluate whether a lightweight adapter can absorb the remaining mismatch between the idealized spectral trajectory and the distribution seen by the pretrained backbone. In contrast to full retraining or architectural changes, our LoRA variant keeps the base model frozen and only trains low-rank adapters on the stage-specific velocity targets derived in Sec.~\ref{sec:method:lora}. This fine-tuning improves the quality--efficiency balance on Z-Image: Ours (LoRA) preserves essentially the same acceleration as the training-free $S=2$ configuration while improving ImageReward and NIQE, closing part of the gap to the full-resolution baseline. We view this as complementary to the training-free result: the schedule alone already gives immediate speedups, while LoRA provides a small-parameter path for models or operating points where the progressive spectral trajectory benefits from adaptation.

%The passthrough row confirms that replacing noise-dominated bands with pure noise at every step is essentially lossless, isolating the speedup purely to the resolution schedule.

%\paragraph{LoRA fine-tuning.}
% Augmenting the backbone with a LoRA~\citep{hu2021lora} adapter trained on our per-frequency flow trajectory (Section~\ref{sec:method:lora}) recovers---and in many cells exceeds---the base backbone's quality at the same or stronger speedup. The adapter is trained only on noise-dominated-band rewrites at the post-transition resolution; the pretrained backbone is fully frozen. Compared to full retraining or architectural surgery, the LoRA approach adds at most a few percent trainable parameters to the original backbone and trains in a single-digit number of GPU-hours on a single H100; exact parameter counts and compute are reported in the appendix. The Ours-LoRA row in Table~\ref{tab:flux-main} and the rightmost column of Fig.~\ref{fig:qual} make the effect directly comparable against the training-free variants.

\vspace{-5pt}
\subsection{Latent-space Video Generation}
\label{sec:exp:video}

We further demonstrate the flexibility and robustness of our training-free acceleration approach on latent-space video generation (Table \ref{tab:wan_vbench_720p}, Fig. \ref{fig:qual}). To our knowledge, there are no diffusion acceleration approaches that were validated on all three generation modalities presented in Secs.~\ref{sec:exp:image}--\ref{sec:exp:video}. We observe trends consistent with image generation, where our method achieves significant speedups while preserving generation quality. We include additional implementation details and results in Appendix Sec.~\ref{sec:video_generation}.


\vspace{-5pt}
\subsection{Ablation Studies}
\label{sec:exp:ablations}

We ablate our core design choices, including error threshold $\delta$, resolution stages $S$, and spectral transformation $T_\Phi$, on latent-space image generation using FLUX.1-dev. We show the speed-quality Pareto frontier in Fig. \ref{fig:ablation}(a) for $\delta \in [0.001, 0.1]$, $S \in \{2, 3, 4\}$, and different spectral transformations $T_\Phi$. Consistent with our theory, increasing $\delta$ or $S$ further accelerates generation at the cost of gradual image quality loss. The Discrete Cosine Transform (DCT), Discrete Wavelet Transform (DWT), and Fourier Transform (FFT) achieve similar high-fidelity performance. Qualitative comparisons are given in Appendix Sec.~\ref{sec:ablation}.

\begin{figure}[t]
  \centering
  \includegraphics[width=\linewidth]{figures/ablation_freq_edit.pdf}
  \caption{\textbf{(a): Ablation Studies on $\delta, S$ and $T_\Phi$.} We observe a clear tradeoff between image quality and efficiency when varying $\delta$ and $S$ as shown in the top plot. As shown in the bottom plot, all three transforms (DCT, DWT, FFT) achieve comparable quality and follow the same quality--efficiency tradeoff as $\delta$ increases. \textbf{(b): Frequency-based Image Editing.} Our method demonstrates superior prompt alignment and geometric consistency compared to standard SDEdit-style spatial-domain editing under identical step counts. Our method successfully edits the texture and style of the input image, whereas SDEdit fails to do so.}
  \label{fig:ablation}
\end{figure}


\iffalse
% move to supp
\textbf{Resolution Stages $S$:} The primary speedup derives from bypassing native-resolution steps; consequently, $S=2$ configurations occupy the top-right of the Pareto curve, while higher $S$ provides diminishing efficiency gains by further optimizing the already-inexpensive low-resolution trajectory. \textbf{Spectral Transformation $T_\Phi$:} The Fourier Transform (FFT) tends to oversmooth results, whereas DCT and DWT achieve similar high-fidelity performance. We default to DCT for its native support of non-power-of-two resolution scaling.
\fi



\iffalse
\begin{table}[t]
  \centering
  \footnotesize
  \caption{\textbf{Ablation: tolerance $\delta$} on FLUX.1-dev, using the same metric set as Table~\ref{tab:flux-main}. Two blocks (training-free, LoRA). All values placeholder.}
  \label{tab:delta-sweep}
  \setlength{\tabcolsep}{3pt}
  \begin{tabular}{c ccc ccccc}
    \toprule
    $\delta$ & Latency (s) $\downarrow$ & TFLOPs $\downarrow$ & Speed.\ $\uparrow$ & ImageReward $\uparrow$ & CLIP-IQA $\uparrow$ & NIQE $\downarrow$ & T2I-Comp.\ $\uparrow$ & GenEval $\uparrow$ \\
    \midrule
    \multicolumn{9}{l}{\textit{Training-free}} \\
    0.001 & - & - & - & - & - & - & - & - \\
    0.01  & - & - & - & - & - & - & - & - \\
    0.1   & - & - & - & - & - & - & - & - \\
    \midrule
    \multicolumn{9}{l}{\textit{LoRA}} \\
    0.001 & - & - & - & - & - & - & - & - \\
    0.01  & - & - & - & - & - & - & - & - \\
    0.1   & - & - & - & - & - & - & - & - \\
    \bottomrule
  \end{tabular}

  \vspace{1em}

  \small
  \caption{\textbf{Ablation: number of resolution stages} on FLUX.1-dev. All values placeholder.}
  \label{tab:stages-sweep}
  \setlength{\tabcolsep}{8pt}
  \begin{tabular}{c ccc}
    \toprule
    \# Stages & FID $\downarrow$ & FLOPs $\downarrow$ & Speedup $\uparrow$ \\
    \midrule
    2 & - & - & - \\
    3 & - & - & - \\
    4 & - & - & - \\
    \bottomrule
  \end{tabular}
\end{table}
\fi

\subsection{Frequency-based Image Editing}
\label{sec:exp:applications}
We demonstrate a unique \textit{frequency-based} image editing capability enabled by our method. Specifically, given an input $\mathbf{x}_\text{in}$ to edit, we add noise to the low-frequency part of $T_\Phi(\mathbf{x}_\text{in})$, perform spectral noise expansion and timestep alignment (Sec.~\ref{sec:method:noise-injection}), and then resume denoising from the timestep $t$ corresponding to the resolution schedule with the editing prompt. Fig.~\ref{fig:ablation}(b) shows our method achieves substantially improved prompt alignment and geometric consistency compared to a spatial-domain SDEdit-style baseline \cite{meng2022sdedit}. Additional details and editing results are in Appendix Sec.~\ref{sec:image_editing}.

% Please refer to the supplementary materials for implementation details and additional editing results.

\iffalse
, where we corrupt the input to the same diffusion timestep $t$ when we perform spectral noise expansion (approximately half the number of steps along the default schedule) and resume denoising. We find that SDEdit fails to modify texture and style, whereas our method succeeds. Please refer to the supplementary materials for more results.
\fi

\section{Discussion}\vspace{-5pt}
\label{sec:conclusion}
In this work, we introduce Spectral Progressive Diffusion, a generation framework that leverages spectral autoregression in diffusion models to match computation to frequencies where signal emerges from noise. 
Through a principled spectral noise expansion mechanism and an optimal resolution schedule derived from the model’s power spectrum, we enable progressive resolution growth along the denoising trajectory of pre-trained image and video models, achieving notable speedups. While our approach applies directly in a training-free setting, it implicitly assumes that the underlying model can generalize across varying resolutions. Though modern DiT models are trained on multiple resolutions, the continuous schedules we employ may include intermediate resolutions not seen during training. In practice, using a small number of transitions near the training distribution suffices to achieve substantial efficiency gains in both image and video generation. Our fine-tuning strategy aligns the model with the progressive-resolution trajectory, reducing the training--inference gap and improving quality, suggesting the use of our method during pretraining as a promising direction. While we focus on spatial spectra, extending the framework to temporal frequencies is a natural direction for accelerating video generation. Overall, our approach provides a principled and complementary approach for improving generation efficiency via the spectral autoregressive structure of diffusion models while maintaining high generation quality.




\iffalse
\paragraph{Limitations and future work.} While our approach applies directly in a training-free setting, it implicitly assumes that the underlying model can generalize across varying resolutions.
In practice, modern DiT architectures are not restricted to a fixed resolution, and most state-of-the-art models are trained on multiple resolutions.
Still, the discrete set of resolutions seen during training does not span the continuous range of resolutions our schedule can operate at, potentially leaving some intermediate resolutions outside the model's data distribution.
Empirically, we find that even with a few transitions, which stay close to the in-distribution regime, our approach delivers significant efficiency gains on both image and video generation.
Our fine-tuning strategy aligns the model with the progressive-resolution trajectory, further closing the training-inference gap and improving quality. This suggests that adapting our framework during pretraining is a promising direction.
While we focus on the spatial spectra, extending our framework to temporal frequencies in video models is a natural avenue for accelerating video generation. Overall, our approach provides a principled, complementary path to improving generative efficiency via the spectral autoregressive structure of diffusion models.
\fi



% We introduce Spectral Progressive Diffusion, a progressive generation framework that eliminates computation on noise-dominated high-frequency components in early denoising, yielding substantial efficiency gains without degrading image quality. Its effectiveness depends on the multi-resolution capabilities of the underlying pretrained model, suggesting adaptation during pretraining as a promising direction. While we focus on the spatial spectra, extending our framework to temporal frequencies is a natural avenue for accelerating video generation. Overall, our approach provides a principled, complementary path to improving generative efficiency via the spectral autoregressive structure of diffusion models.

% freq-domain CFG, 
% limitation: pretrained models that do not support multi-resolution - leading to pretraining strategies. video tempory redudancy, power-law assumption

% future work: pretraining, distillation, CFG,

\section*{Acknowledgments}
Howard Xiao and Brian Chao are supported by Stanford Graduate Fellowships (SGF). Brian Chao is also supported by the NSF Graduate Research Fellowship Program (GRFP). We thank Google and Toyota Research Institute (TRI). We thank Shengqu Cai, Hansheng Chen, Kiyohiro Nakayama, and Zichun Xu for fruitful discussions. Compute resources were provided by the Marlowe cluster at Stanford University~\cite{marlowe2025}.


%%%%%%%%%%%%%%%%%%%%%%%%%%%%%%%%%%%%%%%%%%%%%%%%%%%%%%%%%%%%
% References
%%%%%%%%%%%%%%%%%%%%%%%%%%%%%%%%%%%%%%%%%%%%%%%%%%%%%%%%%%%%
\clearpage   % flush all pending figure floats before the bibliography
{
\small
\bibliographystyle{plainnat}
\bibliography{references_arxiv}
}

%%%%%%%%%%%%%%%%%%%%%%%%%%%%%%%%%%%%%%%%%%%%%%%%%%%%%%%%%%%%
% !TEX root = ../arxiv.tex
\makeatletter
\let\addcontentsline\arxiv@savedaddcontentsline
\makeatother

\newpage
\setcounter{tocdepth}{2}
\renewcommand{\contentsname}{Appendix Contents}
\tableofcontents
\newpage
\appendix

\section{Supporting Definitions and Proofs for Section~\ref{sec:method:schedule:signal}}
\label{app:proofs-schedule}

This appendix provides the justification of timestep alignment in Spectral Progressive Diffusion, the deferred proofs of Propositions~1 and~2 from Sec.~\ref{sec:method:schedule:signal}, along with supporting definitions.

\subsection{Definitions}
\label{app:prelim-schedule}

\paragraph{Per-frequency SNR.}
Following~\citet{kingma2021vdm}, we define the \emph{per-frequency signal-to-noise ratio (SNR)} as the ratio of the clean- and noise-component second moments of the per-frequency forward process in Eq.~\eqref{eq:forward-mode}:
\begin{equation}
  \mathrm{SNR}_\omega(t)
  \;:=\;
  \frac{\mathbb{E}[\,|(1-t)\,x_0^{(\omega)}|^2\,]}{\mathbb{E}[\,|t\,\epsilon^{(\omega)}|^2\,]}
  \;=\;
  \underbrace{\frac{(1-t)^2}{t^2}}_{\mathrm{SNR}(t)}\;\cdot\;P_\omega,
  \label{eq:snr}
\end{equation}
which factorizes into a frequency-independent timestep term $\mathrm{SNR}(t) = (1-t)^2/t^2$ and the frequency-specific signal power $P_\omega$.

\paragraph{Resolution grids and representable frequencies.}
By the Nyquist--Shannon sampling theorem~\citep{nyquist1928certain,shannon1949communication}, a spatial grid of size $H \times W$ can represent 2D frequencies in the rectangle $|\omega_x| \le W/2$, $|\omega_y| \le H/2$. We define the maximum (fully) representable radial frequency as
\begin{equation}
  \omega_{\max}(H, W) \;=\; \frac{\min(H, W)}{2}.
  \label{eq:nyquist-cap}
\end{equation}
Note that our definition of the maximum (fully) representable frequency above is conservative, as it captures only the maximum radial frequency fully supported on the spatial grid $H \times W$.
A stage operating at reduced resolution $(sH, sW)$ for $s \in (0,1]$ can represent only a subset of the representable frequencies at full-resolution:
\begin{equation}
  \Omega_s \;:=\; \{\,\omega \in \Omega : |\omega| \le s \cdot \omega_{\max}(H, W)\,\} \;\subseteq\; \Omega.
  \label{eq:nyquist-set}
\end{equation}

\subsection{Justification of Timestep Alignment}
\label{app:time-align}

We verify that applying the scalar rescaling in Eq.~\eqref{eq:scale-align} and querying the pretrained model at the aligned timestep $\tilde{t}_i$ in Eq.~\eqref{eq:time-align} maps the spectral noise expansion output $\mathbf{x}^{s_{i+1}}_{t_i}$ to a valid flow-matching state $\tilde{\mathbf{x}}^{s_{i+1}}_{\tilde{t}_i}$ at the new resolution scale $s_{i+1} > s_i$. Throughout this subsection we write $r := s_{i+1}/s_i > 1$ for the resolution ratio; the final formulas are stated in both the $r$ and $s_{i+1}/s_i$ forms to tie directly back to Eqs.~\eqref{eq:scale-align}--\eqref{eq:time-align} in Sec.~\ref{sec:method:noise-injection}. The derivations below assume the ideal setting in which the network outputs match the ground-truth flow targets.

\paragraph{1. Setup.}
The spectral noise expansion produces an expanded state $\mathbf{x}^{s_{i+1}}_{t_i}$ in the frequency domain of the expanded resolution grid. For a grid with $N_{s_{i+1}}$ points, the coefficients are:
$$
\left(x^{s_{i+1}}_{t_i}\right)^{(\omega)} =
\begin{cases}
(1-t_i)\left(x_{0}^{s_i}\right)^{(\omega)} + t_i\epsilon^{(\omega)}, & \omega \in \Omega_{s_i}, \\[2pt]
t_i\epsilon'^{(\omega)}, & \omega \in \Omega_{s_{i+1}} \setminus \Omega_{s_i},
\end{cases}
$$
where $\epsilon^{(\omega)}, \epsilon'^{(\omega)} \sim \mathcal{N}(0, 1)$ are independent standard Gaussian variables. Under the orthonormal identification of spectra across resolutions, the clean-signal coefficient at a shared frequency $\omega \in \Omega_{s_i}$ relates to its scale-$s_{i+1}$ counterpart by the resolution ratio $r$:
$$
\left(x_{0}^{s_{i+1}}\right)^{(\omega)} = r \cdot \left(x_{0}^{s_i}\right)^{(\omega)} = (s_{i+1}/s_i) \cdot \left(x_{0}^{s_i}\right)^{(\omega)}, \quad \omega \in \Omega_{s_i}.
$$
For the newly exposed high-frequency slots $\omega \in \Omega_{s_{i+1}} \setminus \Omega_{s_i}$, the clean coefficients are treated as noise-dominated at timestep $t_i$ per Proposition~1, so spectral noise expansion initializes them with the correct noise level.

\paragraph{2. Timestep alignment.}
We seek a scaling factor $\kappa_i$ and an aligned timestep $\tilde{t}_i \in (0, 1)$ such that, for every $\omega \in \Omega_{s_{i+1}}$, the transformed state satisfies:
$$
\left(\tilde{x}^{s_{i+1}}_{\tilde{t}_i}\right)^{(\omega)} = \kappa_i\left(x^{s_{i+1}}_{t_i}\right)^{(\omega)} = (1-\tilde{t}_i)\left(x_{0}^{s_{i+1}}\right)^{(\omega)} + \tilde{t}_i\tilde\epsilon^{(\omega)}, \quad \tilde\epsilon^{(\omega)} \sim \mathcal{N}(0,1) \text{ i.i.d.}
$$
By matching the noise coefficients across all bands and the signal coefficients on the shared band $\Omega_{s_i}$, we derive two scalar conditions:
$$
\kappa_i \cdot t_i = \tilde{t}_i, \quad \kappa_i(1-t_i) = r(1-\tilde{t}_i).
$$
Substituting $\kappa_i = \tilde{t}_i/t_i$ from the noise match into the signal match equation gives:
$$
\frac{\tilde{t}_i}{t_i}(1-t_i) = r(1-\tilde{t}_i)
$$
Solving for $\tilde{t}_i$ and re-expressing the result in $s_{i+1}/s_i$ notation gives:
$$
\tilde{t}_i = \frac{r \cdot t_i}{1 + (r - 1)t_i} = \frac{(s_{i+1}/s_i) \cdot t_i}{1 + ((s_{i+1}/s_i) - 1)t_i}.
$$
Consequently, the scaling factor $\kappa_i$ is:
$$
\kappa_i = \frac{r}{1 + (r - 1)t_i} = \frac{s_{i+1}/s_i}{1 + ((s_{i+1}/s_i) - 1)t_i}.
$$
Substituting these back into the setup confirms that $\left(\tilde{x}^{s_{i+1}}_{\tilde{t}_i}\right)^{(\omega)} = (1-\tilde{t}_i)\left(x_{0}^{s_{i+1}}\right)^{(\omega)} + \tilde{t}_i\tilde\epsilon^{(\omega)}$ holds for all $\omega \in \Omega_{s_{i+1}}$. This proves that the resulting scaled and timestep-aligned state $\tilde{\mathbf{x}}^{s_{i+1}}_{\tilde{t}_i}$ is a valid flow-matching state at resolution $s_{i+1}$ and timestep $\tilde{t}_i$. $\square$

\subsection{Proof of Proposition 1: Per-frequency $\delta$-optimal Activation Time}
\label{app:proof-prop1}

We first establish a lemma giving the Bayes-optimal velocity error in closed form under a simplified modelling assumption $x_0^{(\omega)} \sim \mathcal{N}(0, P_\omega)$. In practice, we measure the power spectrum on centered (i.e. zero-mean) inputs $\mathbf{x}_0$ to get the power spectrum $P_\omega$, and this modelling assumption serves as a second moment approximation that allows the analytical derivation of our $\delta$-optimal resolution schedule. The proof of Proposition~1 then follows by solving the $\delta$-bound on the Bayes-optimal velocity error for $t$ to obtain the activation time $t_\omega$ in Eq.~\eqref{eq:activation}.

\begin{lemma}[Bayes-optimal velocity error] Let $v^{*(\omega)}$ be the per-frequency component of the Bayes-optimal velocity predictor $\mathbf{v}^{*}$ defined in Sec.~\ref{sec:method:prelim:basics}:

\begin{equation}
  \mathbf{v}^*(\mathbf{x}_t, t) \;=\; \mathbb{E}\!\left[\,\boldsymbol{\epsilon} - \mathbf{x}_0 \,\big|\, \mathbf{x}_t\,\right].
  \label{eq:optimal-v}
\end{equation}
Under the modelling assumption $x_0^{(\omega)} \sim \mathcal{N}(0, P_\omega)$, $v^{*(\omega)}$ deviates from the noise-only prediction $\epsilon^{(\omega)}$ in expectation by
\begin{equation}
  \mathbb{E}\!\left[\bigl|v^{*(\omega)}(\mathbf{x}_t, t) - \epsilon^{(\omega)}\bigr|^2\right]
  \;=\; \frac{\mathrm{SNR}_\omega(t)\,(1 + P_\omega)}{1 + \mathrm{SNR}_\omega(t)}.
  \label{eq:velocity-error-optimal}
\end{equation}
\end{lemma}


%\paragraph{Remark.} Real-data latents are not perfectly Gaussian, so Eq.~\eqref{eq:velocity-error-optimal} plays the role of an analytically tractable reference curve for our resolution schedule calculations. In practice, classifier-free guidance (CFG) from the text prompt during the generation process can also allow lower frequencies to be generated faster. We empirically measure this error in Sec.~\ref{sec:motivation}. \hx{TODO!}

\begin{proof}
\leavevmode

\paragraph{1. Assumptions.} The proof assumes a linear-Gaussian model, where:
\begin{enumerate}
  \item[(A1)] we model the clean data spectrum as zero-mean Gaussian $x_0^{(\omega)} \sim \mathcal{N}(0,\, P_\omega)$ with second moment $P_\omega > 0$;
  \item[(A2)] we assume $\epsilon^{(\omega)} \sim \mathcal{N}(0, 1)$, by orthonormality of $\Phi$ applied to $\boldsymbol{\epsilon} \sim \mathcal{N}(0, I)$;
  \item[(A3)] we assume $x_0^{(\omega)} \perp \epsilon^{(\omega)}$, i.e. the signal and noise coefficients are independent.
\end{enumerate}
(A1) is the simplified modelling assumption; (A2) and (A3) follow directly from the flow-matching forward process. Under (A1), the triple $(x_0^{(\omega)},\, \epsilon^{(\omega)},\, x_t^{(\omega)})$ is jointly Gaussian, so conditional expectations are linear and closed-form.

\paragraph{2. Second-moment quantities.} From the per-frequency forward process in Eq.~\eqref{eq:forward-mode} and (A1)--(A3),
\begin{align}
  \mathbb{E}[x_t^{(\omega)}]                     &= 0, \qquad
  \mathrm{Var}(x_t^{(\omega)}) \;=\; (1-t)^2 P_\omega + t^2,  \\
  \mathrm{Cov}(x_t^{(\omega)},\, x_0^{(\omega)})      &= (1-t)\,P_\omega, \qquad
  \mathrm{Cov}(x_t^{(\omega)},\, \epsilon^{(\omega)}) \;=\; t.
\end{align}

\paragraph{3. Bayes-optimal velocity as a linear predictor.}
Under (A1)--(A3), the triplet $(x_0^{(\omega)},\, \epsilon^{(\omega)},\, x_t^{(\omega)})$ is jointly Gaussian with zero mean, since $x_t^{(\omega)} = (1-t) x_0^{(\omega)} + t\,\epsilon^{(\omega)}$ is a linear combination of two independent zero-mean Gaussians. For any pair of zero-mean jointly Gaussian random variables $(X, Y)$ with $\mathrm{Var}(Y) > 0$, the conditional expectation is a linear function of $Y$:
\begin{equation}
  \mathbb{E}[X \,|\, Y] \;=\; \frac{\mathrm{Cov}(X, Y)}{\mathrm{Var}(Y)}\,Y.
  \label{eq:gauss-cond}
\end{equation}

We apply \eqref{eq:gauss-cond} twice, using the second moments computed above. First, with $X = x_0^{(\omega)}$ and $Y = x_t^{(\omega)}$,
\begin{equation}
  \mathbb{E}\bigl[x_0^{(\omega)} \,\big|\, x_t^{(\omega)}\bigr]
  \;=\; \frac{(1-t)\,P_\omega}{(1-t)^2 P_\omega + t^2}\, x_t^{(\omega)}.
  \label{eq:cond-x0}
\end{equation}
Second, with $X = \epsilon^{(\omega)}$ and $Y = x_t^{(\omega)}$,
\begin{equation}
  \mathbb{E}\bigl[\epsilon^{(\omega)} \,\big|\, x_t^{(\omega)}\bigr]
  \;=\; \frac{t}{(1-t)^2 P_\omega + t^2}\, x_t^{(\omega)}.
  \label{eq:cond-eps}
\end{equation}

The Bayes-optimal velocity at frequency $\omega$ is $v^{*(\omega)} = \mathbb{E}\bigl[\epsilon^{(\omega)} - x_0^{(\omega)} \,\big|\, \mathbf{x}_t\bigr]$, the $\phi_\omega$-projection of the full-input optimal velocity in Eq.~\eqref{eq:optimal-v}. Under the modelling assumption that $\mathbf{x}_0$ has diagonal covariance in basis $\Phi$ (i.e., (A1) holds independently for every $\omega$), the joint Gaussian $(\mathbf{x}_0, \boldsymbol{\epsilon}, \mathbf{x}_t)$ factorizes across $\Omega$, so the posterior collapses to its single-coefficient marginal:
\begin{equation}
  \mathbb{E}\bigl[\epsilon^{(\omega)} - x_0^{(\omega)} \,\big|\, \mathbf{x}_t\bigr]
  \;=\; \mathbb{E}\bigl[\epsilon^{(\omega)} - x_0^{(\omega)} \,\big|\, x_t^{(\omega)}\bigr].
  \label{eq:full-to-scalar}
\end{equation}
Subtracting \eqref{eq:cond-x0} from \eqref{eq:cond-eps} and combining with \eqref{eq:full-to-scalar},
\begin{equation}
  v^{*(\omega)} \;=\; A \cdot x_t^{(\omega)},
  \qquad
  A \;:=\; \frac{t - (1-t)\,P_\omega}{(1-t)^2 P_\omega + t^2}.
  \label{eq:vstar-linear}
\end{equation}

\paragraph{4. Substitute and simplify.} Plugging $x_t^{(\omega)} = (1-t)\, x_0^{(\omega)} + t\,\epsilon^{(\omega)}$ into \eqref{eq:vstar-linear} and collecting terms in $x_0^{(\omega)}$ and $\epsilon^{(\omega)}$,
\begin{equation}
  v^{*(\omega)} - \epsilon^{(\omega)}
  \;=\; A\bigl[(1-t)\,x_0^{(\omega)} + t\,\epsilon^{(\omega)}\bigr] - \epsilon^{(\omega)}
  \;=\; A(1-t)\,x_0^{(\omega)} \,+\, (At - 1)\,\epsilon^{(\omega)}.
\end{equation}
Squaring and taking expectation, using (A1)--(A3) (so $\mathbb{E}[(x_0^{(\omega)})^2] = P_\omega$, $\mathbb{E}[(\epsilon^{(\omega)})^2] = 1$, and $\mathbb{E}[x_0^{(\omega)}\,\epsilon^{(\omega)}] = 0$), the cross term vanishes:
\begin{equation}
  \mathbb{E}\!\left[\bigl|v^{*(\omega)} - \epsilon^{(\omega)}\bigr|^2\right]
  \;=\; A^2 (1-t)^2\,P_\omega \,+\, (At - 1)^2.
\end{equation}
Expanding $(At-1)^2 = A^2 t^2 - 2At + 1$ and grouping the $A^2$ terms,
\begin{equation}
  \mathbb{E}\!\left[\bigl|v^{*(\omega)} - \epsilon^{(\omega)}\bigr|^2\right]
  \;=\; A^2\!\left[(1-t)^2 P_\omega + t^2\right] - 2At + 1.
  \label{eq:verror-expanded}
\end{equation}
Denote $D := (1-t)^2 P_\omega + t^2 = \mathrm{Var}(x_t^{(\omega)})$; from Eq.~\eqref{eq:vstar-linear},
\begin{equation}
  A \;=\; \frac{t - (1-t)\, P_\omega}{D}.
  \label{eq:A-in-D}
\end{equation}
Substituting into Eq.~\eqref{eq:verror-expanded} term by term,
\begin{align}
  A^2\!\left[(1-t)^2 P_\omega + t^2\right] \;&=\; A^2 D \;=\; \frac{(t - (1-t)\, P_\omega)^2}{D}, \\
  2At \;&=\; \frac{2t\,(t - (1-t)\, P_\omega)}{D},
\end{align}
so the right-hand side of Eq.~\eqref{eq:verror-expanded} combines into a single fraction over $D$:
\begin{equation}
  \mathbb{E}\!\left[\bigl|v^{*(\omega)} - \epsilon^{(\omega)}\bigr|^2\right]
  \;=\; \frac{(t - (1-t)\, P_\omega)^2 \,-\, 2t\,(t - (1-t)\, P_\omega) \,+\, D}{D}.
  \label{eq:verror-singlefrac}
\end{equation}
Expanding $(t - (1-t)\, P_\omega)^2 = t^2 - 2t(1-t) P_\omega + (1-t)^2 P_\omega^2$, the numerator of Eq.~\eqref{eq:verror-singlefrac} simplifies as
\begin{align}
  & \bigl[t^2 - 2t(1-t) P_\omega + (1-t)^2 P_\omega^2\bigr]
    - \bigl[2t^2 - 2t(1-t) P_\omega\bigr]
    + \bigl[(1-t)^2 P_\omega + t^2\bigr] \nonumber\\
  ={}& (1-t)^2 P_\omega^2 + (1-t)^2 P_\omega
    \;=\; (1-t)^2\,P_\omega\,(1 + P_\omega).
\end{align}
Substituting back into Eq.~\eqref{eq:verror-singlefrac},
\begin{equation}
  \mathbb{E}\!\left[\bigl|v^{*(\omega)} - \epsilon^{(\omega)}\bigr|^2\right]
  \;=\; \frac{(1-t)^2\,P_\omega\,(1 + P_\omega)}{(1-t)^2 P_\omega + t^2}.
\end{equation}

\paragraph{5. Rewrite via $\mathrm{SNR}_\omega(t)$.} Using $\mathrm{SNR}_\omega(t) = (1-t)^2 P_\omega / t^2$ from Eq.~\eqref{eq:snr}, the denominator is $(1-t)^2 P_\omega + t^2 = t^2(1 + \mathrm{SNR}_\omega(t))$ and the numerator equals $t^2\,\mathrm{SNR}_\omega(t)\,(1 + P_\omega)$. The $t^2$ cancels, recovering the right-hand side of Eq.~\eqref{eq:velocity-error-optimal} and completing the proof of the lemma.

\end{proof}

\begin{proposition}[Per-frequency $\delta$-optimal activation time] Denote the per-frequency Bayes-optimal velocity predictor by $v^{*(\omega)}(\mathbf{x}_t, t)$, whose spatial-domain counterpart is defined in Eq.~\eqref{eq:optimal-v}. Under the simplified modelling assumption $x_0^{(\omega)} \sim \mathcal{N}(0, P_\omega)$ with $P_\omega > 0$, we have
\begin{equation}
  \mathbb{E}\!\left[\bigl|v^{*(\omega)}(\mathbf{x}_t, t) - \epsilon^{(\omega)}\bigr|^2\right] \;\le\; \delta
  \label{eq:noise-dominated-appendix}
\end{equation}
for all $t \ge t_\omega$ and $\delta\in(0, 1)$, where
\begin{equation}
  t_\omega \;:=\; \frac{1}{1 + \sqrt{\dfrac{\delta}{P_\omega\,(1 + P_\omega - \delta)}}}.
  \label{eq:activation-appendix}
\end{equation}
\end{proposition}

\begin{proof}
\leavevmode

\paragraph{Solving the $\delta$-bound for $t$.} Substituting the closed-form error into the tolerance bound $\mathbb{E}[|v^{*(\omega)} - \epsilon^{(\omega)}|^2] \le \delta$ of Proposition~1 yields
\begin{equation}
  \frac{\mathrm{SNR}_\omega(t)\,(1 + P_\omega)}{1 + \mathrm{SNR}_\omega(t)} \;\le\; \delta.
\end{equation}
Clearing the (positive) denominator and collecting terms in $\mathrm{SNR}_\omega(t)$,
\begin{equation}
  \mathrm{SNR}_\omega(t)\,(1 + P_\omega - \delta) \;\le\; \delta,
  \qquad \text{i.e.,} \qquad
  \mathrm{SNR}_\omega(t) \;\le\; \frac{\delta}{1 + P_\omega - \delta},
\end{equation}
where the rearrangement is valid because $\delta < 1 \le 1 + P_\omega$ keeps $1 + P_\omega - \delta > 0$. Substituting $\mathrm{SNR}_\omega(t) = (1-t)^2 P_\omega / t^2$ and taking the positive square root (all quantities positive for $t \in (0, 1)$),
\begin{equation}
  \frac{1-t}{t} \;\le\; \sqrt{\frac{\delta}{P_\omega\,(1 + P_\omega - \delta)}}.
\end{equation}
Adding $1$ to both sides gives $1/t \le 1 + \sqrt{\delta / (P_\omega(1 + P_\omega - \delta))}$, which inverts to
\begin{equation}
  t \;\ge\; \frac{1}{1 + \sqrt{\dfrac{\delta}{P_\omega\,(1 + P_\omega - \delta)}}} \;=\; t_\omega,
\end{equation}
recovering Eq.~\eqref{eq:activation-appendix}.
\end{proof}


\subsection{Proof of Proposition 2: Per-resolution $\delta$-optimal Transition Time}
\label{app:proof-prop2}

\begin{proposition}[Per-resolution $\delta$-optimal transition time] Under the setting of Proposition~1, and assuming $P_\omega$ is monotonically decreasing in $|\omega|$ based on the power-law decay in Eq.~\eqref{eq:power-law}, for any resolution scale $s_i, s_{i+1} \in (0, 1]$, the optimal transition time from scale $s_i$ up to $s_{i+1}$ is
\begin{equation}
  t_i^* \;:=\; \min_{\omega \in \Omega_{s_i}}\, t_\omega \;=\; t_{\omega=s_i \cdot \omega_{\max}(H,W)},
  \label{eq:stage-transition-appendix}
\end{equation}
where $\Omega_{s_i}$ is the set of frequencies representable on the $(s_iH, s_iW)$ grid.
\end{proposition}

\begin{proof}
\leavevmode

We show that the activation time $t_\omega$ of Eq.~\eqref{eq:activation} is strictly decreasing in $|\omega|$ under the assumptions of the proposition; the claim then follows immediately.

\paragraph{1. $t_\omega$ is strictly increasing in $P_\omega$.}
For $\delta \in (0, 1)$ and $P_\omega > 0$, the quadratic $P_\omega\,(1 + P_\omega - \delta) = P_\omega^2 + (1-\delta)\, P_\omega$ has derivative $2 P_\omega + (1-\delta) > 0$, hence is strictly increasing in $P_\omega$. Therefore the radicand $\delta / \bigl(P_\omega(1 + P_\omega - \delta)\bigr)$ in Eq.~\eqref{eq:activation} is strictly decreasing in $P_\omega$, and so is its square root. Hence
\begin{equation}
  t_\omega \;=\; \frac{1}{1 + \sqrt{\dfrac{\delta}{P_\omega\,(1 + P_\omega - \delta)}}}
\end{equation}
is strictly increasing in $P_\omega$.

\paragraph{2. $P_\omega$ is strictly decreasing in $|\omega|$.}
The power-law assumption in Eq.~\eqref{eq:power-law} gives $P_\omega \propto |\omega|^{-\beta}$ with $\beta > 0$. Thus $P_\omega$ is strictly decreasing in $|\omega|$.

\paragraph{3. Minimum over $\Omega_s$.}
Combining the above observations, $t_\omega$ is strictly decreasing in $|\omega|$. The set $\Omega_s = \{\omega \in \Omega : |\omega| \le s \cdot \omega_{\max}(H, W)\}$ consists exactly of the frequencies with radius at most $s \cdot \omega_{\max}(H, W)$, so
\begin{equation}
  \min_{\omega \in \Omega_s}\, t_\omega \;=\; t_{\,s \cdot \omega_{\max}(H, W)},
\end{equation}
which is Eq.~\eqref{eq:stage-transition}.
\end{proof}

\clearpage
\section{Empirical Validation of the Per-frequency Activation Criterion}
\label{sec:motivation}

We empirically validate Proposition 1 with a simple \textit{spectral noise passthrough} experiment.
While maintaining full-resolution processing at all denoising steps, for each frequency $\omega \in \Omega$, if the current timestep $t$ is before its activation time $t_\omega$ (i.e., $t > t_\omega$), we replace that frequency component with $t \cdot T_\Phi(\boldsymbol{\epsilon})^{(\omega)}$, where $\boldsymbol{\epsilon} \sim \mathcal{N}(0, I)$ is the initial noise at $t = 1$. We use the Discrete Cosine Transform (DCT) as $T_\Phi$ in this experiment and sweep over multiple $\delta$ values. As described in Sec.~\ref{sec:method:schedule:signal}, a larger $\delta$ corresponds to transitioning to the high-resolution stage later in denoising with more speedup, while a smaller $\delta$ indicates transitioning to high-resolution earlier in denoising with less speedup.

From Fig. \ref{fig:spectral-passthrough}, we observe that at smaller $\delta \in [0.0001, 0.001]$, there is almost no observable difference compared to native full-resolution generation. As $\delta$ increases, we see a clear transition to blurry and lower-quality images around $\delta=0.05$, and larger values further degrade generation quality. This finding provides empirical support for the spectral autoregression property of diffusion models and Proposition 1, which serves as the foundation of our Spectral Progressive Diffusion framework and the $\delta$-optimal resolution schedule. It also implies that the Gaussian modelling assumption in Proposition 1 serves as a reasonable approximation for subsequent analytical resolution schedule derivations.

\begin{figure}[h!]
    \centering
    \includegraphics[width=\linewidth]{supp_figures/spectral_passthrough.pdf}
    \caption{\textbf{Spectral noise passthrough experiment.} At smaller $\delta \in [0.0001, 0.001]$, there is almost no observable difference compared to native full-resolution generation. As larger $\delta$ values cause high-frequency replacement to persist later in the denoising trajectory, we observe increasingly blurry and distorted results (i.e., ghosting artifacts and ``CHOOLBUS'' instead of ``SCHOOLBUS''). }
    \label{fig:spectral-passthrough}
\end{figure}


\clearpage
\section{Additional Details on Power Spectrum Measurement}
\label{sec:power_spec}
To remain consistent with the zero-mean Gaussian modelling assumption in Sec.~\ref{sec:method:schedule:signal}, we center each data sample before measuring the radially averaged power spectrum of the VAE latent space of FLUX.1-dev~\citep{black2024flux} and Z-Image~\citep{zimage} for image generation, and the VAE latent space of WAN 2.1~\citep{wan} for video generation.
We fit the power-law model $P_\omega \approx A \cdot |\omega|^{-\beta}$ and report the fitted $A$ and $\beta$ values for image latents, video latents, and image pixels in Table~\ref{tab:power-spec-fit}.
These fitted values instantiate the per-frequency signal power $P_\omega$ in the activation-time and transition-time formulas in Sec.~\ref{sec:method:schedule:signal}, which we use to compute the $\delta$-optimal resolution schedule for each model family.

For image latents, we run the FLUX.1-dev VAE on 100K images from the Aesthetics-Train-V2~\citep{zhang2025diffusion4k} dataset, center-cropped and resized to $1024^2$ resolution. We then average power over channels and radial frequency bins before fitting the spectrum. For video latents, we use the same procedure with the WAN 2.1 VAE on 10K videos from the VChitect-T2V-Dataverse~\citep{fan2025vchitect} dataset, averaging over channels, latent frames, and radial frequency bins to estimate the spatial power spectrum. We further measure the pixel-space power spectrum on images from the Aesthetics-Train-V2 dataset, center-cropped and resized to $512^2$, the native resolution supported by the PixelGen model.

\begin{table}[h!]
  \centering
  \small
  \caption{\textbf{Power-law fits to measured power spectra.} We fit $P_\omega \approx A \cdot |\omega|^{-\beta}$ on centered samples for image latents, video latents, and image pixels.}
  \label{tab:power-spec-fit}
  \resizebox{\linewidth}{!}{
  \begin{tabular}{lcccc}
    \toprule
    Spectrum & $\beta$ & $A$ & $R^2$ & Samples \\
    \midrule
    FLUX.1-dev (Z-Image) latent & 1.9155 & 203.62 & 0.978 & 100K images \\
    WAN latent & 2.4227 & 219.48 & 0.999 & 10K videos \\
    Pixel $512^2$ & 2.4493 & 3745.46 & 1.000 & 100K images \\
    \bottomrule
  \end{tabular}
  }
\end{table}


\clearpage
\section{Additional Details and Results on Image Generation}
\label{sec:image_generation}


\subsection{Training-free Acceleration of Latent-space Image Generation}
\label{subsec:image_latent}

For training-free inference acceleration of image generation, we use NVIDIA A100 GPUs with 80~GB of VRAM to match RALU~\citep{ralu}'s evaluation protocol. For training-free inference acceleration of video generation, as well as image editing and ablation studies, we use NVIDIA H100 GPUs with 80~GB of VRAM. For training-free acceleration of latent-space image generation, we use FLUX.1-dev~\citep{black2024flux} as the base model and follow the evaluation protocol of RALU~\citep{ralu}. Table~\ref{tab:flux-extended} extends Table~\ref{tab:flux-main} with additional training-free acceleration baselines from RALU. We organize the comparison into three wall-clock speedup tiers:

\begin{itemize}
    \item \textbf{Small speedup}: the 50-step native full-resolution baseline, compared against RALU (50 steps) and our method (40 steps), yielding approximately $1.7\times$ speedup;
    \item \textbf{Medium speedup}: the 10-step native full-resolution baseline, used as an iso-compute reference for RALU (15 steps) and our method (12 steps), yielding approximately $5\times$ speedup; and
    \item \textbf{Large speedup}: the 7-step native full-resolution baseline, used as an iso-compute reference for RALU (10 steps) and our method (10 steps), yielding approximately $7\times$ speedup.
\end{itemize}
All speedups are calculated relative to the wall-clock latency of the 50-step FLUX.1-dev baseline on a single A100 GPU, following RALU~\citep{ralu}. As in Table~\ref{tab:flux-main}, we evaluate image quality using CLIP-IQA~\citep{wang2023clipiqa} and NIQE~\citep{mittal2013niqe}, prompt alignment using T2I-CompBench~\citep{t2icompbench} and GenEval~\citep{geneval}, and overall quality using ImageReward~\citep{imagereward}.
Following RALU's evaluation protocol, ImageReward, CLIP-IQA, and NIQE are averaged over 5,000 images generated from MS-COCO validation prompts. For compositional alignment, GenEval is evaluated with four random seeds per prompt on its object-focused prompt suite, and T2I-CompBench is evaluated on its spatial, non-spatial, and complex prompt subsets with four random seeds per prompt. Unless otherwise specified, all metrics are computed using their official or default evaluation settings.

We keep the error threshold fixed at $\delta=0.01$, use $T_\Phi = \text{DCT}$, and compare $S=2$ and $S=3$ resolution-stage settings. To align with the multi-resolution training distribution of FLUX.1-dev, we use latent dimensions $64^2$ ($s_1 = 0.5$) and $128^2$ $(s_2 = 1.0)$ (pixel dimensions $512^2$ and $1024^2$) for $S=2$, and latent dimensions $32^2$ $(s_1 = 0.25)$, $64^2$ $(s_2 = 0.5)$, and $128^2$ ($s_3 = 1.0$) (pixel dimensions $256^2$, $512^2$, and $1024^2$) for $S=3$.

By combining spectral transformations for progressive-resolution generation with our derived resolution schedule, our method achieves a better speed-quality tradeoff than the reduced-step FLUX.1-dev baselines and existing temporal and spatial diffusion acceleration methods, including progressive-resolution methods such as RALU~\citep{ralu}. In the largest speedup tier, our $S=3$ setting reaches $7.09\times$ wall-clock speedup and $7.36\times$ FLOPs speedup while preserving competitive image quality. We further provide extended qualitative baseline comparisons in Figures~\ref{fig:flux-1}--\ref{fig:flux-3}. FLUX.1-dev with reduced steps degrades image quality and exhibits over-saturation artifacts, while RALU introduces noticeable noise artifacts in reduced-step settings.


\begin{table*}[t!]
    \centering
    \caption{\textbf{Extended quantitative comparisons on training-free latent-space image generation.} We evaluate FLUX.1-dev at $1024^2$ resolution and group methods into the same $1\times$, $5\times$, and $7\times$ wall-clock speedup tiers as Table~\ref{tab:flux-main}. Rows shared with Table~\ref{tab:flux-main} use the same values. Speedup is calculated relative to the 50-step FLUX.1-dev baseline wall-clock runtime on a single A100 GPU. \textcolor{blue}{\textbf{T}} and \textcolor{red}{\textbf{S}} denote temporal and spatial acceleration methods, respectively. Baseline metrics are copied from RALU~\citep{ralu} under the same evaluation protocol.}
    \label{tab:flux-extended}
    \resizebox{\textwidth}{!}{
    \begin{scriptsize}
    \begin{tabular}{cccccccccc}
         \toprule
        \multirow{2}{*}{Method} & \multirow{2}{*}{Accel.} & \multirow{2}{*}{Latency (s) $\downarrow$} & \multirow{2}{*}{Speedup (s) $\uparrow$} & \multirow{2}{*}{TFLOPs $\downarrow$} & Overall & \multicolumn{2}{c}{Image quality} & \multicolumn{2}{c}{Text alignment} \\
        \cmidrule(lr){6-10}
        & & & & & ImageReward $\uparrow$ & CLIP-IQA $\uparrow$ & NIQE $\downarrow$ & T2I-Comp. $\uparrow$ & GenEval $\uparrow$  \\
        \midrule
        FLUX (50 steps) & - & 25.1 & 1.00$\times$ & 2991.01 & \textbf{1.095} & 0.707 & 6.75 & \textbf{0.634} & \textbf{0.698}  \\
        RALU~\citep{ralu}   & \textcolor{red}{\textbf{S}}  & 15.85 & 1.58$\times$ & \underline{1749.94} & 1.028 & \underline{0.712} & \textbf{6.07} & 0.613 & 0.648 \\
        \textbf{Ours ($S=2$)}     & \textcolor{red}{\textbf{S}}       & \underline{15.16} & \underline{1.66}$\times$ & 1755.22 & \underline{1.049} & \textbf{0.719} & 6.43 & \underline{0.617} & \underline{0.654} \\
        \textbf{Ours ($S=3$)}     & \textcolor{red}{\textbf{S}}       & \textbf{14.54} & \textbf{1.73}$\times$ & \textbf{1672.04} & 1.015 & 0.711 & \underline{6.33} & 0.593 & 0.640 \\
        \cmidrule(lr){1-10}
        FLUX (10 steps) & \textcolor{blue}{\textbf{T}} & 5.18 & 4.84$\times$ & 610.02 & 0.981 & 0.679 & 6.93 & 0.618 & 0.647 \\
        $\Delta$-DiT~\citep{deltadit}      & \textcolor{blue}{\textbf{T}} & 7.42 & 3.38$\times$ & 772.10 & 0.102 & 0.487 & 9.60 & 0.306 & 0.397 \\
        ToCa~\citep{toca}                  & \textcolor{blue}{\textbf{T}} & 15.5 & 1.62$\times$ & 601.12 & -1.827 & 0.253 & 10.6 & 0.259 & 0.137 \\
        TeaCache~\citep{teacache}          & \textcolor{blue}{\textbf{T}} & 5.23 & 4.80$\times$ & 610.59 & 0.944 & 0.665 & 7.92 & 0.620 & 0.647 \\
        TaylorSeer~\citep{taylorseer}      & \textcolor{blue}{\textbf{T}} & 9.34 & 2.69$\times$ & 556.72 & 0.972 & 0.684 & 6.77 & 0.594 & 0.619 \\
        Bottleneck~\citep{bottleneck}      & \textcolor{red}{\textbf{S}}  & 5.37 & 4.67$\times$ & 571.23 & 0.889 & 0.661 & 9.16 & 0.620 & \textbf{0.687} \\
        RALU~\citep{ralu}                  & \textcolor{red}{\textbf{S}}  & 5.04 & 4.98$\times$ & 540.47 & 1.022 & \underline{0.700} & \textbf{6.43} & \textbf{0.626} & 0.652 \\
        \textbf{Ours ($S=2$)}     & \textcolor{red}{\textbf{S}}       & \underline{4.35} & \underline{5.77}$\times$ & \underline{500.34} & \textbf{1.059} & 0.696 & 6.69 & \underline{0.624} & \underline{0.655} \\
        \textbf{Ours ($S=3$)}     & \textcolor{red}{\textbf{S}}       & \textbf{4.12} & \textbf{6.09}$\times$ & \textbf{469.15} & \underline{1.042} & \textbf{0.701} & \underline{6.53} & 0.623 & 0.637 \\
        \cmidrule(lr){1-10}
        FLUX (7 steps) & \textcolor{blue}{\textbf{T}} & 3.79 & 6.62$\times$ & 431.45 & 0.920 & 0.660 & 8.25 & 0.594 & 0.583 \\
        TeaCache~\citep{teacache}          & \textcolor{blue}{\textbf{T}} & 4.21 & 5.96$\times$ & 431.83 & 0.733 & 0.623 & 13.7 & 0.599 & 0.594 \\
        TaylorSeer~\citep{taylorseer}      & \textcolor{blue}{\textbf{T}} & 7.00 & 3.59$\times$ & 431.74 & 0.660 & 0.646 & 9.43 & 0.514 & 0.446 \\
        Bottleneck~\citep{bottleneck}      & \textcolor{red}{\textbf{S}}  & 3.78 & 6.64$\times$ & 431.52 & 0.792 & 0.631 & 8.71 & 0.605 & \underline{0.672} \\
        RALU~\citep{ralu}                  & \textcolor{red}{\textbf{S}}  & 3.75 & 6.69$\times$ & \underline{426.01} & 0.999 & 0.681 & 6.87 & \textbf{0.633} & \textbf{0.682} \\
        \textbf{Ours ($S=2$)}     & \textcolor{red}{\textbf{S}}       & \underline{3.70} & \underline{6.78}$\times$ & 427.03 & \textbf{1.039} & \underline{0.689} & \underline{6.78} & 0.620 & 0.667 \\
        \textbf{Ours ($S=3$)}     & \textcolor{red}{\textbf{S}}       & \textbf{3.54} & \textbf{7.09}$\times$ & \textbf{406.24} & \underline{1.015} & \textbf{0.694} & \textbf{5.99} & \underline{0.627} & 0.637 \\
        \bottomrule
    \end{tabular}
    \end{scriptsize}}
\end{table*}


\subsection{Spectral-transformation-based Fine-tuning for Latent- and Pixel-Space Image Generation}
\label{subsec:image_pixel}
\paragraph{Additional latent-space fine-tuning results.} Since FLUX.1-dev is a guidance-distilled model and its base model is not open-sourced, we additionally conduct latent-space image generation fine-tuning experiments on Z-Image~\citep{zimage}, which employs the same VAE architecture as FLUX.1-dev and hence the same $\delta$-optimal resolution schedules. We adopt Low-Rank Adaptation (LoRA) of rank 32 and train for 2000 steps following the spectral-transformation-based fine-tuning procedure in Sec.~\ref{sec:method:lora}. For fair comparison with the full-resolution and training-free acceleration baselines, we fine-tune on synthetic images generated by full-resolution Z-Image with 50 denoising steps using a set of 5K MS-COCO prompts that is disjoint from the 5K MS-COCO validation prompts used for evaluation. Training is conducted on our internal cluster of 8 NVIDIA H100 GPUs, each with 80~GB of VRAM, at an effective batch size of 8 and takes approximately 2 hours. We keep the optimizer and all other LoRA fine-tuning hyperparameters at the default settings of the \texttt{DiffSynth Studio} library, including a constant learning rate of 1e-4. Fine-tuning compute scales linearly with the number of sampled training images, training iterations, and batch size. In our experiments we fix the synthetic fine-tuning set to 5K prompts and train LoRA adapters for 2000 iterations, so increasing the fine-tuning dataset size would only affect compute through additional sampled training examples and does not change the inference-time cost of our method.

All latent-image generation experiments are conducted at a native resolution of $1024^2$ pixels, which is the default configuration for both FLUX.1-dev and Z-Image. The progressive resolution settings for the $S = 2$ and $S = 3$ cases are the same as in the FLUX.1-dev training-free acceleration experiments in Sec.~\ref{subsec:image_latent}.

Further model fine-tuning allows a more aggressive error tolerance $\delta$, leading to an even higher speedup. For fine-tuning experiments, we additionally evaluate $\delta=0.05$, which is larger than the default $\delta=0.01$. As in Sec.~\ref{subsec:image_latent}, we organize the comparison into three wall-clock speedup tiers, all calculated relative to the 50-step Z-Image baseline on a single H100 GPU (we adopt H100 wall-clock measurements for Z-Image and all subsequent pixel-space image generation and video generation experiments):

\begin{itemize}
    \item \textbf{Small speedup}: the 50-step native full-resolution Z-Image baseline, compared against our training-free acceleration and LoRA fine-tuned variants at $S=2$ and $S=3$ ($\delta=0.01$, 50 steps), as well as our LoRA fine-tuned variant with the configuration ($\delta=0.05$, $S=2$, 50 steps), yielding $1.65\times$ to $2.01\times$ speedup;
    \item \textbf{Medium speedup}: the 10-step native full-resolution Z-Image baseline, used as an iso-compute reference for our training-free acceleration and LoRA fine-tuned variants at $S=2$ and $S=3$ with reduced denoising steps, yielding approximately $5\times$ speedup; and
    \item \textbf{Large speedup}: the 7-step native full-resolution Z-Image baseline, used as an iso-compute reference for our training-free acceleration and LoRA fine-tuned variants at $S=2$ and $S=3$ with further reduced denoising steps, yielding approximately $7\times$ to $8\times$ speedup.
\end{itemize}
For quantitative evaluation, we use the same metric evaluation settings as Sec.~\ref{subsec:image_latent}, following the RALU evaluation protocol, and report the numbers in Table~\ref{tab:zimage-supp}. We show extensive qualitative comparisons in Figures~\ref{fig:zimage-1}--\ref{fig:zimage-3}. Across all three speedup tiers, our training-free acceleration variants outperform the reduced-step Z-Image baselines, and our spectral-transformation-based LoRA fine-tuning further closes the gap to the 50-step full-resolution baseline while attaining up to $7.81\times$ wall-clock speedup. In particular, our LoRA fine-tuned model at $S=3$ outperforms the training-free acceleration variant at $S=2$ while improving efficiency, and our LoRA fine-tuned model at $S=2$ either with reduced denoising steps or at a larger $\delta=0.05$ attains image quality comparable to the training-free acceleration $S=2$, 50-step setting, providing additional efficiency at no quality cost.


\begin{table*}[t!]
    \centering
    \caption{\textbf{Quantitative comparisons on latent-space image generation with fine-tuning.} We evaluate Z-Image at $1024^2$ resolution and group methods into wall-clock speedup tiers. Speedup is calculated relative to the 50-step Z-Image baseline wall-clock runtime on a single H100 GPU. \textcolor{blue}{\textbf{T}} and \textcolor{red}{\textbf{S}} denote temporal and spatial acceleration methods, respectively. All metrics are evaluated using the same protocol as Table~\ref{tab:flux-main}.}
    \label{tab:zimage-supp}
    \resizebox{\textwidth}{!}{
    \begin{scriptsize}
    \begin{tabular}{cccccccccc}
         \toprule
        \multirow{2}{*}{Method} & \multirow{2}{*}{Accel.} & \multirow{2}{*}{Latency (s) $\downarrow$} & \multirow{2}{*}{Speedup (s) $\uparrow$} & \multirow{2}{*}{TFLOPs $\downarrow$} & Overall & \multicolumn{2}{c}{Image quality} & \multicolumn{2}{c}{Text alignment} \\
        \cmidrule(lr){6-10}
        & & & & & ImageReward $\uparrow$ & CLIP-IQA $\uparrow$ & NIQE $\downarrow$ & T2I-Comp. $\uparrow$ & GenEval $\uparrow$  \\
        \midrule
        Z-Image (50 steps) & - & 21.25 & 1.00$\times$ & 4941.23 & 0.965 & 0.700 & 5.41 & 0.731 & 0.745  \\
        \cmidrule(lr){1-10}
        Z-Image (32 steps) & \textcolor{blue}{\textbf{T}} & 13.59 & 1.56$\times$ & 3166.62 & \underline{0.957} & 0.697 & \textbf{5.44} & 0.686 & 0.725  \\
        \textbf{Ours (TF, $\delta=0.01$, $S=2$)}     & \textcolor{red}{\textbf{S}}       & 12.90 & 1.65$\times$ & 3132.03 & 0.904 & 0.688 & 5.87 & 0.658 & \underline{0.730} \\
        \textbf{Ours (TF, $\delta=0.01$, $S=3$)}     & \textcolor{red}{\textbf{S}}       & \underline{12.19} & \underline{1.74$\times$} & \underline{2871.09} & 0.875 & 0.690 & \underline{5.59} & 0.650 & 0.682 \\
        \textbf{Ours (LoRA, $\delta=0.01$, $S=2$)}     & \textcolor{red}{\textbf{S}}       & 12.87 & 1.65$\times$ & 3132.03 & \textbf{0.982} & \textbf{0.699} & 5.72 & \textbf{0.725} & \textbf{0.731} \\
        \textbf{Ours (LoRA, $\delta=0.01$, $S=3$)}     & \textcolor{red}{\textbf{S}}       & 12.23 & 1.74$\times$ & \underline{2871.09} & 0.954 & \underline{0.697} & 5.75 & \underline{0.717} & 0.728 \\
        \textbf{Ours (LoRA, $\delta=0.05$, $S=2$)}     & \textcolor{red}{\textbf{S}}       & \textbf{10.59} & \textbf{2.01$\times$} & \textbf{2436.19} & 0.919 & 0.684 & 6.12 & 0.661 & 0.718 \\
        \cmidrule(lr){1-10}
        Z-Image (10 steps) & \textcolor{blue}{\textbf{T}} & 4.26 & 4.99$\times$ & 997.68 & 0.851 & 0.659 & 5.95 & \underline{0.678} & \underline{0.705} \\
        \textbf{Ours (TF, $\delta=0.01$, $S=2$)}     & \textcolor{red}{\textbf{S}}       & \underline{4.22} & \underline{5.04$\times$} & \underline{962.95} & \underline{0.860} & \underline{0.668} & 6.17 & 0.677 & 0.693 \\
        \textbf{Ours (TF, $\delta=0.01$, $S=3$)}     & \textcolor{red}{\textbf{S}}       & \textbf{4.01} & \textbf{5.30$\times$} & \textbf{875.97} & 0.827 & 0.662 & \underline{5.78} & 0.676 & 0.655 \\
        \textbf{Ours (LoRA, $\delta=0.01$, $S=2$)}     & \textcolor{red}{\textbf{S}}       &  4.24 & 5.01$\times$ & \underline{962.95} & \textbf{0.923} & \textbf{0.683} & \textbf{5.58} & \textbf{0.706} & \textbf{0.738} \\
        \cmidrule(lr){1-10}
        Z-Image (7 steps) & \textcolor{blue}{\textbf{T}} & 3.56 & 5.97$\times$ & 701.92 & 0.763 & 0.631 & \textbf{6.14} & 0.648 & 0.667 \\
        \textbf{Ours (TF, $\delta=0.01$, $S=2$)}     & \textcolor{red}{\textbf{S}}       & 2.75 & 7.73$\times$ & \underline{678.77} & \underline{0.804} & \underline{0.641} & 6.46 & \underline{0.665} & \underline{0.680} \\
        \textbf{Ours (TF, $\delta=0.01$, $S=3$)}     & \textcolor{red}{\textbf{S}}       & \textbf{2.63} & \textbf{8.08$\times$} & \textbf{609.18} & 0.759 & 0.630 & \underline{6.16} & 0.658 & 0.645 \\
        \textbf{Ours (LoRA, $\delta=0.01$, $S=2$)}     & \textcolor{red}{\textbf{S}}       & \underline{2.72} & \underline{7.81$\times$} & \underline{678.77} & \textbf{0.918} & \textbf{0.658} & 6.57 & \textbf{0.683} & \textbf{0.744} \\
        \bottomrule
    \end{tabular}
    \end{scriptsize}
    }
\end{table*}


% \subsection{Spectral Transform-based Fine-tuning for Latent- and Pixel-Space Image Generation}
% \label{subsec:image_pixel}

\paragraph{Additional pixel-space fine-tuning results.} For pixel-space image generation, we evaluate our method on PixelGen-XXL/16 T2I~\citep{pixelgen} at its native resolution of $512^2$ pixels. We use the same fine-tuning data construction and LoRA training procedure as the Z-Image latent-space experiments above, with synthetic fine-tuning images generated by full-resolution PixelGen with 25 denoising steps from a set of 5K MS-COCO prompts disjoint from the 5K MS-COCO validation prompts used for evaluation. Since PixelGen is pretrained with $x_0$ prediction rather than flow matching, we convert the spectral-transform velocity targets from the fine-tuning procedure in Sec.~\ref{sec:method:lora} to $x_0$ targets and supervise the LoRA adapter with the converted $x_0$.

We use $S=2$ with $\delta=0.01$ and, following Sec.~\ref{subsec:image_latent}, organize the comparison into two wall-clock speedup tiers, all calculated relative to the 25-step PixelGen baseline on a single H100 GPU:

\begin{itemize}
    \item \textbf{Small speedup}: the 25-step native full-resolution PixelGen baseline, used as the reference for the reduced 13-step PixelGen variant and our training-free acceleration and LoRA fine-tuned variants at $S=2$ (25 steps), yielding approximately $1.5\times$ to $1.9\times$ speedup;
    \item \textbf{Medium speedup}: the 8-step native full-resolution PixelGen variant, used as an iso-compute reference for our training-free acceleration and LoRA fine-tuned variants at $S=2$ with reduced denoising steps, yielding approximately $2\times$ to $3\times$ speedup.
\end{itemize}

For quantitative evaluation, we follow the same protocol used in latent-space image generation and report the numbers in Table~\ref{tab:pixelgen-main}. We show extensive qualitative comparisons in Figures~\ref{fig:pixelgen-1}--\ref{fig:pixelgen-3}. Because PixelGen does not natively support a high-fidelity $256^2$ resolution generation, our training-free acceleration variant falls short of the reduced-step PixelGen baselines on overall image quality across both speedup tiers. Our spectral-transformation-based LoRA fine-tuning closes this gap: at the same FLOPs as the training-free acceleration variant, the LoRA fine-tuned model matches or surpasses the reduced-step PixelGen baseline on most image-quality and prompt-alignment metrics, including ImageReward, CLIP-IQA, and T2I-CompBench. The qualitative comparisons in Figures~\ref{fig:pixelgen-1}--\ref{fig:pixelgen-3} corroborate this trend, showing that our LoRA fine-tuned model's outputs match the visual quality of the 25-step native-resolution PixelGen baseline at substantially higher throughput.


\begin{table*}[t!]
    \centering
    \caption{\textbf{Quantitative comparisons on pixel-space image generation.} The number in parentheses next to PixelGen indicates the total number of inference steps. $\uparrow / \downarrow$ denotes that a higher / lower metric is favourable. Speedup is calculated relative to the 25-step PixelGen baseline wall-clock runtime on a single H100 GPU. \textcolor{blue}{\textbf{T}} and \textcolor{red}{\textbf{S}} denote temporal and spatial acceleration, respectively.}
    \label{tab:pixelgen-main}
    \resizebox{\textwidth}{!}{
    \begin{scriptsize}
    \begin{tabular}{cccccccccc}
         \toprule
        \multirow{2}{*}{Method} & \multirow{2}{*}{Accel.} & \multirow{2}{*}{Latency (s) $\downarrow$} & \multirow{2}{*}{Speedup (s) $\uparrow$} & \multirow{2}{*}{TFLOPs $\downarrow$} & Overall & \multicolumn{2}{c}{Image quality} & \multicolumn{2}{c}{Text alignment} \\
        \cmidrule(lr){6-10}
        & & & & & ImageReward $\uparrow$ & CLIP-IQA $\uparrow$ & NIQE $\downarrow$ & T2I-Comp. $\uparrow$ & GenEval $\uparrow$  \\
        \midrule
        PixelGen (25 steps) & - & 0.48 & 1.00$\times$ & 65.36 & 0.921 & 0.734 & 5.95 & 0.574 & 0.794 \\
        \cmidrule(lr){1-10}
        PixelGen (13 steps) & \textcolor{blue}{\textbf{T}} & \textbf{0.25} & \textbf{1.92$\times$} & \underline{34.16} & \underline{0.886} & \underline{0.726} & \underline{5.89} & \underline{0.571} & \underline{0.781} \\
        \textbf{Ours (TF, $\delta=0.01$, $S = 2$)}     & \textcolor{red}{\textbf{S}}       & \underline{0.30} & \underline{1.60$\times$} & \textbf{33.72} & 0.799 & 0.718 & 6.10 & 0.568 & \textbf{0.782} \\
        \textbf{Ours (LoRA, $\delta=0.01$, $S = 2$)}   & \textcolor{red}{\textbf{S}}       & 0.31 & 1.55$\times$ & \textbf{33.72} & \textbf{0.913} & \textbf{0.728} & \textbf{5.87} & \textbf{0.580} & 0.776 \\
        \cmidrule(lr){1-10}
        PixelGen (8 steps) & \textcolor{blue}{\textbf{T}} & \textbf{0.16} & \textbf{3.00$\times$} & \underline{21.16} & \underline{0.858} & \underline{0.715} & \textbf{5.79} & \underline{0.572} & 0.756 \\
        \textbf{Ours (TF, $\delta=0.01$, $S = 2$)}     & \textcolor{red}{\textbf{S}}       & \underline{0.21} & \underline{2.29$\times$} & \textbf{20.75} & 0.779 & 0.713 & 6.02 & 0.567 & \underline{0.765} \\
        \textbf{Ours (LoRA, $\delta=0.01$, $S = 2$)}   & \textcolor{red}{\textbf{S}}       & 0.22 & 2.18$\times$ & \textbf{20.75} & \textbf{0.908} & \textbf{0.724} & \underline{5.85} & \textbf{0.577} & \textbf{0.770} \\
        \bottomrule
    \end{tabular}
    \end{scriptsize}
    }
\end{table*}


\clearpage
\begin{figure}
    \centering
    \includegraphics[width=\linewidth]{supp_figures/flux_supp_1.pdf}
    \caption{\textbf{Qualitative comparisons on latent-space image generation.} We compare our method against default-step generation, reduced-step native-resolution generation on FLUX.1-dev \citep{black2024flux}, and RALU~\citep{ralu}, a state-of-the-art acceleration baseline matched to similar speedups. Our method outperforms both baselines. FLUX.1-dev with reduced steps degrades image quality and exhibits over-saturation artifacts, while RALU introduces noticeable noise artifacts.}
    \label{fig:flux-1}
\end{figure}

\clearpage
\begin{figure}
    \centering
    \includegraphics[width=\linewidth]{supp_figures/flux_supp_2.pdf}
    \caption{\textbf{Qualitative comparisons on latent-space image generation.} We compare our method against default-step generation, reduced-step native-resolution generation on FLUX.1-dev \citep{black2024flux}, and RALU~\citep{ralu}, a state-of-the-art acceleration baseline matched to similar speedups. Our method outperforms both baselines. FLUX.1-dev with reduced steps degrades image quality and exhibits over-saturation artifacts, while RALU introduces noticeable noise artifacts.}
    \label{fig:flux-2}
\end{figure}

\clearpage
\begin{figure}
    \centering
    \includegraphics[width=\linewidth]{supp_figures/flux_supp_3.pdf}
    \caption{\textbf{Qualitative comparisons on latent-space image generation.} We compare our method against default-step generation, reduced-step native-resolution generation on FLUX.1-dev \citep{black2024flux}, and RALU~\citep{ralu}, a state-of-the-art acceleration baseline matched to similar speedups. Our method outperforms both baselines. FLUX.1-dev with reduced steps degrades image quality and exhibits over-saturation artifacts, while RALU introduces noticeable noise artifacts.}
    \label{fig:flux-3}
\end{figure}

\iffalse
\clearpage
\begin{figure}
    \centering
    \includegraphics[width=\linewidth]{supp_figures/flux_supp_4.pdf}
    \caption{\textbf{Qualitative comparisons on latent-space image generation.} We compare our method against default-step generation, reduced-step native-resolution generation on FLUX.1-dev \citep{black2024flux}, and RALU~\citep{ralu}, a state-of-the-art acceleration baseline matched to similar speedups. Our method outperforms both baselines. FLUX.1-dev with reduced steps degrades image quality and exhibits over-saturation artifacts, while RALU introduces noticeable noise artifacts.}
    \label{fig:flux-4}
\end{figure}

\clearpage
\begin{figure}
    \centering
    \includegraphics[width=\linewidth]{supp_figures/flux_supp_5.pdf}
    \caption{\textbf{Qualitative comparisons on latent-space image generation.} We compare our method against default-step generation, reduced-step native-resolution generation on FLUX.1-dev \citep{black2024flux}, and RALU~\citep{ralu}, a state-of-the-art acceleration baseline matched to similar speedups. Our method outperforms both baselines. FLUX.1-dev with reduced steps degrades image quality and exhibits over-saturation artifacts, while RALU introduces noticeable noise artifacts.}
    \label{fig:flux-5}
\end{figure}
\fi

\clearpage
\begin{figure}
     \centering
     \includegraphics[width=\linewidth]{supp_figures/zimage_supp_1.pdf}
     \caption{\textbf{Qualitative comparisons on latent-space image generation (fine-tuned).} We compare our method against default-step generation, reduced-step native-resolution generation on Z-Image \citep{zimage} matched to similar speedups. Our fine-tuned model (Ours$^*$) achieves even higher image quality compared to our training-free acceleration variant and outperforms the reduced-step baseline.}
     \label{fig:zimage-1}
\end{figure}

\clearpage
\begin{figure}
     \centering
     \includegraphics[width=\linewidth]{supp_figures/zimage_supp_2.pdf}
     \caption{\textbf{Qualitative comparisons on latent-space image generation (fine-tuned).} We compare our method against default-step generation, reduced-step native-resolution generation on Z-Image \citep{zimage} matched to similar speedups. Our fine-tuned model (Ours$^*$) achieves even higher image quality compared to our training-free acceleration variant and outperforms the reduced-step baseline.}
     \label{fig:zimage-2}
\end{figure}

\clearpage
\begin{figure}
     \centering
     \includegraphics[width=\linewidth]{supp_figures/zimage_supp_3.pdf}
     \caption{\textbf{Qualitative comparisons on latent-space image generation (fine-tuned).} We compare our method against default-step generation, reduced-step native-resolution generation on Z-Image \citep{zimage} matched to similar speedups. Our fine-tuned model (Ours$^*$) achieves even higher image quality compared to our training-free acceleration variant and outperforms the reduced-step baseline.}
     \label{fig:zimage-3}
\end{figure}


\iffalse
\clearpage
\begin{figure}
     \centering
     \includegraphics[width=\linewidth]{supp_figures/zimage_supp_4.pdf}
     \caption{\textbf{Qualitative comparisons on latent-space image generation (fine-tuned).} We compare our method against default-step generation, reduced-step native-resolution generation on Z-Image \citep{zimage} matched to similar speedups. Our fine-tuned model (Ours$^*$) achieves even higher image quality compared to our training-free acceleration variant and outperforms the reduced-step baseline.}
     \label{fig:zimage-4}
\end{figure}

\clearpage
\begin{figure}
     \centering
     \includegraphics[width=\linewidth]{supp_figures/zimage_supp_5.pdf}
     \caption{\textbf{Qualitative comparisons on latent-space image generation (fine-tuned).} We compare our method against default-step generation, reduced-step native-resolution generation on Z-Image \citep{zimage} matched to similar speedups. Our fine-tuned model (Ours$^*$) achieves even higher image quality compared to our training-free acceleration variant and outperforms the reduced-step baseline.}
     \label{fig:zimage-5}
\end{figure}
\fi

\clearpage
\begin{figure}
    \centering
    \includegraphics[width=0.67\linewidth]{supp_figures/pixelgen_supp_1.pdf}
    \caption{\textbf{Qualitative comparisons on pixel-space image generation.} We compare our method against default-step generation and reduced-step native-resolution generation on PixelGen~\citep{pixelgen}, matched to comparable speedups. An asterisk (Ours$^*$) marks the fine-tuned model. Our method achieves similar quality to full-resolution generation while attaining higher speedups.}
    \label{fig:pixelgen-1}
\end{figure}

\clearpage
\begin{figure}
    \centering
    \includegraphics[width=0.67\linewidth]{supp_figures/pixelgen_supp_2.pdf}
    \caption{\textbf{Qualitative comparisons on pixel-space image generation.} We compare our method against default-step generation and reduced-step native-resolution generation on PixelGen~\citep{pixelgen}, matched to comparable speedups. An asterisk (Ours$^*$) marks the fine-tuned model. Our method achieves similar quality to full-resolution generation while attaining higher speedups.}
    \label{fig:pixelgen-2}
\end{figure}

\clearpage
\begin{figure}
    \centering
    \includegraphics[width=0.67\linewidth]{supp_figures/pixelgen_supp_3.pdf}
    \caption{\textbf{Qualitative comparisons on pixel-space image generation.} We compare our method against default-step generation and reduced-step native-resolution generation on PixelGen~\citep{pixelgen}, matched to comparable speedups. An asterisk (Ours$^*$) marks the fine-tuned model. Our method achieves similar quality to full-resolution generation while attaining higher speedups.}
    \label{fig:pixelgen-3}
\end{figure}

\iffalse
\clearpage
\begin{figure}
    \centering
    \includegraphics[width=0.67\linewidth]{supp_figures/pixelgen_supp_4.pdf}
    \caption{\textbf{Qualitative comparisons on pixel-space image generation.} We compare our method against default-step generation and reduced-step native-resolution generation on PixelGen~\citep{pixelgen}, matched to comparable speedups. An asterisk (Ours$^*$) marks the fine-tuned model. Our method achieves similar quality to reduced-step generation while attaining higher speedups.}
    \label{fig:pixelgen-4}
\end{figure}


\clearpage
\begin{figure}
    \centering
    \includegraphics[width=0.67\linewidth]{supp_figures/pixelgen_supp_5.pdf}
    \caption{\textbf{Qualitative comparisons on pixel-space image generation.} We compare our method against default-step generation and reduced-step native-resolution generation on PixelGen~\citep{pixelgen}, matched to comparable speedups. An asterisk (Ours$^*$) marks the fine-tuned model. Our method achieves similar quality to reduced-step generation while attaining higher speedups.}
    \label{fig:pixelgen-5}
\end{figure}
\fi


\clearpage
\section{Additional Video Generation Results}
\label{sec:video_generation}

\paragraph{Training-free acceleration of latent-space video generation.} For training-free inference acceleration of video generation, we use our internal cluster of NVIDIA H100 GPUs with 80~GB of VRAM. For quantitative evaluation, we generate 200 videos using 200 test prompts and report the standard VBench scores across six evaluation dimensions, including Subject Consistency, Background Consistency, Motion Smoothness, Dynamic Degree, Aesthetic Quality, and Image Quality, in Table~\ref{tab:wan_vbench_720p}. All videos are generated in 720P resolution. For additional video results and qualitative comparisons, please see the project website at~\url{https://howardxiao.ca/speed/}.


\clearpage
\section{Extended Ablation Studies}
\label{sec:ablation}

We perform extended ablation studies on latent-space image generation. Specifically, we vary the error tolerance $\delta$, the number of resolution stages $S$, and the spectral transformation $T_\Phi$, and report wall-clock latency, speedup, and ImageReward in Table~\ref{tab:ablation-quan}, with qualitative comparisons in Figures~\ref{fig:ablation-T} to \ref{fig:ablation-s}.

\begin{table*}[t!]
    \centering
    \caption{\textbf{Ablation studies on the spectral transformation $T_\Phi$, error tolerance $\delta$, and number of resolution stages $S$.} Speedup is calculated relative to the 7-step FLUX.1-dev baseline wall-clock runtime on a single H100 GPU.}
    \label{tab:ablation-quan}
    \centering
    \begin{tiny}
    \begin{tabular}{cccccc}
        \toprule
        $T_\Phi$ & $\delta$ & $S$ & Latency (s) $\downarrow$ & Speedup (s) $\uparrow$ & ImageReward $\uparrow$ \\
        \midrule
        FLUX (7 steps) & - & - & 2.08 & 1.00$\times$ & 0.920 \\
        \midrule
        DCT & 0.001 & 2 & \textbf{2.08} & \textbf{1.00$\times$} & 1.032 \\
        DWT & 0.001 & 2 & \underline{2.08} & \textbf{1.00$\times$} & \textbf{1.038} \\
        FFT & 0.001 & 2 & 2.08 & \textbf{1.00$\times$} & \underline{1.036} \\
        \midrule
        DCT & 0.001 & 3 & \textbf{2.05} & \textbf{1.01$\times$} & \underline{1.034} \\
        DWT & 0.001 & 3 & \underline{2.06} & \textbf{1.01$\times$} & \underline{1.034} \\
        FFT & 0.001 & 3 & 2.06 & \textbf{1.01$\times$} & \textbf{1.040} \\
        \midrule
        DCT & 0.001 & 4 & \textbf{2.05} & \textbf{1.02$\times$} & \textbf{1.008} \\
        DWT & 0.001 & 4 & \underline{2.05} & \underline{1.01$\times$} & \underline{1.002} \\
        FFT & 0.001 & 4 & 2.06 & \underline{1.01$\times$} & 0.992 \\
        \midrule
        DCT & 0.002 & 2 & \textbf{2.07} & \textbf{1.01$\times$} & \underline{1.032} \\
        DWT & 0.002 & 2 & \underline{2.08} & \underline{1.00$\times$} & 1.031 \\
        FFT & 0.002 & 2 & 2.08 & \underline{1.00$\times$} & \textbf{1.036} \\
        \midrule
        DCT & 0.002 & 3 & \textbf{2.05} & \textbf{1.02$\times$} & \underline{1.034} \\
        DWT & 0.002 & 3 & \underline{2.06} & \underline{1.01$\times$} & \underline{1.034} \\
        FFT & 0.002 & 3 & 2.06 & \underline{1.01$\times$} & \textbf{1.040} \\
        \midrule
        DCT & 0.002 & 4 & \textbf{2.05} & \textbf{1.02$\times$} & \textbf{1.008} \\
        DWT & 0.002 & 4 & \underline{2.05} & \underline{1.01$\times$} & \underline{1.002} \\
        FFT & 0.002 & 4 & 2.06 & \underline{1.01$\times$} & 0.992 \\
        \midrule
        DCT & 0.005 & 2 & \textbf{1.91} & \textbf{1.09$\times$} & \underline{1.019} \\
        DWT & 0.005 & 2 & \underline{1.92} & \textbf{1.09$\times$} & \textbf{1.026} \\
        FFT & 0.005 & 2 & 1.92 & \underline{1.08$\times$} & 1.014 \\
        \midrule
        DCT & 0.005 & 3 & \textbf{1.87} & \textbf{1.12$\times$} & \underline{1.009} \\
        DWT & 0.005 & 3 & \underline{1.88} & \underline{1.11$\times$} & \underline{1.009} \\
        FFT & 0.005 & 3 & 1.88 & \underline{1.11$\times$} & \textbf{1.028} \\
        \midrule
        DCT & 0.005 & 4 & \textbf{1.87} & \textbf{1.12$\times$} & \underline{1.005} \\
        DWT & 0.005 & 4 & \underline{1.87} & \underline{1.11$\times$} & \textbf{1.008} \\
        FFT & 0.005 & 4 & 1.87 & \underline{1.11$\times$} & 0.994 \\
        \midrule
        DCT & 0.01 & 2 & \textbf{1.75} & \textbf{1.19$\times$} & \textbf{1.039} \\
        DWT & 0.01 & 2 & \underline{1.76} & \underline{1.18$\times$} & \underline{1.025} \\
        FFT & 0.01 & 2 & 1.76 & \underline{1.18$\times$} & 1.016 \\
        \midrule
        DCT & 0.01 & 3 & \textbf{1.71} & \textbf{1.22$\times$} & 1.015 \\
        DWT & 0.01 & 3 & \underline{1.72} & \underline{1.21$\times$} & \underline{1.018} \\
        FFT & 0.01 & 3 & 1.72 & \underline{1.21$\times$} & \textbf{1.030} \\
        \midrule
        DCT & 0.01 & 4 & \textbf{1.71} & \textbf{1.22$\times$} & \textbf{1.014} \\
        DWT & 0.01 & 4 & \underline{1.71} & \textbf{1.22$\times$} & \underline{1.006} \\
        FFT & 0.01 & 4 & 1.72 & \underline{1.21$\times$} & 0.996 \\
        \midrule
        DCT & 0.02 & 2 & \textbf{1.60} & \textbf{1.30$\times$} & \underline{1.009} \\
        DWT & 0.02 & 2 & \underline{1.61} & \textbf{1.30$\times$} & \textbf{1.010} \\
        FFT & 0.02 & 2 & \underline{1.61} & \textbf{1.30$\times$} & 1.002 \\
        \midrule
        DCT & 0.02 & 3 & \textbf{1.53} & \textbf{1.36$\times$} & \underline{0.967} \\
        DWT & 0.02 & 3 & \underline{1.53} & \textbf{1.36$\times$} & 0.966 \\
        FFT & 0.02 & 3 & 1.53 & \textbf{1.36$\times$} & \textbf{0.980} \\
        \midrule
        DCT & 0.02 & 4 & \textbf{1.53} & \textbf{1.36$\times$} & \textbf{0.964} \\
        DWT & 0.02 & 4 & \underline{1.53} & \textbf{1.36$\times$} & \underline{0.955} \\
        FFT & 0.02 & 4 & 1.53 & \textbf{1.36$\times$} & 0.949 \\
        \midrule
        DCT & 0.05 & 2 & \textbf{1.44} & \textbf{1.44$\times$} & \underline{0.989} \\
        DWT & 0.05 & 2 & \underline{1.45} & \textbf{1.44$\times$} & \textbf{0.991} \\
        FFT & 0.05 & 2 & 1.45 & \underline{1.43$\times$} & 0.972 \\
        \midrule
        DCT & 0.05 & 3 & \textbf{1.34} & \textbf{1.55$\times$} & \underline{0.936} \\
        DWT & 0.05 & 3 & \underline{1.34} & \textbf{1.55$\times$} & \textbf{0.937} \\
        FFT & 0.05 & 3 & 1.35 & \underline{1.54$\times$} & 0.933 \\
        \midrule
        DCT & 0.05 & 4 & \textbf{1.33} & \textbf{1.56$\times$} & \textbf{0.918} \\
        DWT & 0.05 & 4 & \underline{1.34} & \textbf{1.56$\times$} & \underline{0.902} \\
        FFT & 0.05 & 4 & 1.34 & \underline{1.55$\times$} & 0.895 \\
        \midrule
        DCT & 0.1 & 2 & \textbf{1.29} & \textbf{1.62$\times$} & \underline{0.966} \\
        DWT & 0.1 & 2 & \underline{1.29} & \textbf{1.62$\times$} & \textbf{0.977} \\
        FFT & 0.1 & 2 & 1.30 & \underline{1.61$\times$} & 0.944 \\
        \midrule
        DCT & 0.1 & 3 & \textbf{1.15} & \textbf{1.81$\times$} & \underline{0.880} \\
        DWT & 0.1 & 3 & 1.17 & 1.78$\times$ & \textbf{0.892} \\
        FFT & 0.1 & 3 & \underline{1.17} & \underline{1.79$\times$} & 0.877 \\
        \midrule
        DCT & 0.1 & 4 & \textbf{1.15} & \textbf{1.82$\times$} & \underline{0.864} \\
        DWT & 0.1 & 4 & \underline{1.15} & \underline{1.81$\times$} & \textbf{0.867} \\
        FFT & 0.1 & 4 & 1.16 & 1.80$\times$ & 0.841 \\
        \bottomrule
    \end{tabular}
    \end{tiny}
\end{table*}


\paragraph{Computational efficiency and scaling.} The inference cost of Spectral Progressive Diffusion is determined by the number of denoising steps and the token count processed at each resolution stage. For a fixed pretrained model and resolution schedule, the per-sample cost is independent of the dataset size, so generating or evaluating more samples scales linearly with the number of samples. The spectral transformation and noise-expansion operations add only lightweight per-transition overhead compared to the DiT forward passes. 

\paragraph{Spectral Transformation $T_\Phi$.} In Fig. \ref{fig:ablation-T}, we observe that DCT, DWT, and FFT achieve similar high-fidelity performance. We default to DCT for its real-valued computation and native support of non-power-of-two resolution scaling.

\paragraph{Error tolerance $\delta$.} In Fig. \ref{fig:ablation-delta}, we observe that increasing $\delta$ improves efficiency, but results in ghosting and halo artifacts near detailed edges.

\paragraph{Resolution stages $S$.} In Fig. \ref{fig:ablation-s}, we find that increasing $S$ leads to marginal speedup improvements and little image quality degradation. Since the primary speedup derives from bypassing native high-resolution steps, $S=2$ achieves the largest speedup while higher $S$ provides diminishing efficiency gains by further optimizing the already-inexpensive low-resolution trajectory.




\newpage
\begin{figure}
    \centering
    \includegraphics[width=\linewidth]{supp_figures/ablation_T.pdf}
    \caption{\textbf{Qualitative ablation on $T_\Phi$.} We see that all spectral transformations attain similar image quality. }
    \label{fig:ablation-T}
\end{figure}

\newpage
\begin{figure}
    \centering
    \includegraphics[width=\linewidth]{supp_figures/ablation_delta.pdf}
    \caption{\textbf{Qualitative ablation on $\delta$.} We observe that increasing $\delta$ improves efficiency, but results in ghosting and halo artifacts near detailed edges. }
    \label{fig:ablation-delta}
\end{figure}

\newpage
\begin{figure}
    \centering
    \includegraphics[width=\linewidth]{supp_figures/ablation_s.pdf}
    \caption{\textbf{Qualitative ablation on $S$.} We find that increasing $S$ leads to marginal speedup improvements and little image quality degradation.}
    \label{fig:ablation-s}
\end{figure}


\clearpage
\section{Additional Frequency-based Image Editing Details}
\label{sec:image_editing}

\paragraph{Detailed Procedure of Frequency-based Image Editing.} In our frequency-based editing pipeline, we perform the following operations given an input image $\mathbf{x}_\text{in}$, resolution scales $s_{1:S}$, and transition times $t_{1:S-1}$. We assume that the image editing process starts from transition time $t_k$, with $k \in \{1, 2, \ldots, S-1\}$.

\begin{enumerate}
    \item Compute the spectrum of the input image $\boldsymbol{\xi}_\text{in} = T_\Phi(\mathbf{x}_\text{in})$.
    \item Extract the low-frequency part of $\boldsymbol{\xi}_\text{in}$ corresponding to the representable spectrum $\Omega_{s_k}$ at scale $s_k$. 
    \item Fill the spectrum $\Omega_{s_{k+1}} \setminus \Omega_{s_k}$ with $t_k\epsilon$, where $\epsilon \sim \mathcal{N}(0, 1)$.
    \item Convert the expanded spectrum $\Omega_{s_{k+1}}$ back to the spatial domain $\mathbf{x}_{t_k}^{s_{k+1}}$ via $T^{-1}_\Phi$ and continue denoising starting from the transition time $t_k$.
\end{enumerate}

The total time of denoising is thus $1-t_k$.

We conduct image editing experiments using Z-Image \citep{zimage} as the backbone, evaluating our frequency-domain image editing pipeline against the SDEdit baseline \citep{meng2022sdedit}. We set $S=2$ and  $\delta=0.01$.

\paragraph{Extended Qualitative Results.} We show extended image editing results on two tasks: \textit{texture editing} and \textit{artistic stylization}. We configure our method with $\delta=0.01$ and $S=2$, which corresponds to a resolution transition at step 26 of the default 50-step denoising trajectory. We set $T_\Phi=\text{DCT}$, and evaluate the effect of different spectral transforms in the next section. For the SDEdit baseline, we initialize the process by corrupting the clean source image to the equivalent noise level at step 26 and subsequently proceed with the denoising schedule. This setting accurately reflects the editing method's capability to inject textural or style information, which are considered high-frequency details, in later denoising steps. 

As shown in Figures \ref{fig:texture_edit_1} to \ref{fig:art_stylization}, our frequency-based editing method achieves superior prompt alignment and geometric consistency compared to SDEdit, successfully changing the texture and artistic style of the input image while retaining the topological structure of the input. While SDEdit-style editing at earlier timesteps (higher noise levels at step 17) can facilitate stylistic changes, it introduces significant structural drift and geometric divergence from the source image. Furthermore, starting from early timesteps incurs substantial computational overhead and requires a brittle renoising to a specific, difficult-to-tune initialization timestep to balance editability with structural fidelity.

\paragraph{Effect of $T_\Phi$ on Image Editing.} As shown in Fig. \ref{fig:T_comparison_editing}, FFT-, DCT- and DWT-based editing achieve similar editing quality.

\begin{figure}
    \centering
    \includegraphics[width=\linewidth]{supp_figures/image_editing_supp_1.pdf}
    \caption{\textbf{Texture editing results.} Our frequency-based editing framework outperforms SDEdit, enabling high-fidelity texture transfer while preserving the geometric structure of the input image.}
    \label{fig:texture_edit_1}
\end{figure}

\begin{figure}
    \centering
    \includegraphics[width=\linewidth]{supp_figures/image_editing_supp_2.pdf}
    \caption{\textbf{Texture editing results.} Our frequency-based editing framework outperforms SDEdit, enabling high-fidelity texture transfer while preserving the geometric structure of the input image.}
    \label{fig:texture_edit_2}
\end{figure}

\begin{figure}
    \centering
    \includegraphics[width=\linewidth]{supp_figures/image_editing_supp_3.pdf}
    \caption{\textbf{Texture editing results.} Our frequency-based editing framework outperforms SDEdit, enabling high-fidelity texture transfer while preserving the geometric structure of the input image.}
    \label{fig:texture_edit_3}
\end{figure}

\begin{figure}
    \centering
    \includegraphics[width=\linewidth]{supp_figures/image_editing_supp_4.pdf}
    \caption{\textbf{Texture editing results.} Our frequency-based editing framework outperforms SDEdit, enabling high-fidelity texture transfer while preserving the geometric structure of the input image.}
    \label{fig:texture_edit_4}
\end{figure}


\begin{figure}
    \centering
    \includegraphics[width=\linewidth]{supp_figures/artistic_supp.pdf}
    \caption{\textbf{Artistic stylization results.} Aside from texture editing, our frequency-based editing approach also supports artistic stylization given stylistic descriptions and a representative artist. }
    \label{fig:art_stylization}
\end{figure}

\begin{figure}
    \centering
    \includegraphics[width=\linewidth]{supp_figures/image_editing_T_comparison.pdf}
    \caption{\textbf{Effect of $T_\Phi$ on image editing.} FFT-, DCT- and DWT-based editing achieve similar editing quality.}
    \label{fig:T_comparison_editing}
\end{figure}


\clearpage
\section{Broader Impacts}
\label{app:broader-impacts}

Spectral Progressive Diffusion is a principled method for improving the efficiency of pretrained image and video generation models, rather than a system targeted at a particular deployment domain. Its main positive impact is to reduce the compute, latency, and energy cost of generation, which can make experimentation and creative workflows more accessible and lower the environmental footprint of repeated generation. At the same time, because the method can accelerate existing generative models, we acknowledge that it could also lower the cost of producing synthetic media for undesirable uses such as disinformation, impersonation, or spam, depending on the pretrained model and release setting. We do not release new datasets or generation services, and any future release of fine-tuned adapters should follow the access terms, safeguards, and usage restrictions of the underlying pretrained models.

\section{Existing Assets and Licenses}
\label{app:assets-licenses}

Our experiments build on publicly available pretrained models, datasets, evaluation benchmarks, and software packages, including FLUX.1-dev~\citep{black2024flux}, Z-Image~\citep{zimage}, PixelGen~\citep{pixelgen}, WAN~\citep{wan}, MS-COCO~\citep{lin2014microsoft}, GenEval~\citep{geneval}, T2I-CompBench~\citep{t2icompbench}, VBench~\citep{vbench}, and the evaluation metrics used in Sec.~\ref{sec:exp}, including ImageReward~\citep{imagereward}, CLIP-IQA~\citep{wang2023clipiqa}, and NIQE~\citep{mittal2013niqe}. We cite the original sources for these assets and use them under their respective licenses and access terms.
The named licenses and access terms are: FLUX.1-dev is released under the FLUX.1 [dev] Non-Commercial License; Z-Image, PixelGen, WAN~2.1, VBench, ImageReward, VChitect-T2V-Dataverse, and the PyIQA implementation used for NIQE are under the Apache License 2.0; GenEval and T2I-CompBench are under the MIT License; and CLIP-IQA is under the NTU S-Lab License 1.0.
MS-COCO annotations are under Creative Commons Attribution 4.0 International (CC BY 4.0), while its images follow their original Flickr licenses; the public page for Aesthetic-Train-V2 does not list an explicit license, so we use it only under its stated access terms and do not redistribute the dataset.

% \subsection{Training-free Acceleration of Latent-space Video Generation}
% \label{app:video_generation}

% We run all latent-space video generation experiments on a single H100 GPU and measure wall-clock time measurements. For quantitative evaluation, we generate 200 videos using 200 test prompts and report the standard VBench scores across six evaluation dimensions, including Subject Consistency, Background Consistency, Motion Smoothness, Dynamic Degree, Aesthetic Quality, and Image Quality, in Table~\ref{tab:wan_vbench_720p}. All videos are generated in 720P resolution. For additional video results and qualitative comparisons, please see the supplementary website (\texttt{website/index.html}).

% \section{Classification of general progressive-resolution transforms}
% \label{app:family}
% \hx{Under construction...}

% Progressive resolution generation requires a principled transition between a low-resolution flow state and a high-resolution flow state. In this appendix we characterize \emph{all} such transitions in closed form: they form an affine-linear family determined by an upsample operator $U$ and an associated noise-injection matrix $B$. Three specific choices of $U$ (block-repeat, bilinear, sinc) recover Haar wavelet, spatial interpolation, and Fourier-domain constructions respectively.

% \subsection{Problem Setup}
% \label{app:family:setup}

% Let $x \in \mathbb{R}^{N_l}$ be a clean data sample at low resolution, and let $y \in \mathbb{R}^{N_h}$ denote the corresponding clean sample at high resolution, related by a fixed linear upsample $U \in \mathbb{R}^{N_h \times N_l}$ with $y = U x$ and $N_h > N_l$. Write $k = N_h / N_l$ for the total count ratio (so $k = c^d$ if each of $d$ spatial dimensions grows by a factor $c$).

% At some diffusion time $t \in (0, 1)$, the low-resolution rectified-flow state is
% \begin{equation}
%   x_t \;=\; (1-t)\, x + t\, \epsilon_l,
%   \qquad \epsilon_l \sim \mathcal{N}(0, I_{N_l}),
%   \label{eq:low-res-state}
% \end{equation}
% and we seek a \emph{resolution lift} $L_t : \mathbb{R}^{N_l} \to \mathbb{R}^{N_h}$ (possibly using auxiliary randomness $\xi$) whose output is a valid flow state of $y = Ux$ at some target time $t'$:
% \begin{equation}
%   L_t(x_t;\, \xi)
%   \;\stackrel{d}{=}\;
%   (1-t')\, U x + t'\, \epsilon_h,
%   \qquad \epsilon_h \sim \mathcal{N}(0, I_{N_h}),
%   \label{eq:high-res-state}
% \end{equation}
% with $\xi$ independent of the low-resolution noise $\epsilon_l$. The requirement that $\epsilon_h$ be i.i.d.\ standard Gaussian is essential: the pretrained denoiser $v_\theta$ was trained on precisely this distribution, and any departure produces out-of-distribution inputs.

% \subsection{Characterization: $L_t$ Must Be Affine-Linear}
% \label{app:family:affine}

% \begin{theorem}[Affine-linearity]
% \label{thm:affine}
% Any $L_t$ satisfying \eqref{eq:high-res-state} for all clean signals $x \in \mathbb{R}^{N_l}$ and all realizations of $\epsilon_l$ is an affine-linear function of its first argument:
% \begin{equation}
%   L_t(y; \xi) \;=\; A\, y + B\, \xi,
%   \label{eq:affine}
% \end{equation}
% where $A \in \mathbb{R}^{N_h \times N_l}$, $B \in \mathbb{R}^{N_h \times m}$ for some $m$, and $\xi \sim \mathcal{N}(0, I_m)$ is an auxiliary Gaussian independent of the argument.
% \end{theorem}

% \begin{proof}[Proof sketch]
% The right-hand side of \eqref{eq:high-res-state} is Gaussian with mean $(1-t')\, U x$ \emph{linear} in $x$ and covariance $t'^2\, I_{N_h}$ \emph{independent of $x$}. Translating $x \mapsto x + \Delta x$ shifts the argument of $L_t$ by $(1-t)\, \Delta x$ and must shift the output distribution by $(1-t')\, U \Delta x$. Since $\xi$ is independent of the argument, this shift property holds pointwise in $\xi$:
% \begin{equation*}
%   L_t\bigl(y + (1-t)\,\Delta x;\; \xi\bigr)
%   \;=\;
%   L_t(y;\; \xi) \,+\, (1-t')\, U\, \Delta x .
% \end{equation*}
% A measurable function satisfying Cauchy's functional equation $f(y + \Delta) - f(y) = \ell(\Delta)$ for all $\Delta$ must be affine-linear in $y$; this yields $L_t(y; \xi) = A\, y + g(\xi)$ for some linear $A$ and measurable $g$. Finally, $L_t(y; \xi)$ must be Gaussian for every fixed $y$ (so that the distributional identity in~\eqref{eq:high-res-state} is preserved under conditioning on $y$), which forces $g(\xi) = B\, \xi$ for a linear $B$.
% \end{proof}

% \subsection{Closed-Form Solution for $A$ and $B$}
% \label{app:family:closedform}

% Given Theorem~\ref{thm:affine}, we substitute $y = x_t = (1-t)\, x + t\, \epsilon_l$ into \eqref{eq:affine} and match moments against the target distribution in \eqref{eq:high-res-state}.

% \paragraph{Mean matching.} Linearity in $x$ gives $A\, (1-t) x = (1-t')\, U x$ for all $x$, so
% \begin{equation}
%   \boxed{\;\;A \;=\; \frac{1-t'}{1-t}\, U\;\;}
%   \label{eq:A-closedform}
% \end{equation}
% The signal-path operator is therefore the upsample $U$, rescaled by the ratio of signal coefficients at the two times.

% \paragraph{Covariance matching.} The output covariance is the sum of contributions from $\epsilon_l$ (via $A$) and from $\xi$ (via $B$), and must equal $t'^2\, I_{N_h}$:
% \begin{equation}
%   t^2\, A A^\top + B B^\top \;=\; t'^2\, I_{N_h}.
%   \label{eq:cov-match}
% \end{equation}
% Substituting \eqref{eq:A-closedform} and solving for $B B^\top$,
% \begin{equation}
%   \boxed{\;\;B B^\top \;=\; t'^2\, I_{N_h} \,-\, t^2 \!\left(\tfrac{1-t'}{1-t}\right)^{\!2}\, U U^\top\;\;}
%   \label{eq:B-closedform}
% \end{equation}
% Any $B$ whose Gram matrix equals the right-hand side of \eqref{eq:B-closedform} (e.g., its symmetric square root) realizes a valid lift.

% \paragraph{$B$ is necessarily nonzero.} Since $U \in \mathbb{R}^{N_h \times N_l}$ maps between spaces of different dimension, $\operatorname{rank}(U U^\top) \le N_l < N_h$. The target $t'^2 I_{N_h}$ has full rank $N_h$. Hence the right-hand side of \eqref{eq:B-closedform} has rank at least $N_h - N_l > 0$: \emph{any genuine resolution lift requires injecting fresh Gaussian noise}, whose variance fills the $N_h - N_l$ directions not reached by $U$. This is a direct structural consequence of the dimension increase and is independent of the specific choice of $U$.

% \paragraph{PSD feasibility and the minimum-$t'$ theorem.} For $B$ to exist as a real matrix, $B B^\top$ must be positive semi-definite. Applying this to \eqref{eq:B-closedform},
% \begin{equation}
%   t'^2\, I_{N_h} \;\succeq\; t^2\! \left(\tfrac{1-t'}{1-t}\right)^{\!2} U U^\top,
%   \label{eq:PSD-condition}
% \end{equation}
% which reduces to a scalar inequality on the maximum eigenvalue of $U U^\top$:
% \begin{equation}
%   t'^2 \;\ge\; t^2\!\left(\tfrac{1-t'}{1-t}\right)^{\!2} \lambda_{\max}(U U^\top).
%   \label{eq:PSD-scalar}
% \end{equation}
% This bound is \emph{always active} for DC-preserving upsamples, as we now show.

% \begin{theorem}[Minimum-$t'$ bound for DC-preserving upsamples]
% \label{thm:min-t-prime}
% Let $U \in \mathbb{R}^{N_h \times N_l}$ be a DC-preserving upsample, i.e., $U\,\mathbf{1}_{N_l} = \mathbf{1}_{N_h}$ where $\mathbf{1}_n \in \mathbb{R}^n$ is the all-ones vector. Let $k = N_h / N_l$. Then $\lambda_{\max}(U U^\top) \ge k$, and the PSD feasibility condition \eqref{eq:PSD-scalar} combined with $t, t' \in (0, 1)$ forces
% \begin{equation}
%   \boxed{\;\; t' \;\ge\; \frac{\sqrt{\lambda_{\max}(U U^\top)} \; t}{1 + \bigl(\sqrt{\lambda_{\max}(U U^\top)} - 1\bigr)\, t} \;\;}
%   \label{eq:min-t-prime}
% \end{equation}
% with equality achievable when $\lambda_{\max}(U U^\top) = k$ (i.e., when $U$ attains the DC-preserving lower bound). In particular, for any $U$ that saturates $\lambda_{\max}(U U^\top) = k$,
% \begin{equation}
%   t'_{\min} \;=\; \frac{\sqrt{k}\, t}{1 + (\sqrt{k} - 1)\, t},
%   \label{eq:min-t-prime-tight}
% \end{equation}
% giving the tightest possible target time while preserving flow-matching structure.
% \end{theorem}

% \begin{proof}[Proof sketch]
% The Rayleigh quotient at the all-ones vector gives $\lambda_{\max}(U U^\top) \ge \|U\,\mathbf{1}_{N_l}\|_2^2 / \|\mathbf{1}_{N_l}\|_2^2 = \|\mathbf{1}_{N_h}\|_2^2 / N_l = N_h/N_l = k$. Let $s := \sqrt{\lambda_{\max}(U U^\top)}$. Taking square roots of \eqref{eq:PSD-scalar} (all quantities positive for $t, t' \in (0,1)$), $t'(1-t) \ge s\, t (1-t')$. Rearranging, $t' + s\,t\,t' \ge s\,t + t\,t'$, so $t'\bigl(1 + (s-1)\,t\bigr) \ge s\,t$, yielding \eqref{eq:min-t-prime}. Equality in \eqref{eq:PSD-scalar} is achieved when the eigenspace of $\lambda_{\max}$ saturates and all other eigenvalues lie below $\lambda_{\max}$; this is exactly the condition that $U U^\top$ has a rank-one eigenvector at the all-ones direction with eigenvalue $k$ and no direction exceeds $k$, i.e., $\|U\|_{\mathrm{op}}^2 = k$.
% \end{proof}

% \emph{Interpretation.} Equation \eqref{eq:min-t-prime-tight} captures a fundamental ``rewind'' required at every resolution transition: since $t'_{\min} > t$ for any $k > 1$, the target state must always sit at a \emph{noisier} time than the source. The rewind grows with $k$: at $t = 0.5$, $t'_{\min} = \sqrt{k}/(\sqrt{k}+1)$ yields $0.67$ for $k{=}4$, $0.75$ for $k{=}9$, and $0.80$ for $k{=}16$. Upsamples that amplify a non-DC direction above $k$ (e.g., bicubic, which exhibits a spectral overshoot) carry $\lambda_{\max} > k$ and therefore pay a strictly larger rewind than block-repeat, bilinear, and sinc---all of which saturate $\lambda_{\max}(U U^\top) = k$ exactly and thus achieve the tight minimum \eqref{eq:min-t-prime-tight}.

% \subsection{Three Canonical Instances of $U$}
% \label{app:family:instances}

% We now specialize \eqref{eq:A-closedform}--\eqref{eq:B-closedform} to three choices of $U$ used throughout the literature. We give closed-form expressions for $A$ and $B B^\top$ in each case, and illustrate them in the elementary $2 \!\times\! 2 \to 4 \!\times\! 4$ example (so $N_l = 4$, $N_h = 16$, and $k = 4$).

% \paragraph{(i) Block-repeat upsample (Haar LL-only IDWT).}
% The nearest-neighbor lift replicates each low-resolution coefficient $x_{(i)}$ to the $k$ high-resolution positions it covers:
% \begin{equation}
%   [U_{\mathrm{BR}}\, x]_{(i,\alpha)} \;=\; x_{(i)},
%   \qquad \alpha \in \{0, \ldots, k-1\}.
%   \label{eq:BR-def}
% \end{equation}
% For the $2\!\times\!2 \to 4\!\times\!4$ case, each input pixel is replicated into its $2\!\times\!2$ output block. The Gram matrix $U_{\mathrm{BR}} U_{\mathrm{BR}}^\top$ is block-diagonal: each $N_l$-group of $k$ replicated positions yields a $k \!\times\! k$ all-ones block $J_k$, with zero off-block entries. Its eigendecomposition has eigenvalue $k$ on the per-block constant vector and $0$ on the $(k{-}1)$-dimensional orthogonal complement (per block). Substituting into \eqref{eq:A-closedform}--\eqref{eq:B-closedform},
% \begin{equation}
%   A_{\mathrm{BR}} = \tfrac{1-t'}{1-t}\, U_{\mathrm{BR}},
%   \qquad
%   B_{\mathrm{BR}} B_{\mathrm{BR}}^\top = t'^2\, I_{N_h} \;-\; t^2\!\left(\tfrac{1-t'}{1-t}\right)^{\!2}\,
%   \bigoplus_{i=1}^{N_l} J_k .
%   \label{eq:BR-AB}
% \end{equation}
% Equivalently, $B_{\mathrm{BR}} B_{\mathrm{BR}}^\top$ is block-diagonal with per-block matrix $t'^2 I_k - t^2\!\left(\tfrac{1-t'}{1-t}\right)^{\!2} J_k$, whose eigenvalues are $t'^2$ on the $(k{-}1)$-dim complement of the constant direction and $t'^2 - k\, t^2\!\left(\tfrac{1-t'}{1-t}\right)^{\!2}$ on the constant direction itself.

% \emph{Connection to the Haar wavelet.} $U_{\mathrm{BR}}$ is exactly $\sqrt{k}$ times the orthonormal Haar inverse DWT applied to the coefficient vector $(x, 0, \ldots, 0)$ (low-frequency slot populated; detail slots zero). The block-diagonal decomposition into \emph{per-block constant direction} and \emph{per-block detail complement} is precisely the LL-vs-\{LH, HL, HH\} split of the Haar basis in the $2\!\times\!2 \to 4\!\times\!4$ case.

% \paragraph{(ii) Bilinear upsample.}
% The bilinear lift convolves $x$ with a separable tent kernel, yielding weighted averages of neighboring coefficients. For the $2\!\times\!2 \to 4\!\times\!4$ case with half-pixel offset, each output pixel receives four weights $w_1 w_1', w_1 w_2', w_2 w_1', w_2 w_2'$ with $(w_1, w_2) = (\tfrac{3}{4}, \tfrac{1}{4})$ in each dimension:
% \begin{equation}
%   [U_{\mathrm{BL}}\, x]_{(2i+a, 2j+b)}
%   \;=\; \sum_{p,q \in \{-1,0\}} w_{a-p}\, w_{b-q}\, x_{(i+p, j+q)}.
%   \label{eq:BL-def}
% \end{equation}
% The Gram matrix $U_{\mathrm{BL}} U_{\mathrm{BL}}^\top$ is now a \emph{banded, non-block-diagonal} matrix: adjacent $2\!\times\!2$ blocks share weights, so correlation leaks across block boundaries. It has the same maximum eigenvalue $k$ (attained on the all-ones direction, since $U_{\mathrm{BL}}$ preserves DC by construction $\sum_\alpha w_\alpha = 1$), but its eigenstructure decays smoothly with frequency rather than collapsing to $\{k, 0\}$. Substituting:
% \begin{equation}
%   A_{\mathrm{BL}} = \tfrac{1-t'}{1-t}\, U_{\mathrm{BL}},
%   \qquad
%   B_{\mathrm{BL}} B_{\mathrm{BL}}^\top = t'^2\, I_{N_h} \,-\, t^2\!\left(\tfrac{1-t'}{1-t}\right)^{\!2}\, U_{\mathrm{BL}} U_{\mathrm{BL}}^\top.
%   \label{eq:BL-AB}
% \end{equation}
% Taking the symmetric square root gives a frequency-shaped noise injector that is strongest where $U_{\mathrm{BL}} U_{\mathrm{BL}}^\top$ has the smallest eigenvalues (high spatial frequencies, where the bilinear transfer function rolls off). Unlike the block-repeat case, $B_{\mathrm{BL}}$ has no exact spatial locality.

% \paragraph{(iii) Sinc (ideal bandlimited) upsample.}
% In the Fourier domain, the ideal lift zero-pads the spectrum: $U_{\mathrm{S}}$ takes $x$, computes its orthonormal DFT $X = \mathcal{F}_{N_l} x$, places $X$ in the low-frequency block of an $N_h$-point spectrum, and inverts. Algebraically,
% \begin{equation}
%   U_{\mathrm{S}} \;=\; \sqrt{k}\, \mathcal{F}_{N_h}^{-1}\, P_L\, \widetilde{\mathcal{F}}_{N_l},
%   \qquad
%   P_L : \mathbb{R}^{N_h} \to \mathbb{R}^{N_h}
%   \text{ is the projection onto the low-frequency support},
%   \label{eq:sinc-def}
% \end{equation}
% where $\widetilde{\mathcal{F}}_{N_l}$ embeds the $N_l$-point spectrum into the low-frequency block of an $N_h$-point grid. For $2\!\times\!2 \to 4\!\times\!4$, $P_L$ selects the two lowest horizontal and two lowest vertical frequency bins (the central $2\!\times\!2$ block of the $4\!\times\!4$ shifted spectrum). The Gram matrix is a \emph{pure projection}:
% \begin{equation}
%   U_{\mathrm{S}} U_{\mathrm{S}}^\top \;=\; k\, P_L,
%   \label{eq:sinc-gram}
% \end{equation}
% with eigenvalues $k$ on the $N_l$-dim low-frequency subspace and $0$ on the $N_h - N_l$-dim high-frequency complement. Substituting,
% \begin{equation}
%   A_{\mathrm{S}} = \tfrac{1-t'}{1-t}\, U_{\mathrm{S}},
%   \qquad
%   B_{\mathrm{S}} B_{\mathrm{S}}^\top = t'^2\, I_{N_h} \,-\, k\, t^2\!\left(\tfrac{1-t'}{1-t}\right)^{\!2}\, P_L
%   \;=\; \underbrace{\Big(t'^2 - k\, t^2\!\left(\tfrac{1-t'}{1-t}\right)^{\!2}\Big)\, P_L}_{\text{low-freq residual}} \;+\; t'^2\, P_H,
%   \label{eq:sinc-AB}
% \end{equation}
% where $P_H = I_{N_h} - P_L$. When the rewind time $t'$ is chosen so that the low-frequency coefficient vanishes (i.e., $t' = \sqrt{k}\, t / (1 + (\sqrt{k}-1) t)$), $B_{\mathrm{S}}$ reduces to a \emph{pure high-pass} noise injector: $B_{\mathrm{S}} = t'\, P_H$, which simply fills the high-frequency half of the $N_h$-spectrum with fresh i.i.d.\ complex Gaussian noise. This is the cleanest form of \eqref{eq:B-closedform} and motivates the Fourier interpretation of the lift.

% \emph{Connection to FFT padding.} The sinc lift is exactly ``FFT, zero-pad spectrum, inverse FFT'' when the detail-band noise is zero. With the prescribed $B_{\mathrm{S}}$, it becomes ``FFT, pad with complex Gaussian noise at amplitude $t'$, inverse FFT''---a single-operator drop-in replacement for the wavelet lift.

% \paragraph{Summary.}
% All three instances share the same $A = \tfrac{1-t'}{1-t}\, U$ in \eqref{eq:A-closedform} and saturate the DC eigenvalue $\lambda_{\max}(U U^\top) = k$; they differ only in how $B B^\top$ distributes the required $N_h - N_l$ units of fresh variance. The block-repeat choice concentrates $B$ in per-block detail directions (Haar-local), bilinear spreads $B$ as a smooth high-pass filter (mildly non-local), and sinc projects $B$ onto the exact high-frequency subspace of the global DFT (fully non-local, frequency-ideal).



% \iffalse
% \section{Additional Quantitative Results}
% \label{app:additional-quant}

% \subsection{SD3 Main Results}
% \label{app:sd3-results}

% Table~\ref{tab:sd3-supp} replicates the main-paper FLUX.1-dev table (Table~\ref{tab:flux-main}) on Stable Diffusion~3~\citep{sd3} using the same RALU~\citep{ralu} evaluation protocol, at SD3's native 28 steps and 1024$\times$1024 resolution. Speedup targets are 2$\times$ and 3$\times$ per RALU. RAS~\citep{ras} is included as an SD3-specific baseline from the RALU table.

% \begin{table}[t]
%   \centering
%   \footnotesize
%   \caption{\textbf{Image results on Stable Diffusion~3}, matching RALU's~\citep{ralu} SD3 table. Placeholder values.}
%   \label{tab:sd3-supp}
%   \setlength{\tabcolsep}{2.5pt}
%   \begin{tabular}{l ccc ccccc}
%     \toprule
%     Method & Latency $\downarrow$ & TFLOPs $\downarrow$ & Speedup $\uparrow$ & FID $\downarrow$ & NIQE $\downarrow$ & CLIP-IQA $\uparrow$ & T2I-Comp.\ $\uparrow$ & GenEval $\uparrow$ \\
%     \midrule
%     SD3 (28 steps, base)                   & - & - & 1$\times$ & - & - & - & - & - \\
%     SD3 (reduced steps)                    & - & - & -         & - & - & - & - & - \\
%     ToCa~\citep{toca}                      & - & - & -         & - & - & - & - & - \\
%     Bottleneck Sampling~\citep{bottleneck} & - & - & -         & - & - & - & - & - \\
%     RAS~\citep{ras}                        & - & - & -         & - & - & - & - & - \\
%     RALU~\citep{ralu}                      & - & - & -         & - & - & - & - & - \\
%     Ours (spectral passthrough)            & - & - & -         & - & - & - & - & - \\
%     \textbf{Ours-TF (Haar)}                & - & - & -         & - & - & - & - & - \\
%     \textbf{Ours-TF (FFT)}                 & - & - & -         & - & - & - & - & - \\
%     \textbf{Ours-LoRA (Haar)}              & - & - & -         & - & - & - & - & - \\
%     \bottomrule
%   \end{tabular}
% \end{table}

% \subsection{PixelGen (Pixel-Space) Main Results}
% \label{app:pixelgen-results}

% Table~\ref{tab:pixelgen-supp} reports the same protocol on the pixel-space backbone PixelGen~\citep{pixelgen}, which operates directly on RGB pixels rather than a VAE latent.

% \begin{table}[t]
%   \centering
%   \footnotesize
%   \caption{\textbf{Pixel-space image results on PixelGen}~\citep{pixelgen}. Same RALU-style columns as Table~\ref{tab:flux-main}; baselines restricted to methods applicable in pixel space. Placeholder values.}
%   \label{tab:pixelgen-supp}
%   \setlength{\tabcolsep}{2.5pt}
%   \begin{tabular}{l ccc ccccc}
%     \toprule
%     Method & Latency $\downarrow$ & TFLOPs $\downarrow$ & Speedup $\uparrow$ & FID $\downarrow$ & NIQE $\downarrow$ & CLIP-IQA $\uparrow$ & T2I-Comp.\ $\uparrow$ & GenEval $\uparrow$ \\
%     \midrule
%     PixelGen (base)                        & - & - & 1$\times$ & - & - & - & - & - \\
%     PixelGen (reduced steps)               & - & - & -         & - & - & - & - & - \\
%     Bottleneck Sampling~\citep{bottleneck} & - & - & -         & - & - & - & - & - \\
%     Ours (spectral passthrough)            & - & - & -         & - & - & - & - & - \\
%     \textbf{Ours-TF (Haar)}                & - & - & -         & - & - & - & - & - \\
%     \textbf{Ours-TF (FFT)}                 & - & - & -         & - & - & - & - & - \\
%     \textbf{Ours-LoRA (Haar)}              & - & - & -         & - & - & - & - & - \\
%     \bottomrule
%   \end{tabular}
% \end{table}

% \fi


%%%%%%%%%%%%%%%%%%%%%%%%%%%%%%%%%%%%%%%%%%%%%%%%%%%%%%%%%%%%

% \newpage
% \input{checklist.tex}


\end{document}
